\chapter{The Lean Library: Path Sums}\label{ch:lean-pathsum}
\begingroup\sloppy

This chapter specifies the eighteen modules of the \tool{Lean~4}{https://lean-lang.org} library
\cite{demoura2021lean4} that certify path-sum computations: \file{PathSum}, \file{PathSumSem},
\file{PathSumRules}, \file{PathSumParity}, \file{PathSumAmp}, \file{PathSumSupport}, \file{PathSumChain},
\file{PathSumNet}, \file{PathSumDigits}, \file{PathSumW}, \file{PathSumNetW},
\file{PathSumNetAmp}, \file{PathSumQudit}, \file{PathSumQuditChain}, \file{PathSumQuditNet},
\file{PathSumQuditDigits}, \file{PathSumMixed}, and \file{PathSumMixedNet}, all under \file{archytaszx/lean/ArchytasZX/}. A \term{path sum} is
Amy's sum-over-paths representation of a linear map \cite{amy2018pathsums,vilmart2021sop}: a sum,
over Boolean (or $\Zd{p}$-valued) path variables, of a phase $\operatorname{ex}(P/M)$ restricted to
the solutions of polynomial constraints, times a constant. The Python engine reduces path sums by
rewrite rules (\cref{ch:semantics}); the modules of this chapter define path sums as syntax in
Lean, give them a value in $\C$, define one \emph{step} per reduction rule as a computable
function on that syntax together with a Boolean side condition, and prove each step sound: when the
side condition evaluates to \textnormal{\texttt{true}}, the step leaves the value unchanged up to a stated
constant factor. A certificate replays the engine's reduction trace as a list of steps; the Lean
kernel evaluates the side conditions by \textnormal{\texttt{decide}}, and the soundness theorems give the equality
of the values. No theorem of the library enumerates the assignments of path variables: every
soundness proof is an induction over lists and \textnormal{\texttt{Function.\allowbreak{}update}}.

The library depends on \tool{Mathlib}{https://github.com/leanprover-community/mathlib4}
\cite{mathlib2020} and on the core modules \file{Net} (tensor networks, \textnormal{\texttt{eval}}, \textnormal{\texttt{Equal}}),
\file{Cyclo} and \file{Ball} (roots of unity), \file{Sym}, \file{Closed}, \file{Positional}, and
\file{GaussSign}, which \cref{ch:lean-core} specifies. The module imports form the chain
\begin{center}\small
\begin{tabular}{@{}l@{}}
\file{PathSum} $\to$ \file{PathSumSem} $\to$ \file{PathSumRules} $\to$ \file{PathSumAmp}
$\to$ \file{PathSumSupport} $\to$ \file{PathSumChain},\\
\file{PathSumRules} $\to$ \file{PathSumNet} $\to$ \file{PathSumDigits} $\to$ \file{PathSumW}
$\to$ \file{PathSumNetW} $\to$ \file{PathSumNetAmp},\\
\file{PathSum} $\to$ \file{PathSumQudit} $\to$ \file{PathSumQuditChain} $\to$
\file{PathSumQuditNet} $\to$ \file{PathSumQuditDigits},
\end{tabular}
\end{center}
with \file{PathSumNetW} also importing \file{PathSumChain}, \file{PathSumNetAmp} also importing
\file{PathSumAmp}, \file{PathSumQuditChain} also importing \file{PathSumChain}, and
\file{PathSumQuditNet} also importing \file{PathSumDigits} and \file{GaussSign},
\file{PathSumParity} importing \file{PathSumRules} and imported by \file{PathSumChain},
\file{PathSumMixed} importing \file{PathSumQudit} and \file{PathSumAmp}, and \file{PathSumMixedNet}
importing \file{PathSumMixed}, \file{PathSumNetAmp}, \file{PathSumQuditNet}, and
\file{PathSumQuditDigits}. The mixed definitions live in namespace \textnormal{\texttt{ArchytasZX.MSem}}. The qubit
definitions live in namespace \textnormal{\texttt{ArchytasZX.\allowbreak{}PSem}}, the qudit definitions in
\textnormal{\texttt{ArchytasZX.\allowbreak{}QSem}}, and the rule identities of \file{PathSum} in \textnormal{\texttt{ArchytasZX}}. The last
section of the chapter specifies the Python exporters that write Lean files against these modules:
\module{archytaszx.semantics.lean_pathsum}, \module{archytaszx.semantics.lean_amp},
\module{archytaszx.semantics.lean_qdigits}, and the part of
\module{archytaszx.semantics.lean_chain} that reuses them.

\section{Notation}\label{sec:lean-pathsum-notation}

The library uses the following constants of the core modules.
\begin{align*}
  \textnormal{\texttt{ex}}\ \theta &= e^{2\pi i\theta} \quad(\theta\in\R), &
  \textnormal{\texttt{om}}\ d &= \omega_d = e^{2\pi i/d}, &
  \textnormal{\texttt{rt}}\ d &= \sqrt{d}\in\C, &
  \textnormal{\texttt{fou}}\ d\ j\ k &= \omega_d^{jk}/\sqrt{d}.
\end{align*}
In this chapter $\operatorname{ex}$, $\omega_d$, and $\sqrt{d}$ denote these terms. The core
lemma \textnormal{\texttt{om\_\allowbreak{}eq\_\allowbreak{}ex}} states $\omega_m^e = \operatorname{ex}(e/m)$, and \textnormal{\texttt{om\_\allowbreak{}sub}} states
$\omega_m = \omega_M^{q}$ for $M = mq$, $0<M$.

An \term{assignment} is a function $\sigma:\N\to\Zd{2}$ (qubit) or $\sigma:\N\to\Zd{p}$ (qudit);
variables are natural numbers. $\sigma[y\mapsto b]$ denotes \textnormal{\texttt{Function.\allowbreak{}update \ensuremath{\sigma} y b}}. For
$x\in\Zd{2}$, $\lfloor x\rfloor\in\{0,1\}\subset\Z$ denotes \textnormal{\texttt{zq x}}
(\textnormal{\texttt{if x = 0 then 0 else 1}}); the library proves by \textnormal{\texttt{decide}} that
$\lfloor x+y\rfloor = \lfloor x\rfloor+\lfloor y\rfloor-2\lfloor x\rfloor\lfloor y\rfloor$
(\textnormal{\texttt{zq\_\allowbreak{}add}}), $\lfloor xy\rfloor=\lfloor x\rfloor\lfloor y\rfloor$ (\textnormal{\texttt{zq\_\allowbreak{}mul}}),
$\lfloor x\rfloor \bmod 2 = x$ (\textnormal{\texttt{zq\_\allowbreak{}cast}}), $x\in\{0,1\}$ (\textnormal{\texttt{z2\_\allowbreak{}cases}}), and $x^2=x$
(\textnormal{\texttt{z2\_\allowbreak{}idem}}). Lists are written $[a_1,\dots,a_n]$; $\ell\mathbin{+\!\!+}\ell'$ is
concatenation. Iverson brackets $[\varphi]\in\{0,1\}$ denote the indicator of a proposition.

A \term{side condition} is a function into \textnormal{\texttt{Bool}}; a theorem hypothesis \textnormal{\texttt{h : xOK .\allowbreak{}.\allowbreak{}.\allowbreak{} = true}}
is discharged in a certificate by \textnormal{\texttt{by decide}}, which makes the Lean kernel evaluate the
function on the literal data. Every side condition below is a conjunction of decidable tests,
listed in the order of the Lean \textnormal{\texttt{\&\&}} chain.

\section{Rule identities: \texorpdfstring{\file{PathSum}}{PathSum}}\label{sec:lean-pathsum-rules0}

\file{PathSum} (imports \file{Sym}, \file{Cyclo}) states, for each rule of the Python reducers
\textnormal{\texttt{archytaszx.\allowbreak{}algebra.\allowbreak{}pathsum.\allowbreak{}reduce}} and \textnormal{\texttt{reduce\_\allowbreak{}qudit}}, the scalar identity the rule applies
to the sum over its eliminated variables, at the Boolean values $0,1$ of the variables its quotient
polynomials read (\textnormal{\texttt{ps\_\allowbreak{}*}}) or at a residue mod $p$ (\textnormal{\texttt{q\_\allowbreak{}*}}). In these statements
$(-1)^q$ stands for $\operatorname{ex}(q/2)$ and $\omega_8^e$ for $\operatorname{ex}(e/8)$. The
lemmas are the targets of the trace comments of \textnormal{\texttt{export\_\allowbreak{}pair}}
(\cref{sec:lean-pathsum-python}); the step-level soundness theorems of later sections do not use
them.

\begin{theorem}[Qubit rule identities]\label{thm:lean-pathsum-ps}
For $c\in\C$, $g:\N\to\C$, and natural numbers with the stated bounds:
\begin{alignat*}{2}
&\text{\textnormal{\texttt{ps\_\allowbreak{}elim}}:}\quad && c+c = 2c,\\
&\text{\textnormal{\texttt{ps\_\allowbreak{}hh}}:}\quad && 1+(-1)^q = \begin{cases}2 & q\equiv 0 \pmod 2\\ 0 & \text{else,}\end{cases}\\
&\text{\textnormal{\texttt{ps\_\allowbreak{}hh\_\allowbreak{}pin}} }(q\le 1)\text{:}\quad && (1+(-1)^{0+q})\,g(0)+(1+(-1)^{1+q})\,g(1) = 2\,g(q),\\
&\text{\textnormal{\texttt{ps\_\allowbreak{}omega0}}:}\quad && 1+\omega_8^2 = \sqrt2\,\omega_8,\\
&\text{\textnormal{\texttt{ps\_\allowbreak{}omega1}}:}\quad && 1+\omega_8^6 = \sqrt2\,\omega_8\,\omega_8^6,\\
&\text{\textnormal{\texttt{ps\_\allowbreak{}omega\_\allowbreak{}conj0}}:}\quad && 1+\omega_8^6 = \sqrt2\,\omega_8^7,\\
&\text{\textnormal{\texttt{ps\_\allowbreak{}omega\_\allowbreak{}conj1}}:}\quad && 1+\omega_8^2 = \sqrt2\,\omega_8^7\,\omega_8^2,\\
&\text{\textnormal{\texttt{ps\_\allowbreak{}hh\_\allowbreak{}gen}} }(a,b\le1,\ 1\le a+b)\text{:}\quad &&
  (1+(-1)^{0a+b})g(0)+(1+(-1)^{1a+b})g(1)\\
 &&& \qquad = (1+(-1)^{b(1+a)})\,g(b),\\
&\text{\textnormal{\texttt{ps\_\allowbreak{}rem}} }(m\le1)\text{:}\quad && (1+(-1)^m)(c^m g) = (1+(-1)^m)\,g,\\
&\text{\textnormal{\texttt{ps\_\allowbreak{}hh\_\allowbreak{}nl}} }(q,q'\le1)\text{:}\quad && (1+(-1)^q)(1+(-1)^{q'}) = 2\,(1+(-1)^{q+q'+qq'}),\\
&\text{\textnormal{\texttt{ps\_\allowbreak{}geom}} }(0<N)\text{:}\quad && \textstyle\sum_{Y<N}\omega_N^{YS} = N\,[S\equiv0\pmod N],\\
&\text{\textnormal{\texttt{ps\_\allowbreak{}digits}}:}\quad && \textstyle\prod_{j<k}(1+(-1)^{b_j}) = 2^k\,[S\equiv 0\pmod{2^k}],\\
&\text{\textnormal{\texttt{ps\_\allowbreak{}dyadic}}:}\quad && \textstyle\sum_{Y<2^k}\omega_{2^k}^{YS} = \prod_{j<k}(1+(-1)^{b_j}),
\end{alignat*}
where $b_j = \lfloor S/2^j\rfloor\bmod 2$ is the $j$-th binary digit of $S$.
\end{theorem}
\noindent The $\omega$ identities follow from \textnormal{\texttt{rt\_\allowbreak{}sub2}} ($\sqrt2=\omega_8+\omega_8^7$) and
$\omega_8^8=1$ by \textnormal{\texttt{linear\_\allowbreak{}combination}}; \textnormal{\texttt{ps\_\allowbreak{}hh\_\allowbreak{}pin}}, \textnormal{\texttt{ps\_\allowbreak{}hh\_\allowbreak{}gen}}, \textnormal{\texttt{ps\_\allowbreak{}rem}}, and
\textnormal{\texttt{ps\_\allowbreak{}hh\_\allowbreak{}nl}} by \textnormal{\texttt{interval\_\allowbreak{}cases}} on the bits; \textnormal{\texttt{ps\_\allowbreak{}geom}} by the geometric series with
$\omega_N^S\ne1$ when $N\nmid S$; \textnormal{\texttt{ps\_\allowbreak{}digits}} by induction on $k$ with
$S\bmod 2^{k+1} = S\bmod 2^k + 2^k(\lfloor S/2^k\rfloor\bmod2)$.

\begin{theorem}[Qudit rule identities]\label{thm:lean-pathsum-q}
For $p,q,c,a,e\in\N$, $C\in\C$, and $f:\N\to\C$:
\begin{alignat*}{2}
&\text{\textnormal{\texttt{q\_\allowbreak{}elim}}:}\quad && \textstyle\sum_{y<p} C = p\,C,\\
&\text{\textnormal{\texttt{q\_\allowbreak{}hh}} }(0<p)\text{:}\quad && \textstyle\sum_{y<p}\omega_p^{yq} = p\,[q\equiv0\pmod p],\\
&\text{\textnormal{\texttt{q\_\allowbreak{}shift}} }(0<p,\ f(n+p)=f(n))\text{:}\quad && \textstyle\sum_{y<p} f(y+c)=\sum_{y<p} f(y),\\
&\text{\textnormal{\texttt{q\_\allowbreak{}gauss}}}\ (0<p,\ 2ac\equiv q,\ e+ac^2\equiv0 \pmod p)\text{:}\quad &&
 \textstyle\sum_{y<p}\omega_p^{ay^2+qy} = \omega_p^{e}\sum_{y<p}\omega_p^{ay^2}.
\end{alignat*}
\end{theorem}
\noindent \textnormal{\texttt{q\_\allowbreak{}gauss}} completes the square: each summand equals
$\omega_p^e\,\omega_p^{a(y+c)^2}$ (checked in \textnormal{\texttt{ZMod p}}), and \textnormal{\texttt{q\_\allowbreak{}shift}} removes the shift
by $c$ \cite{berndt1998gauss}.

\begin{definition}[Canonical-form records]\label{def:lean-pathsum-canon}
\textnormal{\texttt{PSCanon}} is the canonical-form data of a reduced qubit path sum, with fields
\textnormal{\texttt{rows : List (\allowbreak{}\ensuremath{\mathbb{N}} \ensuremath{\times} List (\allowbreak{}List \ensuremath{\mathbb{N}}))}} (pivot, monomials of the rest),
\textnormal{\texttt{phase : List (\allowbreak{}List \ensuremath{\mathbb{N}} \ensuremath{\times} \ensuremath{\mathbb{Q}})}} (monomial, coefficient in turns),
\textnormal{\texttt{const : List \ensuremath{\mathbb{Q}}}} (coefficients over $\omega_N^0,\dots,\omega_N^{\varphi(N)-1}$), and
\textnormal{\texttt{n : \ensuremath{\mathbb{N}}}} ($N$). \textnormal{\texttt{QCanon}} is the qudit analogue over $\Zd{p}$, with
\textnormal{\texttt{rows : List (\allowbreak{}\ensuremath{\mathbb{N}} \ensuremath{\times} List (\allowbreak{}List (\allowbreak{}\ensuremath{\mathbb{N}} \ensuremath{\times} \ensuremath{\mathbb{N}}) \ensuremath{\times} \ensuremath{\mathbb{N}}))}}, \textnormal{\texttt{phase : List (\allowbreak{}List (\allowbreak{}\ensuremath{\mathbb{N}} \ensuremath{\times} \ensuremath{\mathbb{N}}) \ensuremath{\times} \ensuremath{\mathbb{N}})}}
(monomials as (variable, exponent) pairs, coefficients mod $p$), \textnormal{\texttt{const : List \ensuremath{\mathbb{Q}}}},
\textnormal{\texttt{n : \ensuremath{\mathbb{N}}}}, and \textnormal{\texttt{p : \ensuremath{\mathbb{N}}}}. Both are plain structures; \textnormal{\texttt{export\_\allowbreak{}pair}} proves the equality of
the two sides' records by \textnormal{\texttt{rfl}}.
\end{definition}

\section{Qubit path sums: \texorpdfstring{\file{PathSumSem}}{PathSumSem}}\label{sec:lean-pathsum-sem}

\file{PathSumSem} (imports \file{PathSum}, \file{Ball}) defines the syntax and value of a qubit path
sum, the polynomial normal forms, Boolean substitution, and the steps \textnormal{\texttt{perm}}, \textnormal{\texttt{elim}},
\textnormal{\texttt{hhCons}}, \textnormal{\texttt{hhPivot}}, \textnormal{\texttt{canon}}, \textnormal{\texttt{gather}}, \textnormal{\texttt{bcanon}}, and the final
comparison.

\subsection{Polynomials and their values}

\begin{definition}[Monomials and polynomials]\label{def:lean-pathsum-poly}
\begin{itemize}
\item \textnormal{\texttt{Mono := List \ensuremath{\mathbb{N}}}}, a \term{monomial}: the product of its variables (repeats allowed).
\item \textnormal{\texttt{Poly := List (\allowbreak{}Mono \ensuremath{\times} \ensuremath{\mathbb{N}})}}, a \term{phase polynomial}: terms $(m,k)$ worth $k/M$ turns,
  where $M$ is the \term{phase modulus} passed to every value function.
\item \textnormal{\texttt{F2 := List Mono}}, an \term{$\mathbb F_2$ polynomial}: the sum of its monomials; the
  empty monomial \textnormal{\texttt{[]}} is the constant $1$ and the empty list is $0$.
\end{itemize}
For $\sigma:\N\to\Zd2$ the values are
\begin{align*}
\textnormal{\texttt{mz \ensuremath{\sigma} m}} &= \textstyle\prod_{v\in m}\sigma(v)\in\Zd2, &
\textnormal{\texttt{fz \ensuremath{\sigma} Q}} &= \textstyle\sum_{m\in Q}\textnormal{\texttt{mz \ensuremath{\sigma} m}}\in\Zd2,\\
\textnormal{\texttt{pn \ensuremath{\sigma} P}} &= \textstyle\sum_{(m,k)\in P} k\,\lfloor\textnormal{\texttt{mz \ensuremath{\sigma} m}}\rfloor\in\Z .
\end{align*}
The phase of $P$ is $\textnormal{\texttt{pn \ensuremath{\sigma} P}}/M$ turns. \textnormal{\texttt{peq M a b}} is the proposition
$\exists n\in\Z,\ a = b + nM$ (congruence mod $M$); \textnormal{\texttt{peq\_\allowbreak{}refl}}, \textnormal{\texttt{peq\_\allowbreak{}trans}},
\textnormal{\texttt{peq\_\allowbreak{}add}}, \textnormal{\texttt{peq\_\allowbreak{}mul}} (multiplication by $c\in\Z$), and \textnormal{\texttt{peq\_\allowbreak{}symm}} (in
\file{PathSumRules}) hold, and \textnormal{\texttt{ex\_\allowbreak{}peq}} states that $a\equiv b\pmod M$ with $0<M$ implies
$\operatorname{ex}(a/M)=\operatorname{ex}(b/M)$.
\end{definition}

\noindent The value lemmas are \textnormal{\texttt{pn\_\allowbreak{}nil}} ($0$), \textnormal{\texttt{pn\_\allowbreak{}cons}}, \textnormal{\texttt{pn\_\allowbreak{}append}}
(additivity), \textnormal{\texttt{mz\_\allowbreak{}append}} ($\textnormal{\texttt{mz \ensuremath{\sigma} (\allowbreak{}a ++ b)}} = \textnormal{\texttt{mz \ensuremath{\sigma} a}}\cdot\textnormal{\texttt{mz \ensuremath{\sigma} b}}$),
\textnormal{\texttt{fz\_\allowbreak{}cons}}, \textnormal{\texttt{fz\_\allowbreak{}append}}, \textnormal{\texttt{pn\_\allowbreak{}perm}}, and \textnormal{\texttt{fz\_\allowbreak{}perm}} (invariance under list
permutation). The update lemmas are:
\begin{itemize}
\item \textnormal{\texttt{mz\_\allowbreak{}update\_\allowbreak{}not}}: $z\notin m\Rightarrow \textnormal{\texttt{mz}} (\sigma[z\mapsto w])\,m = \textnormal{\texttt{mz}}\,\sigma\,m$;
\item \textnormal{\texttt{mz\_\allowbreak{}update\_\allowbreak{}in}}: $z\in m\Rightarrow \textnormal{\texttt{mz}}(\sigma[z\mapsto w])\,m = w\cdot\textnormal{\texttt{mz}}\,\sigma\,(m \text{ with every } z \text{ removed})$
  (the filter \textnormal{\texttt{m.\allowbreak{}filter (\allowbreak{}\ensuremath{\cdot} != z)}}; repeats of $z$ collapse by $x^2=x$);
\item \textnormal{\texttt{pn\_\allowbreak{}update\_\allowbreak{}not}}, \textnormal{\texttt{fz\_\allowbreak{}update\_\allowbreak{}not}}: the same invariance for every term of a phase or
  $\mathbb F_2$ polynomial that omits $z$.
\end{itemize}

\subsection{Bitmask monomials}

Normalization of large polynomials runs on \term{masks}: a monomial over an ascending variable list
$vs = [v_0,v_1,\dots]$ is the natural number whose bit $i$ is set when $v_i\in m$.

\begin{definition}[Masks]\label{def:lean-pathsum-masks}
\begin{align*}
\textnormal{\texttt{toM [] m}} &= 0, & \textnormal{\texttt{toM (\allowbreak{}v :: vs) m}} &= [v\in m] + 2\,\textnormal{\texttt{toM vs m}},\\
\textnormal{\texttt{ofM [] k}} &= [\,], & \textnormal{\texttt{ofM (\allowbreak{}v :: vs) k}} &=
 \begin{cases} v :: \textnormal{\texttt{ofM vs (\allowbreak{}k/2)}} & k\bmod 2 = 1\\ \textnormal{\texttt{ofM vs (\allowbreak{}k/2)}} & \text{else.}\end{cases}
\end{align*}
\textnormal{\texttt{mleM a b}} is \textnormal{\texttt{true}} when $a=b$; otherwise, with $d = a\oplus b$ and $t$ the lowest set
bit of $d$ (computed as \textnormal{\texttt{d \&\&\& (\allowbreak{}d \^{}\^{}\^{} (\allowbreak{}d -{} 1))}}), it is $a<t$ when bit $t$ of $a$ is clear and
$t\le b$ when it is set. On the monomials of two masks over one ascending list this is the
lexicographic order \textnormal{\texttt{mle}} below.
\end{definition}

\noindent Lemmas: \textnormal{\texttt{mem\_\allowbreak{}ofM\_\allowbreak{}toM}} ($w\in\textnormal{\texttt{ofM vs (\allowbreak{}toM vs m)}}\iff w\in vs\wedge w\in m$),
\textnormal{\texttt{or\_\allowbreak{}div\_\allowbreak{}two}} ($(a \mathbin{|} b)/2 = a/2 \mathbin{|} b/2$), \textnormal{\texttt{mem\_\allowbreak{}ofM\_\allowbreak{}or}},
\textnormal{\texttt{ofM\_\allowbreak{}zero}}, \textnormal{\texttt{mz\_\allowbreak{}dedup}}, \textnormal{\texttt{mz\_\allowbreak{}congr}} (monomials with the same variable set have the
same value), \textnormal{\texttt{mz\_\allowbreak{}ofM\_\allowbreak{}toM}} (when every variable of $m$ is in $vs$), and \textnormal{\texttt{mz\_\allowbreak{}ofM\_\allowbreak{}or}}
($\textnormal{\texttt{mz \ensuremath{\sigma} (\allowbreak{}ofM vs (\allowbreak{}a ||| b))}} = \textnormal{\texttt{mz \ensuremath{\sigma} (\allowbreak{}ofM vs a)}}\cdot\textnormal{\texttt{mz \ensuremath{\sigma} (\allowbreak{}ofM vs b)}}$: a bitwise
or of masks is a product of monomials).

\subsection{Normal forms}

\begin{definition}[Sorting]\label{def:lean-pathsum-sort}
\textnormal{\texttt{mergeN le n xs ys}} merges two lists by the Boolean order \textnormal{\texttt{le}}, comparing at most $n$
heads and appending the rest; \textnormal{\texttt{halve l}} splits $l$ into the entries at even and at odd
positions; \textnormal{\texttt{msortN le n l}} is merge sort at most $n$ levels deep:
$\textnormal{\texttt{msortN le (\allowbreak{}n+1) (\allowbreak{}a :: b :: l)}} = \textnormal{\texttt{mergeN le (\allowbreak{}|l|+2)}}\,(\textnormal{\texttt{msortN le n}}\,h_1)\,(\textnormal{\texttt{msortN le n}}\,h_2)$
with $(h_1,h_2)=\textnormal{\texttt{halve (\allowbreak{}a :: b :: l)}}$, lists of length $\le1$ returned unchanged and
\textnormal{\texttt{msortN le 0 l = l}}. \textnormal{\texttt{sortB le l := msortN le |l| l}}. The bounds make every function
structurally recursive, so the kernel reduces them under \textnormal{\texttt{decide}}. \textnormal{\texttt{mergeN\_\allowbreak{}perm}},
\textnormal{\texttt{halve\_\allowbreak{}perm}}, \textnormal{\texttt{msortN\_\allowbreak{}perm}}, and \textnormal{\texttt{sortB\_\allowbreak{}perm}} state that each output is a
permutation of its input; no lemma states sortedness.
\end{definition}

\begin{definition}[Insertion and merging]\label{def:lean-pathsum-ins}
\begin{itemize}
\item \textnormal{\texttt{ins M t P}} adds term $t=(m,k)$ to the first term of $P$ with the same monomial (list
  equality), coefficient $(u_2+k)\bmod M$, dropping the term when that is $0$; when no term
  matches, it appends $(m, k\bmod M)$, or nothing when $k\bmod M=0$.
\item \textnormal{\texttt{insH M t P}} compares $t$ with the head term of $P$ only: on equal monomials it merges
  as \textnormal{\texttt{ins}}; otherwise it conses $(m,k\bmod M)$ (or nothing when $k\bmod M = 0$).
\item \textnormal{\texttt{insHM M t P}} is \textnormal{\texttt{insH}} on mask terms \textnormal{\texttt{(\allowbreak{}\ensuremath{\mathbb{N}} \ensuremath{\times} \ensuremath{\mathbb{N}})}}.
\item \textnormal{\texttt{insU v l}} inserts $v$ into an ascending list without repeats; \textnormal{\texttt{vlist ms}} is the
  ascending repeat-free list of the variables of a list of monomials
  (\textnormal{\texttt{foldr}} of \textnormal{\texttt{insU}}); \textnormal{\texttt{mem\_\allowbreak{}vlist}} states that it contains every variable.
\item \textnormal{\texttt{squash l}} removes adjacent duplicates; \textnormal{\texttt{mnorm m := squash (\allowbreak{}insertionSort (\allowbreak{}\ensuremath{\le}) m)}}
  is a monomial sorted without repeats ($\textnormal{\texttt{mz\_\allowbreak{}mnorm}}$: same value).
\item \textnormal{\texttt{mle}} is the lexicographic order on monomials:
  $\textnormal{\texttt{mle [] \_\allowbreak{}}} = \textnormal{\texttt{true}}$, $\textnormal{\texttt{mle (\allowbreak{}\_\allowbreak{} :: \_\allowbreak{}) []}} = \textnormal{\texttt{false}}$,
  $\textnormal{\texttt{mle (\allowbreak{}a :: l) (\allowbreak{}b :: r)}} = (a<b)\vee(a=b\wedge \textnormal{\texttt{mle l r}})$.
\item \textnormal{\texttt{cancelH m Q}} cancels $m$ against the head monomial of $Q$ when equal (removing both),
  else conses; \textnormal{\texttt{cancelHM}} is the mask version.
\item \textnormal{\texttt{conv vs (\allowbreak{}k,\allowbreak{} c) := (\allowbreak{}ofM vs k,\allowbreak{} c)}}; \textnormal{\texttt{pnM \ensuremath{\sigma} vs P := pn \ensuremath{\sigma} (\allowbreak{}P.\allowbreak{}map (\allowbreak{}conv vs))}}.
\end{itemize}
\end{definition}

\begin{definition}[Normalization functions]\label{def:lean-pathsum-norm}
\begin{align*}
\textnormal{\texttt{normP M P}} &= \textnormal{\texttt{foldr (\allowbreak{}insH M) []}}\,\bigl(\textnormal{\texttt{sortB}}_{\textnormal{\texttt{mle}}}\,[(\textnormal{\texttt{mnorm}}\,m, k) : (m,k)\in P]\bigr),\\
\textnormal{\texttt{normPM M P}} &= \textnormal{\texttt{foldr (\allowbreak{}insHM M) []}}\,(\textnormal{\texttt{sortB}}_{\textnormal{\texttt{mleM}}}\,P),\\
\textnormal{\texttt{normFL Q}} &= \textnormal{\texttt{foldr cancelH []}}\,(\textnormal{\texttt{sortB}}_{\textnormal{\texttt{mle}}}\,(Q.\textnormal{\texttt{map mnorm}})),\\
\textnormal{\texttt{normFD Q}} &= \textnormal{\texttt{pickX}}\,(\textnormal{\texttt{xorF Q}},\ \textnormal{\texttt{orF Q}},\ |Q|,\ \textnormal{\texttt{sortB}}_{\textnormal{\texttt{mleM}}}\,[k\le\textnormal{\texttt{orF Q}} : \text{bit } k \text{ of } \textnormal{\texttt{xorF Q}} \text{ set}]),\\
\textnormal{\texttt{normFM Q}} &= \begin{cases}\textnormal{\texttt{normFD Q}} & \textnormal{\texttt{orF Q}} < 4|Q|+64\\
   \textnormal{\texttt{foldr cancelHM []}}\,(\textnormal{\texttt{sortB}}_{\textnormal{\texttt{mleM}}}\,Q) & \text{else,}\end{cases}\\
\textnormal{\texttt{normF Q}} &= \begin{cases}\textnormal{\texttt{normFL Q}} & |Q|<32\\
   (\textnormal{\texttt{normFM}}\,(Q.\textnormal{\texttt{map (\allowbreak{}toM vs)}})).\textnormal{\texttt{map (\allowbreak{}ofM vs)}},\ vs = \textnormal{\texttt{vlist Q}} & \text{else,}\end{cases}
\end{align*}
where $\textnormal{\texttt{xorF Q}} = \bigoplus_{k\in Q} 2^k$ (the set of masks of odd multiplicity, as a bit
set) and $\textnormal{\texttt{orF Q}} = \bigvee_{k\in Q} k$ (an upper bound of every mask of $Q$, \textnormal{\texttt{le\_\allowbreak{}orF}}).
Equal adjacent terms after sorting merge (\textnormal{\texttt{normP}}) or cancel in pairs (\textnormal{\texttt{normF}}), so on
sorted input the output lists each monomial once with a nonzero coefficient.
\end{definition}

\begin{definition}[Submask enumeration: \texttt{pickX}]\label{def:lean-pathsum-pickx}
\begin{itemize}
\item \lean!subsM : List ℕ → List ℕ!: $\textnormal{\texttt{subsM []}}=[0]$ and
  $\textnormal{\texttt{subsM (b :: R)}} = 0 :: \bigl([b\mathbin{|}s : s\in\textnormal{\texttt{subsM R}}]\mathbin{+\!\!+}(\textnormal{\texttt{subsM R}}).\textnormal{\texttt{tail}}\bigr)$,
  the number $0$ and the bitwise or of every nonempty sublist of $bs$; for $bs$ a list of
  ascending powers of two the output is in \lean!mleM! order.
\item \lean!bitsOf (m : ℕ) : List ℕ! lists $2^i$ for the set bits $i\le\lfloor\log_2 m\rfloor$ of $m$,
  ascending.
\item \lean!fromX (x b : ℕ) : List ℕ! is \textnormal{\texttt{subsM (bitsOf b)}} filtered to the $k$ with bit
  $k$ of $x$ set: the submasks of $b$ that $x$ marks, in \lean!mleM! order.
\item \lean!pickX (x : Unit → ℕ) (b n : ℕ) (R : List ℕ) : List ℕ! is \textnormal{\texttt{fromX (x ()) b}}
  when $b<4n+64$ and $\textnormal{\texttt{xorF (fromX (x ()) b)}} = x()$, and $R$ otherwise. The thunk $x$ is
  evaluated only when $b<4n+64$.
\end{itemize}
In \lean!normFD!, $x=\textnormal{\texttt{xorF Q}}$, $b=\textnormal{\texttt{orF Q}}$, $n=|Q|$: the set masks are
read off the submasks of \lean!orF Q! in order, and the merge-sorted bit list is the fallback.
\lean!fz_xorF_eq!: two mask lists with equal \lean!xorF! have equal \lean!fz! values over every
$vs$; \lean!fz_pickX!: if $x()=\textnormal{\texttt{xorF Q}}$ and $R$ has the value of $Q$, then
\textnormal{\texttt{pickX x b n R}} has the value of $Q$ (the test $\textnormal{\texttt{xorF (fromX \ldots)}} = x()$
is the side condition the proof reads).
\end{definition}

\begin{proposition}[Normalization preserves values]\label{prop:lean-pathsum-norm}
For all $\sigma$, $M$, $vs$:
$\textnormal{\texttt{pn \ensuremath{\sigma} (\allowbreak{}normP M P)}} \equiv \textnormal{\texttt{pn \ensuremath{\sigma} P}}\pmod M$ (\textnormal{\texttt{pn\_\allowbreak{}normP}}),
$\textnormal{\texttt{pnM \ensuremath{\sigma} vs (\allowbreak{}normPM M P)}}\equiv\textnormal{\texttt{pnM \ensuremath{\sigma} vs P}}\pmod M$ (\textnormal{\texttt{pnM\_\allowbreak{}normPM}}),
$\textnormal{\texttt{fz \ensuremath{\sigma} (\allowbreak{}normF Q)}} = \textnormal{\texttt{fz \ensuremath{\sigma} Q}}$ (\textnormal{\texttt{fz\_\allowbreak{}normF}}), and
$\textnormal{\texttt{fz \ensuremath{\sigma} (\allowbreak{}(\allowbreak{}normFM Q).\allowbreak{}map (\allowbreak{}ofM vs))}} = \textnormal{\texttt{fz \ensuremath{\sigma} (\allowbreak{}Q.\allowbreak{}map (\allowbreak{}ofM vs))}}$ (\textnormal{\texttt{fz\_\allowbreak{}normFM}}).
\end{proposition}
\noindent The supporting lemmas are \textnormal{\texttt{peq\_\allowbreak{}mod}} ($(k\bmod M)z\equiv kz$), \textnormal{\texttt{pn\_\allowbreak{}ins}},
\textnormal{\texttt{pn\_\allowbreak{}insH}}, \textnormal{\texttt{pn\_\allowbreak{}foldr\_\allowbreak{}insH}}, \textnormal{\texttt{pnM\_\allowbreak{}insHM}}, \textnormal{\texttt{pnM\_\allowbreak{}foldr\_\allowbreak{}insHM}}, \textnormal{\texttt{fz\_\allowbreak{}cancelH}},
\textnormal{\texttt{fz\_\allowbreak{}foldr\_\allowbreak{}cancelH}}, \textnormal{\texttt{fz\_\allowbreak{}cancelHM}}, \textnormal{\texttt{fz\_\allowbreak{}foldr\_\allowbreak{}cancelHM}}, \textnormal{\texttt{fz\_\allowbreak{}toM}} (masks over a
list containing all variables round-trip), and, for the bit-set path, \textnormal{\texttt{sum\_\allowbreak{}xorF}}: for a
repeat-free candidate list containing every element of $Q$, the sum of $f$ over the candidates
whose bit is set in \textnormal{\texttt{xorF Q}} equals $\sum_{k\in Q} f(k)$ in $\Zd2$ (proved with
\textnormal{\texttt{testBit\_\allowbreak{}xorF}}, \textnormal{\texttt{sum\_\allowbreak{}filter\_\allowbreak{}ite}}, \textnormal{\texttt{sum\_\allowbreak{}ite\_\allowbreak{}eq\_\allowbreak{}none}}, \textnormal{\texttt{sum\_\allowbreak{}ite\_\allowbreak{}eq\_\allowbreak{}self}}).

\subsection{The Boolean lift and substitution}

\begin{definition}[Lift]\label{def:lean-pathsum-lift}
The \term{lift} of an $\mathbb F_2$ polynomial is the integer phase polynomial with the same
$\{0,1\}$ value, $-2$ written $M-2$:
\begin{align*}
\textnormal{\texttt{lift M []}} &= [\,],\\
\textnormal{\texttt{lift M (\allowbreak{}m :: B)}} &= (m,1) :: \bigl(L \mathbin{+\!\!+} [(m\mathbin{+\!\!+}u, (M-2)\,k) : (u,k)\in L]\bigr),\quad L=\textnormal{\texttt{lift M B}}.
\end{align*}
\textnormal{\texttt{pn\_\allowbreak{}lift}} ($2\le M$): $\textnormal{\texttt{pn \ensuremath{\sigma} (\allowbreak{}lift M B)}}\equiv\lfloor\textnormal{\texttt{fz \ensuremath{\sigma} B}}\rfloor\pmod M$, by
$\lfloor x+y\rfloor = \lfloor x\rfloor+\lfloor y\rfloor-2\lfloor x\rfloor\lfloor y\rfloor$.
\textnormal{\texttt{liftNM M c B}} (masks) is $c$ times the lift, normalized by \textnormal{\texttt{normPM}} after each mask:
$\textnormal{\texttt{liftNM M c (\allowbreak{}m :: B)}} = \textnormal{\texttt{normPM M (\allowbreak{}(\allowbreak{}m,\allowbreak{} c) :: (\allowbreak{}L ++ L.\allowbreak{}map (\allowbreak{}fun t => (\allowbreak{}m ||| t.\allowbreak{}1,\allowbreak{} (\allowbreak{}M -{} 2) * t.\allowbreak{}2))))}}$.
\textnormal{\texttt{liftN M c B}} is $[(m,c): m\in B]$ when $2c\equiv 0\pmod M$; otherwise, with
$vs=\textnormal{\texttt{vlist B}}$, it is the dense lift \textnormal{\texttt{liftD M c vs B}} when
$2^{|vs|}\le16|B|$, and
\textnormal{\texttt{(\allowbreak{}liftNM M c (\allowbreak{}B.\allowbreak{}map (\allowbreak{}toM vs))).\allowbreak{}map (\allowbreak{}conv vs)}} else.
\textnormal{\texttt{pn\_\allowbreak{}liftN}} ($2\le M$): $\textnormal{\texttt{pn \ensuremath{\sigma} (\allowbreak{}liftN M c B)}}\equiv c\lfloor\textnormal{\texttt{fz \ensuremath{\sigma} B}}\rfloor\pmod M$;
the first case is \textnormal{\texttt{pn\_\allowbreak{}linlift}} (when $2c\equiv0$ the cross terms vanish mod $M$),
the second \lean!pn_liftD!.
\end{definition}

\begin{definition}[The dense lift: \texttt{Tb}, \texttt{liftD}]\label{def:lean-pathsum-liftd}
\lean!Tb! is the inductive type of binary tables with constructors \lean!leaf (k : ℕ)! and
\lean!node (lo hi : Tb)!; over a variable list $v::vs$, \lean!node lo hi! holds \lean!lo! where
$v=0$ and \lean!hi! where $v=1$. \lean!Tb.dep n t! states that $t$ is a complete tree of depth $n$.
\lean!Tb.at σ vs t! reads the leaf at $\sigma$ (a leaf answers at once; a \lean!node! over $[\,]$
reads \lean!lo!). The operations are:
\begin{itemize}
\item \lean!Tb.zip f!: leafwise $f$ on two trees of one shape (the first tree when the shapes
  differ); \lean!Tb.scale c!: every leaf times $c$;
\item \lean!Tb.xorN a b = node a (zip (x+y mod 2) a b)!; \lean!Tb.subN M a b = node a (zip (fun x y => x + (M - y % M)) b a)!;
\item \lean!tabF : List ℕ → F2 → Tb!: $\textnormal{\texttt{tabF [] B}}=\textnormal{\texttt{leaf (|B| mod 2)}}$,
  $\textnormal{\texttt{tabF (v :: vs) B}} = \textnormal{\texttt{xorN (tabF vs A}}_0\textnormal{\texttt{) (tabF vs A}}_1\textnormal{\texttt{)}}$ with
  $(A_1,A_0)=\textnormal{\texttt{splitAt v B}}$: the $0/1$ values of $B$ on the assignments of $vs$;
\item \lean!Tb.mob M!: $\textnormal{\texttt{mob (leaf k)}}=\textnormal{\texttt{leaf k}}$,
  $\textnormal{\texttt{mob (node l h)}}=\textnormal{\texttt{subN M (mob l) (mob h)}}$, the multilinear (M\"obius)
  coefficients mod $M$;
\item \lean!Tb.poly M vs t!: the terms of the leaves with $k\bmod M\ne0$, each monomial the
  variables at $1$ (\lean!node l h! over $v::vs$ gives the terms of $l$, then those of $h$ with $v$
  prepended);
\item \lean!liftD M c vs B = Tb.poly M vs (Tb.mob M (Tb.scale c (tabF vs B)))!.
\end{itemize}
\lean!pn_liftD! ($0<M$, every variable of $B$ in $vs$):
$\textnormal{\texttt{pn \ensuremath{\sigma} (liftD M c vs B)}}\equiv c\lfloor\textnormal{\texttt{fz \ensuremath{\sigma} B}}\rfloor\pmod M$.
It rests on \lean!dep_tabF!, \lean!Tb.dep_zip!, \lean!Tb.dep_scale!, \lean!Tb.dep_mob! (depth
$|vs|$ is kept), \lean!at_tabF! (the table holds $[\textnormal{\texttt{fz \ensuremath{\sigma} B}}\ne0]$), \lean!zn_empty!,
\lean!Tb.at_zip!, \lean!Tb.at_scale!, \lean!pn_mapcons!, \lean!pn_polyLeaf!, \lean!peq_natmod!,
\lean!pn_poly_sub! and \lean!pn_poly_mob! (the polynomial of the M\"obius table takes the table's
value mod $M$).
\end{definition}

\begin{definition}[Substitution]\label{def:lean-pathsum-subst}
For a variable $z$ and $B\in\mathbb F_2[x]$:
\begin{align*}
\textnormal{\texttt{substP M z B P}} &= \textstyle\bigcup_{(m,k)\in P}
  \begin{cases}[(m\setminus z \mathbin{+\!\!+} u, c) : (u,c)\in \textnormal{\texttt{liftN M k B}}] & z\in m\\ [(m,k)] & \text{else,}\end{cases}\\
\textnormal{\texttt{substF z B Q}} &= \textstyle\bigcup_{m\in Q}
  \begin{cases}[m\setminus z\mathbin{+\!\!+}u : u\in B] & z\in m\\ [m] & \text{else,}\end{cases}
\end{align*}
($\bigcup$ is \textnormal{\texttt{flatMap}}, $m\setminus z$ is \textnormal{\texttt{m.\allowbreak{}filter (\allowbreak{}\ensuremath{\cdot} != z)}}).
\textnormal{\texttt{substFV vs z B Q}} is \textnormal{\texttt{normF (\allowbreak{}substF z B Q)}} computed on masks over
$vs$: each monomial holding $z$ becomes the masks $\textnormal{\texttt{toM vs}}(m\setminus z)\mathbin{|}u$ for
$u$ in the masks $B_m$ of $B$, every other monomial its mask; the list goes through
\textnormal{\texttt{pickX}} with $x=\textnormal{\texttt{xorSub z vs B}}_m\textnormal{\texttt{ Q}}$,
$b=\textnormal{\texttt{orF}}(Q\text{'s masks})\mathbin{|}\textnormal{\texttt{orF}}\,B_m$, $n=|Q|(|B|+1)$ and fallback
\textnormal{\texttt{normFM}} of the list, then \textnormal{\texttt{ofM vs}}. \textnormal{\texttt{xorSub}} folds the same masks
into one \lean!xorF! without building the list (\lean!xorSub_eq!, from \lean!xorF_append!).
\textnormal{\texttt{substFN z B Q = substFV (vlist (Q ++ B)) z B Q}}.
\end{definition}

\begin{definition}[Substitution skipping normal forms]\label{def:lean-pathsum-substn}
\lean!ascB m! holds when the monomial is strictly ascending; \lean!normalF Q! when every monomial
is strictly ascending and $Q$ is strictly ascending by \lean!mle!; \lean!normalP M P! when moreover
every coefficient lies in $[1,M)$ (strictly ascending monomials, strictly ascending by \lean!mle!).
\begin{align*}
\textnormal{\texttt{substNF z B q}} &= \begin{cases} q & z\text{ occurs in no monomial of }q\text{ and }\textnormal{\texttt{normalF q}}\\ \textnormal{\texttt{normF (substF z B q)}} & \text{else,}\end{cases}\\
\textnormal{\texttt{substNP M z B P}} &= \begin{cases} P & z\notin\textnormal{\texttt{pvars P}}\text{ and }\textnormal{\texttt{normalP M P}}\\ \textnormal{\texttt{normP M (substP M z B P)}} & \text{else.}\end{cases}
\end{align*}
\end{definition}

\begin{proposition}[Substitution lemmas]\label{prop:lean-pathsum-subst}
With $\sigma' = \sigma[z\mapsto\textnormal{\texttt{fz \ensuremath{\sigma} B}}]$ and $2\le M$:
$\textnormal{\texttt{pn \ensuremath{\sigma} (\allowbreak{}substP M z B P)}}\equiv\textnormal{\texttt{pn \ensuremath{\sigma}' P}}\pmod M$ (\textnormal{\texttt{pn\_\allowbreak{}substP}}),
$\textnormal{\texttt{fz \ensuremath{\sigma} (\allowbreak{}substF z B Q)}} = \textnormal{\texttt{fz \ensuremath{\sigma}' Q}}$ (\textnormal{\texttt{fz\_\allowbreak{}substF}}), and
$\textnormal{\texttt{fz \ensuremath{\sigma} (\allowbreak{}substFN z B Q)}} = \textnormal{\texttt{fz \ensuremath{\sigma}' Q}}$ (\textnormal{\texttt{fz\_\allowbreak{}substFN}}),
$\textnormal{\texttt{fz \ensuremath{\sigma} (substFV vs z B Q)}} = \textnormal{\texttt{fz \ensuremath{\sigma}' Q}}$ when every variable of $Q$
and $B$ is in $vs$ (\lean!fz_substFV!),
$\textnormal{\texttt{fz \ensuremath{\sigma} (substNF z B q)}} = \textnormal{\texttt{fz \ensuremath{\sigma}' q}}$ (\lean!fz_substNF!), and
$\textnormal{\texttt{pn \ensuremath{\sigma} (substNP M z B P)}}\equiv\textnormal{\texttt{pn \ensuremath{\sigma}' P}}\pmod M$ (\lean!pn_substNP!).
\end{proposition}

\subsection{Sums over path variables}

\begin{definition}[Boolean sum]\label{def:lean-pathsum-bsum}
\[
\textnormal{\texttt{bsum [] f \ensuremath{\sigma}}} = f(\sigma),\qquad
\textnormal{\texttt{bsum (\allowbreak{}y :: ys) f \ensuremath{\sigma}}} = \textnormal{\texttt{bsum ys f}}\,\sigma[y\mapsto0]+\textnormal{\texttt{bsum ys f}}\,\sigma[y\mapsto1].
\]
\end{definition}
\noindent Lemmas: \textnormal{\texttt{bsum\_\allowbreak{}congr}} (pointwise equal summands), \textnormal{\texttt{bsum\_\allowbreak{}add}},
\textnormal{\texttt{bsum\_\allowbreak{}smul}}, \textnormal{\texttt{bsum\_\allowbreak{}out}} (for $y\notin ys$, an update at $y$ commutes into the summand),
\textnormal{\texttt{bsum\_\allowbreak{}perm}} (a permutation of the variable list leaves the sum unchanged), and
\textnormal{\texttt{bsum\_\allowbreak{}head}} (for $y\notin ys$,
$\textnormal{\texttt{bsum (\allowbreak{}y :: ys) f \ensuremath{\sigma}}} = \textnormal{\texttt{bsum ys}}\,(\tau\mapsto f(\tau[y\mapsto0])+f(\tau[y\mapsto1]))\,\sigma$).
\file{PathSumRules} adds \textnormal{\texttt{bsum\_\allowbreak{}append}} ($\textnormal{\texttt{bsum (\allowbreak{}xs ++ ys) f = bsum xs (\allowbreak{}bsum ys f)}}$) and
\textnormal{\texttt{bsum\_\allowbreak{}mul\_\allowbreak{}left}} (a factor invariant under updates of $ys$ leaves the sum).

\begin{definition}[Qubit path sum]\label{def:lean-pathsum-data}
\textnormal{\texttt{structure Data}} has fields \textnormal{\texttt{path : List \ensuremath{\mathbb{N}}}} (the path variables), \textnormal{\texttt{phase : Poly}},
\textnormal{\texttt{cons : List F2}} (constraints), and \textnormal{\texttt{rows : List F2}} (support rows); it derives
\textnormal{\texttt{DecidableEq}}. For $M\in\N$, $c\in\C$:
\begin{align*}
\textnormal{\texttt{summand M c s \ensuremath{\sigma}}} &= \Bigl[\textstyle\bigwedge_{q\in \text{cons}\mathbin{+\!\!+}\text{rows}} \textnormal{\texttt{fz \ensuremath{\sigma} q}} = 0\Bigr]\cdot c\cdot\operatorname{ex}\bigl(\textnormal{\texttt{pn \ensuremath{\sigma} s.\allowbreak{}phase}}/M\bigr),\\
\textnormal{\texttt{val M c s \ensuremath{\sigma}}} &= \textnormal{\texttt{bsum s.\allowbreak{}path (\allowbreak{}summand M c s) \ensuremath{\sigma}}}\\
  &= \sum_{x\in\{0,1\}^{\text{path}}}\textnormal{\texttt{summand M c s}}\,(\sigma[\text{path}\mapsto x]).
\end{align*}
Constraints and rows have the same semantics; the steps treat them differently (rows are
solved constraints in canonical form). \textnormal{\texttt{pvars P}} (\textnormal{\texttt{P.\allowbreak{}flatMap (\allowbreak{}\ensuremath{\cdot}.\allowbreak{}1)}}) and
\textnormal{\texttt{fvars C}} (\textnormal{\texttt{C.\allowbreak{}flatMap (\allowbreak{}fun q => q.\allowbreak{}flatMap id)}}) list the variables of a phase and of a
list of $\mathbb F_2$ polynomials. \textnormal{\texttt{not\_\allowbreak{}mem\_\allowbreak{}pvars}}, \textnormal{\texttt{not\_\allowbreak{}mem\_\allowbreak{}fvars}},
\textnormal{\texttt{holds\_\allowbreak{}update\_\allowbreak{}not}}, and \textnormal{\texttt{summand\_\allowbreak{}update\_\allowbreak{}not}} state the invariance of a summand under an
update at a variable absent from its phase, constraints, and rows.
\end{definition}

\noindent In the denotation of \cref{ch:semantics}, a path sum with boundary variables $\vec b$
denotes the matrix entry $\textnormal{\texttt{val M c s \ensuremath{\sigma}}}$ at the assignment $\sigma$ that sets $\vec b$; the
library never separates boundary from path variables beyond the list \textnormal{\texttt{path}}.

\subsection{Steps of \texorpdfstring{\file{PathSumSem}}{PathSumSem}}

Each step has a side condition \textnormal{\texttt{xOK}}, a transformer \textnormal{\texttt{xS}}, and a soundness theorem
\textnormal{\texttt{x\_\allowbreak{}sound}}. ``Head'' variables are the first entries of \textnormal{\texttt{s.\allowbreak{}path}}.

\paragraph{Permutation.} \textnormal{\texttt{permS l s := \{ s with path := l \}}}. \textnormal{\texttt{perm\_\allowbreak{}sound}}: if
\textnormal{\texttt{decide (\allowbreak{}s.\allowbreak{}path.\allowbreak{}Perm l) = true}} then $\textnormal{\texttt{val M c s \ensuremath{\sigma}}} = \textnormal{\texttt{val M c (\allowbreak{}permS l s) \ensuremath{\sigma}}}$.
\textnormal{\texttt{frontS x s := \{ s with path := x :: s.path.erase x \}}} moves $x$ to the front;
\lean!front_sound!: if \textnormal{\texttt{decide (x \ensuremath{\in} s.path) = true}} then the value is kept.

\paragraph{[Elim].} \textnormal{\texttt{elimOK s}} holds when $\text{path} = y :: ys$ with $y\notin ys$,
$y\notin\textnormal{\texttt{pvars phase}}$, and $y\notin\textnormal{\texttt{fvars (\allowbreak{}cons ++ rows)}}$. \textnormal{\texttt{elimS}} drops the head.
\begin{theorem}[\textnormal{\texttt{elim\_\allowbreak{}sound}}]
$\textnormal{\texttt{elimOK s}}\Rightarrow \textnormal{\texttt{val M c s \ensuremath{\sigma}}} = \textnormal{\texttt{val M (\allowbreak{}2 * c) (\allowbreak{}elimS s) \ensuremath{\sigma}}}$.
\end{theorem}

\paragraph{Quotients.} For a variable $y$:
$\textnormal{\texttt{yTerms y P}}$ is the terms of $P$ whose monomial contains $y$, $\textnormal{\texttt{rest y P}}$ the others,
$\textnormal{\texttt{quot y P}} = [m\setminus y : (m,k)\in\textnormal{\texttt{yTerms y P}}]\in\mathbb F_2[x]$, and
$\textnormal{\texttt{yq y P}} = [(m\setminus y,k) : (m,k)\in\textnormal{\texttt{yTerms y P}}]$. \textnormal{\texttt{halfOK M y P}} holds when
every term of $y$ has $2k=M$ (coefficient one half turn). Then
\begin{equation}\label{eq:lean-pathsum-half}
\textnormal{\texttt{pn \ensuremath{\sigma}[y\ensuremath{\mapsto}b] P}}/M = \lfloor b\rfloor\cdot\textstyle\sum_{m\in\textnormal{\texttt{quot y P}}}\lfloor\textnormal{\texttt{mz \ensuremath{\sigma} m}}\rfloor/2 + \textnormal{\texttt{pn \ensuremath{\sigma} (\allowbreak{}rest y P)}}/M
\end{equation}
(\textnormal{\texttt{pn\_\allowbreak{}half}}, $0<M$), and \textnormal{\texttt{pn\_\allowbreak{}yupd}} states the exact split
$\textnormal{\texttt{pn \ensuremath{\sigma}[y\ensuremath{\mapsto}b] P}} = \lfloor b\rfloor\,\textnormal{\texttt{pn \ensuremath{\sigma} (\allowbreak{}yq y P)}} + \textnormal{\texttt{pn \ensuremath{\sigma} (\allowbreak{}rest y P)}}$ for all
coefficients. \textnormal{\texttt{ex\_\allowbreak{}halfN}} gives $\operatorname{ex}(N/2) = 1$ or $-1$ by the parity of $N$, and
\textnormal{\texttt{pair\_\allowbreak{}half}} gives the pair sum over $y$ of a half-coefficient phase:
$c\,\operatorname{ex}(0\cdot N/2+r)+c\,\operatorname{ex}(1\cdot N/2+r) = 2[N\text{ even}]\,c\,\operatorname{ex}(r)$.
\textnormal{\texttt{pn\_\allowbreak{}allhalf}} and \textnormal{\texttt{zsum\_\allowbreak{}mz}} relate the half-coefficient phase to \textnormal{\texttt{fz}} of the
quotient.

\paragraph{[HH] on a constraint.} \textnormal{\texttt{hhConsOK M s}}: $\text{path}=y::ys$, $0<M$, $y\notin ys$,
\textnormal{\texttt{halfOK M y phase}}, $y\notin\textnormal{\texttt{fvars (\allowbreak{}cons ++ rows)}}$.
\textnormal{\texttt{hhConsS s}} has path $ys$, phase \textnormal{\texttt{rest y phase}}, constraints
\textnormal{\texttt{quot y phase :: cons}}, and the same rows.
\begin{theorem}[\textnormal{\texttt{hhCons\_\allowbreak{}sound}}]
$\textnormal{\texttt{hhConsOK M s}}\Rightarrow\textnormal{\texttt{val M c s \ensuremath{\sigma}}} = \textnormal{\texttt{val M (\allowbreak{}2 * c) (\allowbreak{}hhConsS s) \ensuremath{\sigma}}}$. This is
$\sum_{y}(-1)^{yQ} = 2[Q=0]$ with $Q$ the quotient.
\end{theorem}

\paragraph{[HH] on a path pivot.} \textnormal{\texttt{hhPivotOK M s}}: $\text{path} = y :: z :: ys$, $2\le M$,
$y\ne z$, $y\notin ys$, $z\notin ys$, \textnormal{\texttt{halfOK M y phase}}, $y\notin\textnormal{\texttt{fvars (\allowbreak{}cons ++ rows)}}$,
$[z]\in\textnormal{\texttt{quot y phase}}$, and $z$ in no other monomial of the quotient
(\textnormal{\texttt{z \ensuremath{\notin} (\allowbreak{}(\allowbreak{}quot y phase).\allowbreak{}erase [z]).\allowbreak{}flatMap id}}). With $B = (\textnormal{\texttt{quot y phase}}).\textnormal{\texttt{erase [z]}}$:
\textnormal{\texttt{hhPivotS M s}} has path $ys$, phase \textnormal{\texttt{substNP M z B (\allowbreak{}rest y phase)}}, and each
constraint and row $q$ replaced by \textnormal{\texttt{substNF z B q}} (\cref{def:lean-pathsum-substn}).
\begin{theorem}[\textnormal{\texttt{hhPivot\_\allowbreak{}sound}}]
$\textnormal{\texttt{hhPivotOK M s}}\Rightarrow\textnormal{\texttt{val M c s \ensuremath{\sigma}}} = \textnormal{\texttt{val M (\allowbreak{}2 * c) (\allowbreak{}hhPivotS M s) \ensuremath{\sigma}}}$:
$\sum_z\sum_y(-1)^{y(z+B)}g(z) = 2g(B)$ (lemma \textnormal{\texttt{pick}}:
$[0+v=0]\cdot2\,G(0)+[1+v=0]\cdot 2\,G(1) = 2G(v)$).
\end{theorem}

\paragraph{Canonical-form pivot.} \textnormal{\texttt{pivotable v q}} holds when $[v]\in q$ and $v$ occurs in no
other monomial of $q$. \textnormal{\texttt{canonOK v s}}: some constraint is pivotable at $v$ (\textnormal{\texttt{find?}}) and
$v\notin\text{path}$. With $q$ the first such constraint and $B = q.\textnormal{\texttt{erase [v]}}$:
\textnormal{\texttt{canonS M v s}} keeps the path, has phase \textnormal{\texttt{normP M (\allowbreak{}substP M v B phase)}}, the
constraints \textnormal{\texttt{(\allowbreak{}cons.\allowbreak{}erase q).\allowbreak{}map (\allowbreak{}fun r => normF (\allowbreak{}substF v B r))}}, and the rows
\textnormal{\texttt{rows.\allowbreak{}map (\allowbreak{}fun r => normF (\allowbreak{}substF v B r))}} followed by \textnormal{\texttt{normF q}}.
\begin{theorem}[\textnormal{\texttt{canon\_\allowbreak{}sound}}]
$2\le M$, $\textnormal{\texttt{canonOK v s}}\Rightarrow\textnormal{\texttt{val M c s \ensuremath{\sigma}}} = \textnormal{\texttt{val M c (\allowbreak{}canonS M v s) \ensuremath{\sigma}}}$.
On the support of $q$, $v = B$ (\textnormal{\texttt{z2\_\allowbreak{}eq\_\allowbreak{}of\_\allowbreak{}add}}: $a+b=0\Rightarrow a=b$), so $v$ may be
substituted everywhere else; the constraint becomes the row $q$.
\end{theorem}

\paragraph{The final comparison.} $\textnormal{\texttt{constK P}}$ is the sum of the coefficients of the terms with
the empty monomial; $\textnormal{\texttt{strip P}}$ is the other terms; \textnormal{\texttt{pn\_\allowbreak{}strip}} gives
$\textnormal{\texttt{pn \ensuremath{\tau} P}} = \textnormal{\texttt{constK P}} + \textnormal{\texttt{pn \ensuremath{\tau} (\allowbreak{}strip P)}}$. \textnormal{\texttt{finalOK sL sR}} holds when both
paths are empty, every constraint of both sides is the zero polynomial \textnormal{\texttt{[]}}, the row lists are
equal, and \textnormal{\texttt{strip sL.\allowbreak{}phase = strip sR.\allowbreak{}phase}}.
\begin{theorem}[\textnormal{\texttt{final\_\allowbreak{}sound}}]
If $0<M$, \textnormal{\texttt{finalOK sL sR = true}}, $\textnormal{\texttt{constK sL.\allowbreak{}phase}} = k_L$,
$\textnormal{\texttt{constK sR.\allowbreak{}phase}} = k_R$, and $c_L\,\omega_M^{k_L} = c_R\,\omega_M^{k_R}$, then
$\textnormal{\texttt{val M cL sL \ensuremath{\sigma}}} = \textnormal{\texttt{val M cR sR \ensuremath{\sigma}}}$ for every $\sigma$.
\end{theorem}

\paragraph{Boolean elimination: \textnormal{\texttt{gather}}.} $\textnormal{\texttt{f2mul a b}} = [m\mathbin{+\!\!+}u: m\in a, u\in b]$
(\textnormal{\texttt{fz\_\allowbreak{}f2mul}}: the value is the product), \textnormal{\texttt{f2mulM}} the mask version with bitwise or, and
\textnormal{\texttt{xorMul a b}} folds $2^{m\mathbin{|}u}$ by XOR over $m\in a$, $u\in b$
($=\textnormal{\texttt{xorF (f2mulM a b)}}$, \lean!xorMul_eq!), and
$\textnormal{\texttt{mulM a b}} = \textnormal{\texttt{pickX}}(\textnormal{\texttt{xorMul a b}},\ \textnormal{\texttt{orF a ||| orF b}},\ |a|\,|b|,\ \textnormal{\texttt{normFM (f2mulM a b)}})$
(\lean!fz_mulM!: the value is the product).
$\textnormal{\texttt{f2mulN a b}} = \textnormal{\texttt{normF (\allowbreak{}f2mul a b)}}$ on masks over \textnormal{\texttt{vlist (\allowbreak{}a ++ b)}},
computed as \textnormal{\texttt{mulM}} of the masks.
$\textnormal{\texttt{gatherM vs C}} = \textnormal{\texttt{foldr}}\,(c,acc\mapsto\textnormal{\texttt{if c = [] then acc else mulM (\allowbreak{}0 :: c') acc}})\,[0]$ with
$c'$ the masks of $c$, i.e.\ $\prod_{c\in C}(1+c)$ normalized after each factor, a zero constraint
skipped;
\textnormal{\texttt{gatherPoly C}} is its monomial form over \textnormal{\texttt{vlist (\allowbreak{}C.\allowbreak{}flatMap id)}}.
\textnormal{\texttt{fz\_\allowbreak{}gather}}: $\textnormal{\texttt{fz \ensuremath{\tau} (\allowbreak{}gatherPoly C)}} = 1\iff\forall q\in C,\ \textnormal{\texttt{fz \ensuremath{\tau} q}} = 0$.
\textnormal{\texttt{gatherOK s}} holds when \textnormal{\texttt{s.\allowbreak{}path = []}}.
\[
\textnormal{\texttt{gatherS s}} = \{s \text{ with cons} := [\textnormal{\texttt{normF (\allowbreak{}[] :: gatherPoly s.\allowbreak{}cons)}}]\},
\]
the single constraint $1+\prod_{c}(1+c)$. \textnormal{\texttt{gather\_\allowbreak{}sound}}:
$\textnormal{\texttt{val M c s \ensuremath{\sigma}}} = \textnormal{\texttt{val M c (\allowbreak{}gatherS s) \ensuremath{\sigma}}}$.

\paragraph{Boolean elimination: \textnormal{\texttt{bcanon}}.} $\textnormal{\texttt{splitAt v g}} = (A,B)$ with $A$ the monomials
of $g$ containing $v$, $v$ removed, and $B$ the others; \textnormal{\texttt{fz\_\allowbreak{}splitAt}}:
$\textnormal{\texttt{fz \ensuremath{\tau} g}} = \tau(v)\,\textnormal{\texttt{fz \ensuremath{\tau} A}} + \textnormal{\texttt{fz \ensuremath{\tau} B}}$. For the first constraint $q$, let
$g = \textnormal{\texttt{normF (\allowbreak{}[] :: q)}}$ (the indicator $1+q$) and $(A,B) = \textnormal{\texttt{splitAt v g}}$;
$\textnormal{\texttt{detValue v q}} = A\mathbin{+\!\!+}B$. \textnormal{\texttt{bcanonOK v s}} holds when cons $= q :: \_$,
$v\notin\text{path}$, and $\textnormal{\texttt{f2mulN B (\allowbreak{}detValue v q) = []}}$, i.e.\ $B(A+B)=0$: on the support
$g=vA+B=1$, $v$ is determined as $A+B$ (\textnormal{\texttt{det\_\allowbreak{}bits}}: $xa+b=1\wedge b(a+b)=0\Rightarrow x=a+b$).
With $w=\textnormal{\texttt{detValue v q}}$ and cons $= q :: rest$:
\textnormal{\texttt{bcanonS M v s}} keeps the path, has phase \textnormal{\texttt{normP M (\allowbreak{}substP M v w phase)}}, replaces
every constraint $r$ of $q::rest$ (the first one included) by \textnormal{\texttt{substFN v w r}}, and has the
rows \textnormal{\texttt{substFN v w r}} for each old row followed by \textnormal{\texttt{normF (\allowbreak{}[v] :: w)}} ($v+A+B$).
\textnormal{\texttt{bcanon\_\allowbreak{}sound}} ($2\le M$): $\textnormal{\texttt{val M c s \ensuremath{\sigma}}} = \textnormal{\texttt{val M c (\allowbreak{}bcanonS M v s) \ensuremath{\sigma}}}$.

\section{Further qubit steps and step lists: \texorpdfstring{\file{PathSumRules}}{PathSumRules}}\label{sec:lean-pathsum-steps}

\file{PathSumRules} (imports \file{PathSumSem}) adds the steps [HHgen], [Rem], [omega], [HHnl],
[Pin], normalization, [Dyadic], phase-term dropping, and boundary substitution, and composes all
steps into the step language \textnormal{\texttt{Step}}. Each step reads the head path variables. Auxiliary
facts: \textnormal{\texttt{zq0}}, \textnormal{\texttt{zq1}}, \textnormal{\texttt{summand\_\allowbreak{}path}} (the summand ignores the path),
\textnormal{\texttt{summand\_\allowbreak{}pair}} (for $y$ in no constraint or row, the summand at $y=0$ plus at $y=1$ equals
$(1+\operatorname{ex}(\textnormal{\texttt{pn \ensuremath{\tau} (\allowbreak{}yq y phase)}}/M))$ times the summand with phase
\textnormal{\texttt{rest y phase}}), \textnormal{\texttt{yq\_\allowbreak{}map}}, \textnormal{\texttt{one\_\allowbreak{}add\_\allowbreak{}half}} (for \textnormal{\texttt{halfOK}},
$1+\operatorname{ex}(\textnormal{\texttt{pn \ensuremath{\tau} (\allowbreak{}yq y P)}}/M) = 2[\textnormal{\texttt{fz \ensuremath{\tau} (\allowbreak{}quot y P)}} = 0]$), and
\textnormal{\texttt{ex\_\allowbreak{}zqN}}.

\begin{definition}[Substitution throughout]\label{def:lean-pathsum-substD}
$\textnormal{\texttt{substD M z B s}}$ keeps the path and replaces $z$ by $B$ in the phase
(\textnormal{\texttt{normP M (\allowbreak{}substP M z B phase)}}), in every constraint, and in every row
(\textnormal{\texttt{normF (\allowbreak{}substF z B q)}}). \textnormal{\texttt{summand\_\allowbreak{}substD}} ($2\le M$):
$\textnormal{\texttt{summand M c (\allowbreak{}substD M z B s) \ensuremath{\tau}}} = \textnormal{\texttt{summand M c s}}\,(\tau[z\mapsto\textnormal{\texttt{fz \ensuremath{\tau} B}}])$.
\end{definition}

\paragraph{[HHgen].} \textnormal{\texttt{hhGenOK M s}}: path $= y::z::ys$, $2\le M$, $y\ne z$, $y\notin ys$,
$z\notin ys$, \textnormal{\texttt{halfOK M y phase}}, $y\notin\textnormal{\texttt{fvars (\allowbreak{}cons ++ rows)}}$, and, with
$(A,B) = \textnormal{\texttt{splitAt z (\allowbreak{}quot y phase)}}$ (the quotient is $zA+B$),
$\textnormal{\texttt{normF (\allowbreak{}f2mul (\allowbreak{}[] :: A) (\allowbreak{}[] :: B)) = []}}$, i.e.\ $(1+A)(1+B)=0$. The step sets $z := B$ and
removes $z$ from the path, keeping $y$:
$\textnormal{\texttt{hhGenS M s}} = \textnormal{\texttt{substD M z B \{s with path := y :: ys\}}}$. Factor $1$.
\begin{theorem}[\textnormal{\texttt{hhGen\_\allowbreak{}sound}}]
$\textnormal{\texttt{hhGenOK M s}}\Rightarrow\textnormal{\texttt{val M c s \ensuremath{\sigma}}} = \textnormal{\texttt{val M c (\allowbreak{}hhGenS M s) \ensuremath{\sigma}}}$. The core is
\textnormal{\texttt{hhgen\_\allowbreak{}pick}}: for $a,b\in\Zd2$ with $(1+a)(1+b)=0$,
$[0a+b=0]\,2R(0)+[1a+b=0]\,2R(1) = [ba+b=0]\,2R(b)$.
\end{theorem}

\paragraph{[Rem].} \textnormal{\texttt{remKeep y m t}} keeps a term $t$ when it contains $y$ or misses some
variable of $m$. \textnormal{\texttt{remOK M s}}: path $=y::ys$, $0<M$, $y\notin ys$, \textnormal{\texttt{halfOK M y phase}},
$y\notin\textnormal{\texttt{fvars (\allowbreak{}cons ++ rows)}}$, and the quotient has exactly one monomial.
$\textnormal{\texttt{remS s}}$ filters the phase by \textnormal{\texttt{remKeep y m}} with $m$ the quotient monomial
(\textnormal{\texttt{headD []}}). Factor $1$. On the support of the sum over $y$ the quotient monomial is $0$, so
every term without $y$ divisible by it vanishes.
\begin{theorem}[\textnormal{\texttt{rem\_\allowbreak{}sound}}]
$\textnormal{\texttt{remOK M s}}\Rightarrow\textnormal{\texttt{val M c s \ensuremath{\sigma}}} = \textnormal{\texttt{val M c (\allowbreak{}remS s) \ensuremath{\sigma}}}$.
\end{theorem}
\noindent Lemmas: \textnormal{\texttt{yTerms\_\allowbreak{}filter}}, \textnormal{\texttt{rest\_\allowbreak{}filter}}, \textnormal{\texttt{quot\_\allowbreak{}filter}},
\textnormal{\texttt{yq\_\allowbreak{}filter}}, \textnormal{\texttt{halfOK\_\allowbreak{}filter}} (a filter keeping every term of $y$ does not change the
quotient data), \textnormal{\texttt{pn\_\allowbreak{}filter\_\allowbreak{}zero}}, \textnormal{\texttt{mz\_\allowbreak{}zero\_\allowbreak{}mem}}, \textnormal{\texttt{mz\_\allowbreak{}zero\_\allowbreak{}of\_\allowbreak{}mem}}.

\paragraph{[omega].} \textnormal{\texttt{omega\_\allowbreak{}id}} states, for $e\in\{1,3\}$ and $b\in\Zd2$,
\[
1+\omega_8^{2e}\,(-1)^{b} = \sqrt2\,\omega_8^{f(e)}\cdot\begin{cases}1 & b=0\\ \omega_8^{2(4-e)} & b=1,\end{cases}
\qquad f(1)=1,\ f(3)=7,
\]
with the four instances \textnormal{\texttt{omega\_\allowbreak{}e1F0}}, \textnormal{\texttt{omega\_\allowbreak{}e1F1}}, \textnormal{\texttt{omega\_\allowbreak{}e3F0}},
\textnormal{\texttt{omega\_\allowbreak{}e3F1}} and $\omega_8^4=-1$ (\textnormal{\texttt{om8\_\allowbreak{}four}}). \textnormal{\texttt{omegaOK M e s}}: path $=y::ys$,
$0<M$, $4\mid M$, $y\notin ys$, the constant term of \textnormal{\texttt{yq y phase}} is $e\,M/4 \bmod M$
(\textnormal{\texttt{constK (\allowbreak{}yq y phase) \% M = e * (\allowbreak{}M / 4)}}), every nonconstant term of $y$ has $2k=M$, and
$y\notin\textnormal{\texttt{fvars (\allowbreak{}cons ++ rows)}}$. So the phase of $y$ is $y\,(e/4+Q/2)$ with
$Q=\textnormal{\texttt{omQ y phase}}$, the monomials of \textnormal{\texttt{strip (\allowbreak{}yq y phase)}}.
\[
\textnormal{\texttt{omegaS M e s}} = \{ys,\ \textnormal{\texttt{normP M}}\bigl(\textnormal{\texttt{rest y phase}} \mathbin{+\!\!+} (4-e)\tfrac{M}{4}\cdot\textnormal{\texttt{lift M Q}}\bigr),\ \text{cons},\ \text{rows}\}.
\]
\begin{theorem}[\textnormal{\texttt{omega\_\allowbreak{}sound}}]
For $2\le M$, $e\in\{1,3\}$, and \textnormal{\texttt{omegaOK M e s = true}},
$\textnormal{\texttt{val M c s \ensuremath{\sigma}}} = \textnormal{\texttt{val M (\allowbreak{}rt 2 * om 8 \^{} f(\allowbreak{}e) * c) (\allowbreak{}omegaS M e s) \ensuremath{\sigma}}}$ (\textnormal{\texttt{ex\_\allowbreak{}div4}}
gives $\operatorname{ex}(er/M) = \omega_8^{2e}$ for $M=4r$).
\end{theorem}

\paragraph{[HHnl].} \textnormal{\texttt{hhNlOK M s}}: path $= y::y'::ys$, $2\le M$, $y\ne y'$, $y,y'\notin ys$,
both \textnormal{\texttt{halfOK}}, neither in a constraint or row, $y$ not in the quotient of $y'$, and $y'$ not
in the quotient of $y$. \textnormal{\texttt{nlValue y P}} $= [y] :: [[y]\mathbin{+\!\!+}u : u\in\textnormal{\texttt{quot y P}}]$, the
polynomial $y+yQ$. $\textnormal{\texttt{hhNlS M s}}$ sets $y' := y+yQ$ in the phase
(\textnormal{\texttt{normP M (\allowbreak{}substP M y' (\allowbreak{}nlValue y phase) phase)}}) and leaves path $y::ys$.
\begin{theorem}[\textnormal{\texttt{hhNl\_\allowbreak{}sound}}]
$\textnormal{\texttt{hhNlOK M s}}\Rightarrow\textnormal{\texttt{val M c s \ensuremath{\sigma}}} = \textnormal{\texttt{val M (\allowbreak{}2 * c) (\allowbreak{}hhNlS M s) \ensuremath{\sigma}}}$, from
$\sum_{y,y'}(-1)^{yQ+y'Q'} = 2\sum_y(-1)^{y(Q\vee Q')}$ (\textnormal{\texttt{yTerms\_\allowbreak{}rest}}, \textnormal{\texttt{not\_\allowbreak{}mem\_\allowbreak{}flat}}).
\end{theorem}

\paragraph{[Pin].} \textnormal{\texttt{pinAt v B q}} holds when \textnormal{\texttt{pivotable v q}} and $q.\textnormal{\texttt{erase [v]}} = B$.
\textnormal{\texttt{pinOK B s}}: path $= v :: ys$, $v\notin ys$, and some constraint satisfies \textnormal{\texttt{pinAt v B}}.
With $q$ the first such constraint, $\textnormal{\texttt{pinS M B s}}$ drops $v$ from the path, removes $q$, and
substitutes $v:=B$ in the phase (\textnormal{\texttt{substNP M v B}}), the other constraints, and the rows
(\textnormal{\texttt{substNF v B}}). Factor $1$.
\begin{theorem}[\textnormal{\texttt{pin\_\allowbreak{}sound}}]
$2\le M$, $\textnormal{\texttt{pinOK B s}}\Rightarrow\textnormal{\texttt{val M c s \ensuremath{\sigma}}} = \textnormal{\texttt{val M c (\allowbreak{}pinS M B s) \ensuremath{\sigma}}}$.
\end{theorem}

\paragraph{Normalization.} $\textnormal{\texttt{normS M s}}$ applies \textnormal{\texttt{normP M}} to the phase and \textnormal{\texttt{normF}}
to every constraint and row. \textnormal{\texttt{norm\_\allowbreak{}sound}} ($0<M$, no side condition):
$\textnormal{\texttt{val M c s \ensuremath{\sigma}}} = \textnormal{\texttt{val M c (\allowbreak{}normS M s) \ensuremath{\sigma}}}$. \textnormal{\texttt{val\_\allowbreak{}smul}}:
$\textnormal{\texttt{val M (\allowbreak{}a * c) s \ensuremath{\sigma}}} = a\cdot\textnormal{\texttt{val M c s \ensuremath{\sigma}}}$.

\subsection{[Dyadic]}\label{sec:lean-pathsum-dyadic}

[Dyadic] sums $k$ digit variables $d_0,\dots,d_{k-1}$ whose phase is $YS/2^k$ turns with
$Y = \sum_a 2^a d_a$ and $S$ a phase polynomial free of the digits, using
$\sum_{Y<2^k}\omega_{2^k}^{YS} = 2^k[2^k\mid S] = \prod_{j<k}\sum_{f_j}(-1)^{f_jb_j}$ with $b_j$
the binary digits of $S$, over $k$ fresh variables $f_j$.

\begin{definition}[Digit data]\label{def:lean-pathsum-dydata}
\begin{itemize}
\item \textnormal{\texttt{negP N P}} negates every coefficient mod $N$: $(m,k)\mapsto(m, N-(k\bmod N))$;
  \textnormal{\texttt{pn\_\allowbreak{}negP}}: $\textnormal{\texttt{pn \ensuremath{\tau} (\allowbreak{}negP N P)}}\equiv-\textnormal{\texttt{pn \ensuremath{\tau} P}}\pmod N$.
\item \textnormal{\texttt{dyBits k S : Option (\allowbreak{}List F2)}} computes the binary digits of $S$ modulo $2^k$, low digit
  first: $\textnormal{\texttt{dyBits 0 S}} = \textnormal{\texttt{some []}}$; for $k+1$, $b$ is the monomials of the terms of $S$ with
  odd coefficient, $D = \textnormal{\texttt{normP}}_{2^{k+1}}(S\mathbin{+\!\!+}\textnormal{\texttt{negP}}_{2^{k+1}}(\textnormal{\texttt{lift}}_{2^{k+1}}\,b))$;
  when every coefficient of $D$ is even the result is $b :: \textnormal{\texttt{dyBits k}}(D/2)$ (coefficients
  halved), else \textnormal{\texttt{none}}. \textnormal{\texttt{pn\_\allowbreak{}odd}} ($\textnormal{\texttt{pn \ensuremath{\tau} S}} \bmod 2 = \textnormal{\texttt{fz}}$ of the odd
  monomials), \textnormal{\texttt{pn\_\allowbreak{}half\_\allowbreak{}all}}, and \textnormal{\texttt{dy\_\allowbreak{}bits}} give: $\textnormal{\texttt{dyBits k S = some bs}}$ implies
  $2^k\mid\textnormal{\texttt{pn \ensuremath{\tau} S}}\iff\forall b\in bs,\ \textnormal{\texttt{fz \ensuremath{\tau} b}} = 0$; \textnormal{\texttt{dyBits\_\allowbreak{}length}}: $|bs| = k$.
\item $\textnormal{\texttt{dyProd N K k w}} = \prod_{a<k}\bigl(1+\operatorname{ex}(2^a w N/K)\bigr)$ (recursive,
  weight doubling); \textnormal{\texttt{dyProd\_\allowbreak{}double}} ($\textnormal{\texttt{dyProd N (\allowbreak{}2K) k (\allowbreak{}2w)}} = \textnormal{\texttt{dyProd N K k w}}$),
  \textnormal{\texttt{dyProd\_\allowbreak{}geom}}: $\textnormal{\texttt{dyProd N (\allowbreak{}2\^{}k) k 1}} = 2^k[2^k\mid N]$, \textnormal{\texttt{dyProd\_\allowbreak{}succ}}.
\item $\textnormal{\texttt{dyTerms M K S ds w}}$: for the digits $ds = [d_0,\dots]$, the terms
  $(d_a\mathbin{+\!\!+}m,\ 2^a w\,(M/K)\,k)$ for $(m,k)\in S$, i.e.\ $d_a S 2^a w/K$ in units of $1/M$.
\item $\textnormal{\texttt{dyNew M fs bs}}$: for each pair $(f_j,b_j)$, the terms $([f_j]\mathbin{+\!\!+}m, M/2)$ for
  $m\in b_j$, i.e.\ $f_jb_j/2$.
\item \textnormal{\texttt{hasAny ds m}} holds when $m$ contains a variable of $ds$.
\item $\textnormal{\texttt{dnum \ensuremath{\tau} [B\_\allowbreak{}0,\allowbreak{} B\_\allowbreak{}1,\allowbreak{} .\allowbreak{}.\allowbreak{}.\allowbreak{}]}} = \sum_i 2^i\,\textnormal{\texttt{(\allowbreak{}fz \ensuremath{\tau} B\_\allowbreak{}i).\allowbreak{}val}}$, the number the digits
  spell; $\textnormal{\texttt{tLift M w Bs}} = \sum_i w2^i\,\textnormal{\texttt{lift M B\_\allowbreak{}i}}$; \textnormal{\texttt{dnum\_\allowbreak{}lt}},
  \textnormal{\texttt{dnum\_\allowbreak{}update\_\allowbreak{}not}}, \textnormal{\texttt{zq\_\allowbreak{}val}}, \textnormal{\texttt{pn\_\allowbreak{}tLift}} ($\equiv w\cdot\textnormal{\texttt{dnum \ensuremath{\tau} Bs}}$),
  \textnormal{\texttt{dnum\_\allowbreak{}eq\_\allowbreak{}zero}} ($\textnormal{\texttt{dnum}} = 0\iff$ every digit is $0$).
\item \textnormal{\texttt{liftT M d B}} is the lift keeping the products of at most $d$ monomials;
  \textnormal{\texttt{pn\_\allowbreak{}liftT}}: $\textnormal{\texttt{pn \ensuremath{\tau} (\allowbreak{}lift M B)}} = \textnormal{\texttt{pn \ensuremath{\tau} (\allowbreak{}liftT M d B)}} + (M-2)^d z$ for some $z\in\Z$.
  $\textnormal{\texttt{bLift M w d [B\_\allowbreak{}0,\allowbreak{} .\allowbreak{}.\allowbreak{}.\allowbreak{}]}} = \sum_i w2^i\,\textnormal{\texttt{liftT M (\allowbreak{}d-{}i) B\_\allowbreak{}i}}$; \textnormal{\texttt{pn\_\allowbreak{}bLift}}: for
  $1\le k$, $w2^d=2^k$, and $|bs|\le d$, $\textnormal{\texttt{pn \ensuremath{\tau} (\allowbreak{}bLift (\allowbreak{}2\^{}k) w d bs)}}\equiv w\cdot\textnormal{\texttt{dnum \ensuremath{\tau} bs}}\pmod{2^k}$.
\item \textnormal{\texttt{bitsOK k S bs}}: $1\le k$, $|bs|=k$, and
  $\textnormal{\texttt{normP}}_{2^k} S = \textnormal{\texttt{normP}}_{2^k}(\textnormal{\texttt{bLift (\allowbreak{}2\^{}k) 1 k bs}})$. \textnormal{\texttt{bits\_\allowbreak{}iff}} gives the
  same divisibility fact as \textnormal{\texttt{dy\_\allowbreak{}bits}} from \textnormal{\texttt{bitsOK}}.
\end{itemize}
\end{definition}

\begin{definition}[[Dyadic] side conditions]\label{def:lean-pathsum-dyrest}
\textnormal{\texttt{dyRest M ds S fs bs s}} is the conjunction, in this order:
$0<M$; $2^{|ds|}\mid M$; $2\mid M$; the path begins with $ds$; the path has no repeats;
$fs\mathbin{+\!\!+}(\text{path}\text{ after }ds)$ has no repeats; $|fs| = |ds|$; no digit is in a
constraint or row; no monomial of $S$ contains a digit; no fresh variable is in the phase, a
constraint or row, or a digit polynomial ($\textnormal{\texttt{fvars bs}}$); and
$\textnormal{\texttt{normP M}}(\text{the phase terms containing a digit}) = \textnormal{\texttt{normP M (\allowbreak{}dyTerms M (\allowbreak{}2\^{}\{|ds|\}) S ds 1)}}$.
\textnormal{\texttt{dyadicBOK M ds S fs bs s := bitsOK |ds| S bs \&\& dyRest M ds S fs bs s}}.
\end{definition}

\noindent The transformer \textnormal{\texttt{dyadicBS}} (given digits) produces
\[
\{fs\mathbin{+\!\!+}(\text{path after } ds),\ (\text{phase terms free of } ds)\mathbin{+\!\!+}\textnormal{\texttt{dyNew M fs bs}},\ \text{cons},\ \text{rows}\}.
\]
\begin{theorem}[\textnormal{\texttt{dyadic\_\allowbreak{}core}}, \textnormal{\texttt{dyadic\_\allowbreak{}sound}}, \textnormal{\texttt{dyadicB\_\allowbreak{}sound}}]
If $|bs|=|ds|$, $2^{|ds|}\mid\textnormal{\texttt{pn \ensuremath{\tau} S}}\iff\forall b\in bs,\ \textnormal{\texttt{fz \ensuremath{\tau} b}}=0$ for every $\tau$,
and \textnormal{\texttt{dyRest M ds S fs bs s}}, then the value is unchanged. \textnormal{\texttt{dyadic\_\allowbreak{}sound}} and
\textnormal{\texttt{dyadicB\_\allowbreak{}sound}} derive the divisibility fact from \textnormal{\texttt{dy\_\allowbreak{}bits}} and \textnormal{\texttt{bits\_\allowbreak{}iff}}. The
proof evaluates the sum over the digits as \textnormal{\texttt{dyProd}} (\textnormal{\texttt{bsum\_\allowbreak{}dy}}) and the sum over the fresh
variables as $\prod_j 2[\textnormal{\texttt{fz \ensuremath{\tau} b\_\allowbreak{}j}} = 0]$ (\textnormal{\texttt{bsum\_\allowbreak{}new}}, \textnormal{\texttt{prod\_\allowbreak{}ite}}).
\end{theorem}

\paragraph{[Dyadic] by column addition.} A second form of [Dyadic] needs no fresh variables: it
writes the divisibility $2^k\mid S$ as constraints over existing variables $c_0,\dots,c_{k-1}$
that occur in $S$ as linear terms of weights $\pm2^i$, and the digits of the rest of $S$ as $\mathbb
F_2$ polynomials computed by binary column addition.
\begin{itemize}
\item \lean!cntAdd : List F2 → F2 → List F2! adds a bit $x$ to a binary counter of $\mathbb F_2$
  digits, low digit first, the carry past the top dropped:
  $\textnormal{\texttt{cntAdd [] x}}=[\,]$,
  $\textnormal{\texttt{cntAdd (b :: bs) x}} = \textnormal{\texttt{normF (b ++ x)}} :: \textnormal{\texttt{cntAdd bs (f2mulN b x)}}$.
  $\textnormal{\texttt{inum \ensuremath{\tau} [b\_0, b\_1, \ldots]}} = \sum_i 2^i\lfloor\textnormal{\texttt{fz \ensuremath{\tau} b\_i}}\rfloor$ and
  $\textnormal{\texttt{csum \ensuremath{\tau} col}}$ is the number of items of a column with value $1$.
  \lean!inum_cntAdd!: $\textnormal{\texttt{inum (cntAdd bs x)}}\equiv\textnormal{\texttt{inum bs}}+\lfloor\textnormal{\texttt{fz x}}\rfloor\pmod{2^{|bs|}}$;
  \lean!inum_foldl!: folding a column into a counter adds its count mod $2^{|bs|}$; \lean!cntAdd_length!,
  \lean!foldl_cntAdd_length!, \lean!inum_zero!, \lean!inum_bounds! ($0\le\textnormal{\texttt{inum}}<2^{|L|}$).
\item \textnormal{\texttt{carry cs cols}} appends carry $c_i$ to column $i$ (the first carry to the first
  column); $\textnormal{\texttt{cval \ensuremath{\tau} cols}} = \sum_i 2^i\,\textnormal{\texttt{csum}}(\text{column } i)$;
  \lean!cval_carry!: $\textnormal{\texttt{cval (carry cs cols)}} = \textnormal{\texttt{cval cols}}+\textnormal{\texttt{inum cs}}$ when $|cs|\le|cols|$.
\item \textnormal{\texttt{colRun n cols}}: $\textnormal{\texttt{colRun 0 \_}}=[\,]$, $\textnormal{\texttt{colRun (n+1) []}}=[\,]$, and for
  $col::rest$ the column is folded into the zero counter of $n+1$ digits by \lean!cntAdd!; its low
  digit $b$ is the output digit, and its higher digits $bt$ are carried into the remaining columns:
  $b::\textnormal{\texttt{colRun n (carry bt rest)}}$. \lean!colRun_spec!: for $n$ columns,
  $\textnormal{\texttt{cval cols}}\equiv\textnormal{\texttt{inum (colRun n cols)}}\pmod{2^n}$ and the output has length $n$.
\item \textnormal{\texttt{colsOf k T}}: column $j<k$ holds $[m]$ for each term $(m,a)$ of $T$ with bit $j$ of $a$
  set (\lean!pn_halves!, \lean!cval_colsOf!: $\textnormal{\texttt{cval (colsOf k T)}}\equiv\textnormal{\texttt{pn T}}\pmod{2^k}$).
\item \textnormal{\texttt{colRows cs ds}} is $[\textnormal{\texttt{normF ([c\_i] :: d\_i)}}]_i$ (\lean!rows_iff!: all rows vanish iff
  $\textnormal{\texttt{inum}}([[c_i]]_i)=\textnormal{\texttt{inum ds}}$); \textnormal{\texttt{cpoly w cs}} is the terms $([c_i],w2^i)$
  (\lean!pn_cpoly!); \textnormal{\texttt{colP cs S}} and \textnormal{\texttt{colR cs S}} are the terms of $S$ on, and off,
  the monomials $[c_i]$.
\item \textnormal{\texttt{colCheck k S cs pos}}: $|cs|=k$ and
  $\textnormal{\texttt{normP}}_{2^k}(\pm\textnormal{\texttt{colP cs S}}) = \textnormal{\texttt{normP}}_{2^k}(\textnormal{\texttt{cpoly 1 cs}})$, the sign $+$
  for $pos$ (else \lean!negP!).
  \textnormal{\texttt{colBits k S cs pos = colRows cs (colRun k (colsOf k R'))}} with $R'=-\textnormal{\texttt{colR cs S}}$ mod $2^k$
  for $pos$ and $R'=\textnormal{\texttt{colR cs S}}$ else.
  \lean!col_bits! (from \lean!dvd_iff_eq!, \lean!pair_iff!): \lean!colCheck! gives
  $2^k\mid\textnormal{\texttt{pn \ensuremath{\tau} S}}\iff\forall b\in\textnormal{\texttt{colBits k S cs pos}},\ \textnormal{\texttt{fz \ensuremath{\tau} b}}=0$.
\item \textnormal{\texttt{colRest M ds S s}}: $0<M$; $2^{|ds|}\mid M$; the path begins with $ds$; the path has no
  repeats; no digit is in a constraint or row; no monomial of $S$ contains a digit; and
  $\textnormal{\texttt{normP M}}(\text{phase terms with a digit}) = \textnormal{\texttt{normP M (dyTerms M (2\^{}|ds|) S ds 1)}}$.
  \lean!dycol_core!: given the divisibility fact for constraints $bs$ and \lean!colRest!, the value equals
  $2^{|ds|}c$ times the value of the path after $ds$, the phase terms free of $ds$, the constraints
  $cons\mathbin{+\!\!+}bs$ and the same rows.
\item \textnormal{\texttt{pickU K w bad S}} is the least $u\notin bad$ with a term $([u],a)$ of $S$, $a\bmod K=w$;
  \textnormal{\texttt{pickCs K pos bad S n acc}} picks $n$ more variables after $acc$, the next of weight
  $2^{|acc|}\bmod K$ ($pos$) or $(K-2^{|acc|}\bmod K)\bmod K$, each excluded afterwards.
  \textnormal{\texttt{colBad ds S s}} lists the path, the variables of the constraints and rows, of the phase
  terms off the digits, and of the monomials of $S$ of degree at least $2$.
  \textnormal{\texttt{colData M ds s}}: for $ds=d::\_$ and $K=2^{|ds|}$,
  $S=\textnormal{\texttt{normP K}}[(m\setminus d,\ a/(M/K)) : (m,a)\in\textnormal{\texttt{yTerms d phase}}]$; the choice
  \textnormal{\texttt{pickCs K true}} is tried first, then \textnormal{\texttt{pickCs K false}}; \lean!none! for $ds=[\,]$ or no
  choice.
\item \textnormal{\texttt{columnsOK M ds s}}: \textnormal{\texttt{colData}} gives $(S,cs,pos)$, \textnormal{\texttt{colRest M ds S s}} and
  \textnormal{\texttt{colCheck |ds| S cs pos}}. \textnormal{\texttt{columnsS M ds s}} drops $ds$ from the path and their phase
  terms and appends \textnormal{\texttt{colBits |ds| S cs pos}} to the constraints.
\end{itemize}
\begin{theorem}[\textnormal{\texttt{columns\_sound}}]
$\textnormal{\texttt{columnsOK M ds s = true}}\Rightarrow\textnormal{\texttt{val M c s \ensuremath{\sigma}}} = \textnormal{\texttt{val M (2\^{}|ds| * c) (columnsS M ds s) \ensuremath{\sigma}}}$.
The proof is \lean!dycol_core! with \lean!col_bits!.
\end{theorem}

\paragraph{Canonical pivot with the phase kept.} \textnormal{\texttt{canonHS v s}} is \textnormal{\texttt{canonS}} without the
phase substitution: the first constraint $q$ pivotable at $v$ becomes the row \textnormal{\texttt{normF q}}, and
$v:=q.\textnormal{\texttt{erase [v]}}$ is substituted (\textnormal{\texttt{normF (substF v B r)}}) in the other constraints and the
rows. \lean!canonH_sound! (side condition \textnormal{\texttt{canonOK v s}}, no bound on $M$): value unchanged.

\paragraph{Phase terms off the constraints.} \textnormal{\texttt{dropT C m}} holds when for some constraint $q$
the mask xor of $m\cdot(1+q)$, over the variables of $q$, is $0$ ($m(1+q)=0$ identically, so $m$
vanishes where $q$ holds; \textnormal{\texttt{dropT\_\allowbreak{}zero}}, \textnormal{\texttt{fz\_\allowbreak{}xorF\_\allowbreak{}zero}}, \textnormal{\texttt{fz\_\allowbreak{}mapor}}).
\textnormal{\texttt{pdropS s}} drops every phase term whose monomial satisfies \textnormal{\texttt{dropT s.\allowbreak{}cons}}.
\textnormal{\texttt{pdrop\_\allowbreak{}sound}}: no side condition, value unchanged.

\paragraph{Boundary substitution.} $\textnormal{\texttt{bsubS M v s}}$: with $q$ the first constraint pivotable at
$v$, $v$ is replaced by $q.\textnormal{\texttt{erase [v]}}$ in the phase only (\textnormal{\texttt{normP}}); constraints and rows
stay. \textnormal{\texttt{bsub\_\allowbreak{}sound}} ($2\le M$, side condition \textnormal{\texttt{canonOK v s}}): value unchanged.

\subsection{The step language \texorpdfstring{\textnormal{\texttt{Step}}}{Step}}

\begin{definition}[\textnormal{\texttt{Step}}]\label{def:lean-pathsum-Step}
\textnormal{\texttt{inductive Step}} has the constructors listed in \cref{tab:lean-pathsum-steps}.
\textnormal{\texttt{stepOK M : Step \ensuremath{\to} Data \ensuremath{\to} Bool}} is the side condition, \textnormal{\texttt{stepS M}} the transformer,
and \textnormal{\texttt{fac : Step \ensuremath{\to} \ensuremath{\mathbb{C}}}} the constant factor.
\begin{align*}
\textnormal{\texttt{runOK M [] s}} &= \textnormal{\texttt{true}},\\ \textnormal{\texttt{runOK M (\allowbreak{}st :: sts) s}} &= \textnormal{\texttt{stepOK M st s \&\& runOK M sts (\allowbreak{}stepS M st s)}},\\
\end{align*}
and \textnormal{\texttt{run M [] s = s}}, \textnormal{\texttt{run M (\allowbreak{}st :: sts) s = run M sts (\allowbreak{}stepS M st s)}},
\textnormal{\texttt{facs [] = 1}}, \textnormal{\texttt{facs (\allowbreak{}st :: sts) = fac st * facs sts}}.
\end{definition}

\begin{xltabular}{\textwidth}{@{}>{\raggedright\arraybackslash}p{3.6cm}>{\raggedright\arraybackslash}X>{\raggedright\arraybackslash}p{2.3cm}>{\raggedright\arraybackslash}p{2.6cm}@{}}
\caption{The constructors of \texttt{PSem.Step}: side condition, transformer, and factor.}\label{tab:lean-pathsum-steps}\\
\toprule Constructor & Side condition / transformer & Factor & Rule\\\midrule\endfirsthead
\toprule Constructor & Side condition / transformer & Factor & Rule\\\midrule\endhead
\textnormal{\texttt{perm l}} & \textnormal{\texttt{decide (\allowbreak{}s.\allowbreak{}path.\allowbreak{}Perm l)}} / \textnormal{\texttt{permS l}} & $1$ & reorder\\
\textnormal{\texttt{elim}} & \textnormal{\texttt{elimOK}} / \textnormal{\texttt{elimS}} & $2$ & [Elim]\\
\textnormal{\texttt{hhCons}} & \textnormal{\texttt{hhConsOK M}} / \textnormal{\texttt{hhConsS}} & $2$ & [HH], constraint\\
\textnormal{\texttt{hhPivot}} & \textnormal{\texttt{hhPivotOK M}} / \textnormal{\texttt{hhPivotS M}} & $2$ & [HH], pivot\\
\textnormal{\texttt{canon v}} & \textnormal{\texttt{canonOK v}} / \textnormal{\texttt{canonS M v}} & $1$ & canonical pivot\\
\textnormal{\texttt{gather}} & \textnormal{\texttt{gatherOK}} / \textnormal{\texttt{gatherS}} & $1$ & Boolean elimination\\
\textnormal{\texttt{bcanon v}} & \textnormal{\texttt{bcanonOK v}} / \textnormal{\texttt{bcanonS M v}} & $1$ & Boolean pivot\\
\textnormal{\texttt{omega}} & \textnormal{\texttt{omegaOK M 1}} / \textnormal{\texttt{omegaS M 1}} & $\sqrt2\,\omega_8$ & [omega]\\
\textnormal{\texttt{omegaConj}} & \textnormal{\texttt{omegaOK M 3}} / \textnormal{\texttt{omegaS M 3}} & $\sqrt2\,\omega_8^7$ & [omega], conjugate\\
\textnormal{\texttt{hhGen}} & \textnormal{\texttt{hhGenOK M}} / \textnormal{\texttt{hhGenS M}} & $1$ & [HHgen]\\
\textnormal{\texttt{rem}} & \textnormal{\texttt{remOK M}} / \textnormal{\texttt{remS}} & $1$ & [Rem]\\
\textnormal{\texttt{hhNl}} & \textnormal{\texttt{hhNlOK M}} / \textnormal{\texttt{hhNlS M}} & $2$ & [HHnl]\\
\textnormal{\texttt{pin B}} & \textnormal{\texttt{pinOK B}} / \textnormal{\texttt{pinS M B}} & $1$ & [Pin]\\
\textnormal{\texttt{norm}} & \textnormal{\texttt{true}} / \textnormal{\texttt{normS M}} & $1$ & normalization\\
\textnormal{\texttt{pdrop}} & \textnormal{\texttt{true}} / \textnormal{\texttt{pdropS}} & $1$ & phase terms off the support\\
\textnormal{\texttt{dyadicB ds S fs bs}} & \textnormal{\texttt{dyadicBOK M ds S fs bs}} / \textnormal{\texttt{dyadicBS M ds fs bs}} & $1$ & [Dyadic], given digits\\
\textnormal{\texttt{bsub v}} & \textnormal{\texttt{canonOK v}} / \textnormal{\texttt{bsubS M v}} & $1$ & boundary substitution\\
\textnormal{\texttt{front x}} & \textnormal{\texttt{decide (x \ensuremath{\in} s.path)}} / \textnormal{\texttt{frontS x}} & $1$ & reorder\\
\textnormal{\texttt{canonH v}} & \textnormal{\texttt{canonOK v}} / \textnormal{\texttt{canonHS v}} & $1$ & canonical pivot, phase kept\\
\textnormal{\texttt{columns ds}} & \textnormal{\texttt{columnsOK M ds}} / \textnormal{\texttt{columnsS M ds}} & $2^{|ds|}$ & [Dyadic], columns\\
\bottomrule
\end{xltabular}

\begin{theorem}[Step soundness]\label{thm:lean-pathsum-run}
For $2\le M$:
\begin{itemize}
\item \textnormal{\texttt{step\_\allowbreak{}sound}}: $\textnormal{\texttt{stepOK M st s = true}}\Rightarrow\textnormal{\texttt{val M c s \ensuremath{\sigma}}} = \textnormal{\texttt{val M (\allowbreak{}fac st * c) (\allowbreak{}stepS M st s) \ensuremath{\sigma}}}$;
\item \textnormal{\texttt{run\_\allowbreak{}sound}}: $\textnormal{\texttt{runOK M sts s = true}}\Rightarrow\textnormal{\texttt{val M c s \ensuremath{\sigma}}} = \textnormal{\texttt{val M (\allowbreak{}facs sts * c) (\allowbreak{}run M sts s) \ensuremath{\sigma}}}$.
\end{itemize}
\end{theorem}
\noindent \textnormal{\texttt{step\_\allowbreak{}sound}} dispatches on the constructor to the theorems above (\textnormal{\texttt{omega}}
with $e=1$, \textnormal{\texttt{omegaConj}} with $e=3$); \textnormal{\texttt{run\_\allowbreak{}sound}} is induction on the list with
\textnormal{\texttt{val\_\allowbreak{}smul}} and associativity of the factor product.

\begin{example}[A certificate's main theorem]
For the two sides' literal sums \textnormal{\texttt{sL}}, \textnormal{\texttt{sR}} and step lists \textnormal{\texttt{stepsL}},
\textnormal{\texttt{stepsR}}, a certificate rewrites both values with \textnormal{\texttt{run\_\allowbreak{}sound}}, replaces
\textnormal{\texttt{run M stepsL sL}} by its literal result (\cref{sec:lean-pathsum-chain}), and closes the goal
with \textnormal{\texttt{final\_\allowbreak{}sound}}, whose \textnormal{\texttt{finalOK}} and two \textnormal{\texttt{constK}} hypotheses are \textnormal{\texttt{by decide}}
and whose constant identity is the separately proved \textnormal{\texttt{consts}} theorem; see
\cref{sec:lean-pathsum-python}.
\end{example}

\section{Phase differences: \texorpdfstring{\file{PathSumParity}}{PathSumParity}}\label{sec:lean-pathsum-parity}

\file{PathSumParity} (imports \file{PathSumRules}; namespace \textnormal{\texttt{ArchytasZX.PSem}}) compares two
reduced sums whose rows were made by \textnormal{\texttt{canonH}} steps, so that their phases still read the
pivots. The comparison substitutes each row's pivot into the difference of the two phases.

\begin{definition}[\texttt{substRows}, \texttt{rowsOK}, \texttt{diffP}, \texttt{finalPOK}]\label{def:lean-pathsum-diffp}
\begin{itemize}
\item \textnormal{\texttt{substRows M (v :: vs) (r :: rs) D = substRows M vs rs (normP M (substP M v (r.erase [v]) D))}},
  and \textnormal{\texttt{substRows M \_ \_ D = D}} when either list is empty: each pivot replaced by the rest of
  its row, in order.
\item \textnormal{\texttt{rowsOK vs rs}}: the two lists have one length and each row is \textnormal{\texttt{pivotable}} at its
  pivot.
\item \textnormal{\texttt{diffP M vs sL sR = substRows M vs sL.rows (normP M (sL.phase ++ negP M sR.phase))}}.
\item \textnormal{\texttt{finalPOK M vs sL sR}} is the conjunction, in this order: both paths empty; every
  constraint of both sides \textnormal{\texttt{[]}}; equal row lists; \textnormal{\texttt{rowsOK vs sL.rows}}; and
  \textnormal{\texttt{strip (diffP M vs sL sR) = []}}.
\end{itemize}
\end{definition}

\begin{theorem}[\textnormal{\texttt{pn\_substRows}}, \textnormal{\texttt{finalP\_sound}}]
For $2\le M$: if \textnormal{\texttt{rowsOK vs rs}} and every row of $rs$ vanishes at $\sigma$, then
$\textnormal{\texttt{pn \ensuremath{\sigma} (substRows M vs rs D)}}\equiv\textnormal{\texttt{pn \ensuremath{\sigma} D}}\pmod M$. If
\textnormal{\texttt{finalPOK M vs sL sR = true}}, $\textnormal{\texttt{constK (diffP M vs sL sR)}}=k$, and $c_L\,\omega_M^k = c_R$,
then $\textnormal{\texttt{val M cL sL \ensuremath{\sigma}}} = \textnormal{\texttt{val M cR sR \ensuremath{\sigma}}}$ for every $\sigma$.
\end{theorem}
\noindent On the support of the rows each pivot equals the rest of its row
(\textnormal{\texttt{z2\_eq\_of\_add}}), so the substituted difference has the value of the plain difference
there; off the support both summands are $0$. \file{PathSumChain} imports this module.

\section{Amplitude atoms: \texorpdfstring{\file{PathSumAmp}}{PathSumAmp}}\label{sec:lean-pathsum-amp}

\file{PathSumAmp} (imports \file{PathSumRules}) extends path sums by \term{amplitude atoms}: factors
$A_i^{[f]}$ that equal a coefficient $A_i\in\C$ where an $\mathbb F_2$ polynomial $f$ is $1$ and
$1$ where $f$ is $0$. The coefficients form one function $A:\N\to\C$ fixed per certificate. Path
variables that atoms read are \term{held}: summed outside the inner sum.

\begin{definition}[Path sums with atoms]\label{def:lean-pathsum-adata}
\textnormal{\texttt{Atoms := List (\allowbreak{}\ensuremath{\mathbb{N}} \ensuremath{\times} F2)}}. \textnormal{\texttt{structure AData}} has fields \textnormal{\texttt{s : Data}} (the inner sum),
\textnormal{\texttt{held : List \ensuremath{\mathbb{N}}}}, and \textnormal{\texttt{atoms : Atoms}}.
\begin{align*}
\textnormal{\texttt{aval A as \ensuremath{\tau}}} &= \textstyle\prod_{(i,f)\in as}\begin{cases}1 & \textnormal{\texttt{fz \ensuremath{\tau} f}} = 0\\ A_i & \text{else,}\end{cases}\\
\textnormal{\texttt{valA M c A t \ensuremath{\sigma}}} &= \textnormal{\texttt{bsum t.\allowbreak{}held}}\,\bigl(\tau\mapsto\textnormal{\texttt{aval A t.\allowbreak{}atoms \ensuremath{\tau}}}\cdot\textnormal{\texttt{val M c t.\allowbreak{}s \ensuremath{\tau}}}\bigr)\,\sigma .
\end{align*}
\textnormal{\texttt{avars as}} lists the variables the atoms read. Lemmas: \textnormal{\texttt{aval\_\allowbreak{}nil}}, \textnormal{\texttt{aval\_\allowbreak{}cons}},
\textnormal{\texttt{aval\_\allowbreak{}append}}, \textnormal{\texttt{aval\_\allowbreak{}update\_\allowbreak{}not}}, \textnormal{\texttt{aval\_\allowbreak{}perm}}, \textnormal{\texttt{val\_\allowbreak{}update\_\allowbreak{}not}} (\textnormal{\texttt{val}}
is invariant at a variable off the path, phase, constraints, and rows), \textnormal{\texttt{valA\_\allowbreak{}nil}}
($\textnormal{\texttt{valA M c A \ensuremath{\langle}s,\allowbreak{} [],\allowbreak{} []\ensuremath{\rangle} \ensuremath{\sigma}}} = \textnormal{\texttt{val M c s \ensuremath{\sigma}}}$), and \textnormal{\texttt{valA\_\allowbreak{}flat}}: when no path
variable is read by an atom,
$\textnormal{\texttt{valA M c A t \ensuremath{\sigma}}} = \textnormal{\texttt{bsum (\allowbreak{}t.\allowbreak{}held ++ t.\allowbreak{}s.\allowbreak{}path)}}\,(\tau\mapsto\textnormal{\texttt{aval A t.\allowbreak{}atoms \ensuremath{\tau}}}\cdot\textnormal{\texttt{summand M c t.\allowbreak{}s \ensuremath{\tau}}})\,\sigma$.
\end{definition}

\subsection{Steps on path sums with atoms}

Each item gives the side condition, the transformer, the factor, and the soundness theorem.
Hypotheses on the coefficients ($A_j = 0$, $2A_j = \alpha+\beta$, \dots) are explicit theorem
arguments; the step language below collects them as data.

\paragraph{Inner steps.} \textnormal{\texttt{astep\_\allowbreak{}sound}} ($2\le M$, \textnormal{\texttt{stepOK M st t.\allowbreak{}s}}) and
\textnormal{\texttt{arun\_\allowbreak{}sound}} (\textnormal{\texttt{runOK M sts t.\allowbreak{}s}}) carry every \textnormal{\texttt{PSem}} step and step list to the inner
sum, with the same factors.

\paragraph{Holding and releasing.} \textnormal{\texttt{releaseOK t}}: held $= y :: hs$, $y\notin hs$,
$y\notin\text{path}$, $y\notin\textnormal{\texttt{avars atoms}}$; \textnormal{\texttt{releaseS}} moves $y$ to the front of the
inner path (\textnormal{\texttt{release\_\allowbreak{}sound}}). \textnormal{\texttt{holdOK t}}: path $= p :: ps$, $p\notin ps$, $p\notin$ held,
$p\notin\textnormal{\texttt{avars atoms}}$; \textnormal{\texttt{holdS}} moves $p$ to the front of held (\textnormal{\texttt{hold\_\allowbreak{}sound}}).
\textnormal{\texttt{hpermS l}} reorders held under \textnormal{\texttt{decide (\allowbreak{}t.\allowbreak{}held.\allowbreak{}Perm l)}} (\textnormal{\texttt{hperm\_\allowbreak{}sound}}).

\paragraph{Constant atoms.} \textnormal{\texttt{adropS}} removes atoms on the zero polynomial
(\textnormal{\texttt{adrop\_\allowbreak{}sound}}, no condition). \textnormal{\texttt{aconstOK t}}: the head atom's polynomial is
\textnormal{\texttt{[[]]}} ($=1$); \textnormal{\texttt{aconstS}} drops the head atom; \textnormal{\texttt{aconst\_\allowbreak{}sound}}:
$\textnormal{\texttt{valA M c A t \ensuremath{\sigma}}} = \textnormal{\texttt{valA M (\allowbreak{}A (\allowbreak{}aconstIdx t) * c) A (\allowbreak{}aconstS t) \ensuremath{\sigma}}}$, where
\textnormal{\texttt{aconstIdx t}} is the head atom's index (\textnormal{\texttt{(\allowbreak{}t.\allowbreak{}atoms.\allowbreak{}headD (\allowbreak{}0,\allowbreak{} [])).\allowbreak{}1}}).
\textnormal{\texttt{aone\_\allowbreak{}sound}}: if $A_j = 1$ and the head atom has index $j$ (\textnormal{\texttt{aheadOK j t}}), dropping it
leaves the value.

\paragraph{Zero atoms and constraints.} \textnormal{\texttt{azeroOK j t}}: the head atom has index $j$ and its
polynomial is off the path (\textnormal{\texttt{offPath t.\allowbreak{}s f}}: no variable of $f$ is an inner path variable).
\textnormal{\texttt{azeroS}} appends that polynomial to the constraints and drops the atom; \textnormal{\texttt{azero\_\allowbreak{}sound}}
($A_j=0$): value unchanged (\textnormal{\texttt{summand\_\allowbreak{}consApp}}). Conversely \textnormal{\texttt{aconsOK t}}: the last
constraint is off the path; \textnormal{\texttt{aconsS j}} moves it to a new head atom $(j,q)$;
\textnormal{\texttt{acons\_\allowbreak{}sound}} ($A_j=0$).

\paragraph{Merging, picking, normalizing.} \textnormal{\texttt{amergeOK j1 j2 t}}: the head two atoms have indices
$j_1$, $j_2$ and equal polynomials; \textnormal{\texttt{amergeS j}} replaces them by $(j,f)$;
\textnormal{\texttt{amerge\_\allowbreak{}sound}} ($A_j = A_{j_1}A_{j_2}$). \textnormal{\texttt{apickS i}} moves atom $i$ to the front
(\textnormal{\texttt{apick\_\allowbreak{}sound}}, $i<|\text{atoms}|$). \textnormal{\texttt{anormS}} normalizes every atom polynomial
(\textnormal{\texttt{anorm\_\allowbreak{}sound}}).

\paragraph{Substitution in atoms.} \textnormal{\texttt{asubQ v s}} is the first constraint or row (in
\textnormal{\texttt{cons ++ rows}} order) that is pivotable at $v$ and off the path; \textnormal{\texttt{asubOK v t}} holds when it
exists; \textnormal{\texttt{asubS v t}} replaces $v$ by $B=q.\textnormal{\texttt{erase [v]}}$ in every atom, each atom polynomial
by \textnormal{\texttt{substFV vs v B}} over the one variable list $vs=\textnormal{\texttt{vlist (B ++ atoms.flatMap (\ensuremath{\cdot}.2))}}$
(\lean!vlist_atoms!: every variable lies in $vs$; \lean!aval_mapOn!: mapping every atom polynomial by
a value-preserving $F$ keeps \lean!aval!).
\textnormal{\texttt{asub\_\allowbreak{}sound}}: value unchanged; the proof factors the constraint out of the inner sum as an
indicator (\textnormal{\texttt{val\_\allowbreak{}ind}}). \textnormal{\texttt{asubBOK v B t}}: some constraint or row off the path satisfies
\textnormal{\texttt{pinAt v B}}; \textnormal{\texttt{asubBS v B}} substitutes $v:=B$ in every atom by \textnormal{\texttt{substFV}} over
\textnormal{\texttt{vlist (B ++ atoms.flatMap (\ensuremath{\cdot}.2))}} (\textnormal{\texttt{asubB\_\allowbreak{}sound}}).

\paragraph{[Average].} For a held variable $y$ and $g\in\mathbb F_2[x]$:
\textnormal{\texttt{oneY y g}} $=(1+y)g$, \textnormal{\texttt{onlyY y g}} $= yg$, \textnormal{\texttt{reads y a}} holds when atom $a$ reads
$y$, and \textnormal{\texttt{isAlpha y g a}} holds when \textnormal{\texttt{normF a.\allowbreak{}2 = normF (\allowbreak{}oneY y g)}}.
\textnormal{\texttt{averageOK y g t}}: $y$ is the head of held and not in its tail, $y$ is not on the path, in the
phase, a constraint, a row, or $g$, and every atom reading $y$ normalizes to $(1+y)g$ or $yg$.
$\alpha = \textnormal{\texttt{alphaY A y g atoms}}$ is the product of the coefficients of the atoms on $(1+y)g$,
$\beta = \textnormal{\texttt{betaY A y g atoms}}$ the product over the other atoms reading $y$.
$\textnormal{\texttt{averageS y g j t}} = \langle s,\ \text{held tail},\ (\text{atoms not reading }y)\mathbin{+\!\!+}[(j,g)]\rangle$.
\begin{theorem}[\textnormal{\texttt{average\_\allowbreak{}sound}}]
If \textnormal{\texttt{averageOK y g t}} and $2A_j = \alpha+\beta$, then
$\textnormal{\texttt{valA M c A t \ensuremath{\sigma}}} = \textnormal{\texttt{valA M (\allowbreak{}2 * c) A (\allowbreak{}averageS y g j t) \ensuremath{\sigma}}}$. At $y=0$ the atoms reading
$y$ give $1$ where $g=0$ and $\alpha$ where $g=1$; at $y=1$, $1$ and $\beta$ (\textnormal{\texttt{aval\_\allowbreak{}reads}}).
\end{theorem}

\paragraph{[Average] through a half phase.} \textnormal{\texttt{val\_\allowbreak{}half\_\allowbreak{}held}}: a held $y$ with half
coefficients, off the path, constraints, and rows, whose quotient $Q$ is off the path, gives
$\textnormal{\texttt{val M c s \ensuremath{\sigma}[y\ensuremath{\mapsto}b]}} = (-1)^{bQ}\,\textnormal{\texttt{val M c \{s with phase := rest y phase\} \ensuremath{\sigma}}}$.
\textnormal{\texttt{averageHOK M y t}}: $0<M$, $y$ the head of held, not in its tail, not on the path, in no
constraint or row, \textnormal{\texttt{halfOK M y phase}}, the quotient off the path, and every atom reading $y$
normalizes to $1+y$ or $y$ (\textnormal{\texttt{isAlpha y [[]]}}).
$\textnormal{\texttt{averageHS y j0 j1 t}}$ removes the terms of $y$ from the phase, drops $y$ from held, and
replaces the atoms reading $y$ by $(j_0, \textnormal{\texttt{normF (\allowbreak{}1 + Q)}})$ and $(j_1,\textnormal{\texttt{normF Q}})$.
\begin{theorem}[\textnormal{\texttt{averageH\_\allowbreak{}sound}}]
If \textnormal{\texttt{averageHOK M y t}}, $2A_{j_0} = \alpha+\beta$, and $2A_{j_1}=\alpha-\beta$ (with $g = 1$),
the value equals \textnormal{\texttt{valA M (\allowbreak{}2 * c) A (\allowbreak{}averageHS y j0 j1 t) \ensuremath{\sigma}}}. The sum over $y$ is
$\alpha+\beta(-1)^Q$.
\end{theorem}

\paragraph{[Quarter].} $\textnormal{\texttt{qLift M B0 B1}} = \frac M4\textnormal{\texttt{lift M B0}} \mathbin{+\!\!+}\frac M2\textnormal{\texttt{lift M B1}}$
(\textnormal{\texttt{pn\_\allowbreak{}qLift}}: $\equiv r(\lfloor B_0\rfloor+2\lfloor B_1\rfloor)$ for $M=4r$).
\textnormal{\texttt{quarterOK M B0 B1 t}}: inner path $=y::ys$, $0<M$, $4\mid M$, $y\notin ys$, $y$ in no constraint
or row, $B_0$ and $B_1$ off the path, and $\textnormal{\texttt{normP M (\allowbreak{}yq y phase)}} = \textnormal{\texttt{normP M (\allowbreak{}qLift M B0 B1)}}$:
the phase of $y$ is $y(b_0+2b_1)/4$.
\textnormal{\texttt{quarterD M B0 B1 s}} has path $ys$, phase
$\textnormal{\texttt{normP M}}\bigl(\textnormal{\texttt{rest y phase}}\mathbin{+\!\!+}\tfrac{3M}{4}\,\textnormal{\texttt{lift M (\allowbreak{}f2mul B0 B1)}}\bigr)$,
constraints \textnormal{\texttt{cons}} followed by the normalized $(1+B_0)B_1$, and the same rows; $\textnormal{\texttt{quarterS M B0 B1 j t}} = \langle\textnormal{\texttt{quarterD M B0 B1 t.\allowbreak{}s}},\ \text{held},\ \text{atoms}\mathbin{+\!\!+}[(j,B_0)]\rangle$.
\begin{theorem}[\textnormal{\texttt{quarter\_\allowbreak{}val}}, \textnormal{\texttt{quarter\_\allowbreak{}sound}}]
With $2a = 1+\operatorname{ex}(1/4)$,
\textnormal{\texttt{quarter\_\allowbreak{}val}} gives $\textnormal{\texttt{val M c s \ensuremath{\tau}}} = 2\,a^{[B_0]}\,\textnormal{\texttt{val M c (\allowbreak{}quarterD M B0 B1 s) \ensuremath{\tau}}}$;
\textnormal{\texttt{quarter\_\allowbreak{}sound}}: if \textnormal{\texttt{quarterOK M B0 B1 t}} and $2A_j = 1+\omega_M^{M/4}$, the value equals
\textnormal{\texttt{valA M (\allowbreak{}2 * c) A (\allowbreak{}quarterS M B0 B1 j t) \ensuremath{\sigma}}}. The pointwise identity is \textnormal{\texttt{quarter\_\allowbreak{}pt}}:
$1+i^{b_0+2b_1} = 2a^{[b_0]}\,[(1+b_0)b_1=0]\,\operatorname{ex}(3b_0b_1/4)$.
\end{theorem}

\paragraph{Complementary atoms.} \textnormal{\texttt{acompOK j t}}: the head two atoms have index $j$ and
$\textnormal{\texttt{normF (\allowbreak{}f ++ g) = [[]]}}$ ($f+g=1$); \textnormal{\texttt{acompS}} drops both; \textnormal{\texttt{acomp\_\allowbreak{}sound}}: factor $A_j$.

\paragraph{[Turn].} $\textnormal{\texttt{tBit b B}} = B$ if $b$ else $1+B$;
$\textnormal{\texttt{tInd [] s}} = [[]]$, $\textnormal{\texttt{tInd (\allowbreak{}B :: Bs) s}} = \textnormal{\texttt{f2mul (\allowbreak{}tBit (\allowbreak{}s \% 2 == 1) B) (\allowbreak{}tInd Bs (\allowbreak{}s / 2))}}$,
the indicator that the digits spell $s$ (\textnormal{\texttt{fz\_\allowbreak{}tInd}}: $=[\textnormal{\texttt{dnum \ensuremath{\tau} Bs}} = s]$ for
$s<2^{|Bs|}$). $\textnormal{\texttt{tList k}} = [s<2^k : s\ne0,\ s\ne2^{k-1}]$ (\textnormal{\texttt{tList\_\allowbreak{}nodup}},
\textnormal{\texttt{mem\_\allowbreak{}tList}}). $\textnormal{\texttt{tAtoms Bs ats}} = [(j,\textnormal{\texttt{normF (\allowbreak{}tInd Bs s)}}) : (s,j)\in ats]$.
\textnormal{\texttt{turnOK M Bs ats t}}: inner path $=y::ys$, $1\le|Bs|$, $2\le M$, $2^{|Bs|}\mid M$, $y\notin ys$,
$y$ in no constraint or row, every digit off the path, the first entries of $ats$ equal
\textnormal{\texttt{tList |Bs|}}, and $\textnormal{\texttt{normP M (\allowbreak{}yq y phase)}} = \textnormal{\texttt{normP M (\allowbreak{}tLift M (\allowbreak{}M / 2\^{}|Bs|) Bs)}}$:
the phase of $y$ is $yN/2^k$ with $N = \textnormal{\texttt{dnum \ensuremath{\tau} Bs}}$.
$\textnormal{\texttt{turnD M Bs s}}$ removes $y$ and its terms and appends the constraint
$\textnormal{\texttt{normF (\allowbreak{}tInd Bs 2\^{}(\allowbreak{}k-{}1))}}$ ($N\ne 2^{k-1}$); $\textnormal{\texttt{turnS M Bs ats t}}$ applies \textnormal{\texttt{turnD}} to
the inner sum and appends \textnormal{\texttt{tAtoms Bs ats}}.
\begin{theorem}[\textnormal{\texttt{turn\_\allowbreak{}pt}}, \textnormal{\texttt{turn\_\allowbreak{}val}}, \textnormal{\texttt{turn\_\allowbreak{}sound}}]
\textnormal{\texttt{turn\_\allowbreak{}pt}}: for $1\le k$, $N<2^k$, and $2A_j = 1+\operatorname{ex}(s/2^k)$ for every $(s,j)\in ats$,
$1+\operatorname{ex}(N/2^k) = 2\prod_{(s,j)}A_j^{[N=s]}\cdot[N\ne2^{k-1}]$.
\textnormal{\texttt{turn\_\allowbreak{}sound}}: if \textnormal{\texttt{turnOK M Bs ats t}} and $2A_j = 1+\omega_M^{s(M/2^k)}$ for every
$(s,j)\in ats$, the value equals \textnormal{\texttt{valA M (\allowbreak{}2 * c) A (\allowbreak{}turnS M Bs ats t) \ensuremath{\sigma}}}.
\end{theorem}
\noindent Lemmas: \textnormal{\texttt{fz\_\allowbreak{}tBit}}, \textnormal{\texttt{prod\_\allowbreak{}miss}}, \textnormal{\texttt{prod\_\allowbreak{}hit}}, \textnormal{\texttt{aval\_\allowbreak{}tAtoms}},
\textnormal{\texttt{ex\_\allowbreak{}turns}}, \textnormal{\texttt{ex\_\allowbreak{}quarter\_\allowbreak{}sq}}, \textnormal{\texttt{ex\_\allowbreak{}quarters}}, \textnormal{\texttt{om\_\allowbreak{}quarter}}.

\paragraph{Phases on a support.} \textnormal{\texttt{mHas a m}} holds when $m$ contains every variable of $a$;
$\textnormal{\texttt{psupP a b P}}$ multiplies every monomial containing $a$ by $b$. \textnormal{\texttt{psupOK a b s}}: some
constraint normalizes to $a+ab$ (so $a(1+b)=0$ on the support, \textnormal{\texttt{pn\_\allowbreak{}psupP}}).
\textnormal{\texttt{psupS a b s}} applies \textnormal{\texttt{psupP}} to the phase; \textnormal{\texttt{psup\_\allowbreak{}sound}}: value unchanged.

\paragraph{Phases into atoms.} \textnormal{\texttt{aphaseOK m e t}}: the inner path is empty, $m\ne[\,]$, and
$(m,e)$ is a phase term. \textnormal{\texttt{aphaseS m e j t}} erases $(m,e)$ from the phase and appends the
atom $(j,[m])$. \textnormal{\texttt{aphase\_\allowbreak{}sound}}: if $A_j=\omega_M^e$, value unchanged.

\paragraph{Flipped atoms.} \textnormal{\texttt{aflipOK j t}} is \textnormal{\texttt{aheadOK j t}}; \textnormal{\texttt{aflipS j' t}} replaces the
head atom $(j,f)$ by $(j',\textnormal{\texttt{normF (\allowbreak{}1 + f)}})$. \textnormal{\texttt{aflip\_\allowbreak{}sound}}: if $A_jA_{j'} = 1$, the factor
is $A_j$ ($A_j^{[f]} = A_j\,(A_j^{-1})^{[1+f]}$).

\paragraph{Atom monomials off the constraints.} \textnormal{\texttt{adropCS t}} removes from every atom the
monomials $m$ with \textnormal{\texttt{dropT t.\allowbreak{}s.\allowbreak{}cons m}}; \textnormal{\texttt{adropC\_\allowbreak{}sound}} (inner path empty): value unchanged
(\textnormal{\texttt{fz\_\allowbreak{}filter\_\allowbreak{}zero}}, \textnormal{\texttt{mz\_\allowbreak{}of\_\allowbreak{}dropT}}).

\paragraph{[Average] on a context.} For a list of polynomials $Cs$ (low first) and
$Z\subseteq\N$: $\textnormal{\texttt{tSel Cs Z}}$ is the indicator that $Cs$ spells a number in $Z$
(\textnormal{\texttt{fz\_\allowbreak{}tSel}}). For triples $(j,Z_0,Z_1)$: \textnormal{\texttt{ctx0 ats n}} are the indices whose $Z_0$ holds
$n$, \textnormal{\texttt{ctx1 ats n}} those whose $Z_1$ holds $n$. \textnormal{\texttt{ctxOK y Cs as zs}} pairs atoms and
triples one to one: equal indices, $\textnormal{\texttt{normF (\allowbreak{}substF y [] f)}} = \textnormal{\texttt{normF (\allowbreak{}tSel Cs Z0)}}$, and
$\textnormal{\texttt{normF (\allowbreak{}substF y [[]] f)}} = \textnormal{\texttt{normF (\allowbreak{}tSel Cs Z1)}}$ (lists of unequal length fail).
\textnormal{\texttt{avctxOK M y Cs ats Qs t}}: $0<M$; $y$ the head of held and not in its tail, not on the path, in
no constraint or row; \textnormal{\texttt{halfOK M y phase}}; the quotient off the path and equal after
\textnormal{\texttt{normF}} to \textnormal{\texttt{tSel Cs Qs}}; and \textnormal{\texttt{ctxOK y Cs (\allowbreak{}atoms reading y) ats}}.
$\textnormal{\texttt{avctxS y Cs js t}}$ removes the terms of $y$, drops $y$ from held, and replaces the atoms
reading $y$ by $(js_n, \textnormal{\texttt{normF (\allowbreak{}tInd Cs n)}})$ for $n<2^{|Cs|}$ (\textnormal{\texttt{js.\allowbreak{}getD n 0}}).
\begin{theorem}[\textnormal{\texttt{avctx\_\allowbreak{}sound}}]
If \textnormal{\texttt{avctxOK M y Cs ats Qs t}} and, for every $n<2^{|Cs|}$,
$2A_{js_n} = \prod_{i\in\textnormal{\texttt{ctx0 ats n}}}A_i + (-1)^{[n\in Qs]}\prod_{i\in\textnormal{\texttt{ctx1 ats n}}}A_i$,
then the value equals \textnormal{\texttt{valA M (\allowbreak{}2 * c) A (\allowbreak{}avctxS y Cs js t) \ensuremath{\sigma}}}.
\end{theorem}
\noindent \textnormal{\texttt{avctxReqs k ats Qs js jn}} lists these identities as coefficient requirements
(below): first $(0,jn,[[],[jn]])$ ($A_{jn}=-1$), then for each $n<2^k$
$(2, js_n, [\textnormal{\texttt{ctx0 ats n}}, \textnormal{\texttt{ctx1 ats n}} \mathbin{+\!\!+} ([jn]\text{ if } n\in Qs)])$.

\paragraph{Atoms split on a support.} $\textnormal{\texttt{pairsOK C [m\_\allowbreak{}1,\allowbreak{}\textbackslash{}dots]}}$ holds when every product
$m_im_{i'}$ ($i<i'$) satisfies \textnormal{\texttt{dropT C}}. \textnormal{\texttt{asplitOK t}}: inner path empty and the head
atom's normalized monomials satisfy \textnormal{\texttt{pairsOK t.\allowbreak{}s.\allowbreak{}cons}}; \textnormal{\texttt{asplitS}} replaces the head atom
$(j,f)$ by one atom $(j,[m])$ per monomial of \textnormal{\texttt{normF f}}. \textnormal{\texttt{asplit\_\allowbreak{}sound}}: value
unchanged (\textnormal{\texttt{aval\_\allowbreak{}split1}}: on the support at most one monomial is $1$).

\paragraph{Final comparison.} \textnormal{\texttt{afinal\_\allowbreak{}sound}} ($0<M$): both held lists empty, equal atom lists,
\textnormal{\texttt{finalOK tL.\allowbreak{}s tR.\allowbreak{}s}}, the \textnormal{\texttt{constK}} values $k_L,k_R$, and
$c_L\omega_M^{k_L} = c_R\omega_M^{k_R}$ give $\textnormal{\texttt{valA M cL A tL \ensuremath{\sigma}}} = \textnormal{\texttt{valA M cR A tR \ensuremath{\sigma}}}$.
\textnormal{\texttt{afinal\_\allowbreak{}perm}} replaces equal atoms by \textnormal{\texttt{decide (\allowbreak{}tL.\allowbreak{}atoms.\allowbreak{}Perm tR.\allowbreak{}atoms)}}.

\subsection{The step language \texorpdfstring{\textnormal{\texttt{AStep}}}{AStep}}

\begin{definition}[Coefficient requirements]\label{def:lean-pathsum-req}
A \term{coefficient requirement} \textnormal{\texttt{Req := \ensuremath{\mathbb{N}} \ensuremath{\times} \ensuremath{\mathbb{N}} \ensuremath{\times} List (\allowbreak{}List \ensuremath{\mathbb{N}})}} is a triple $(k,j,ps)$
asserting \textnormal{\texttt{reqHolds A r}}: $k\,A_j = \sum_{p\in ps}\prod_{i\in p}A_i$. A
\term{root requirement} is a triple $(e,j,b)$ asserting \textnormal{\texttt{qreqHolds M A p}}:
$(b+1)A_j = b+\omega_M^e$.
\end{definition}

\begin{xltabular}{\textwidth}{@{}>{\raggedright\arraybackslash}p{3.7cm}>{\raggedright\arraybackslash}X>{\raggedright\arraybackslash}p{1.5cm}>{\raggedright\arraybackslash}p{3.9cm}@{}}
\caption{The constructors of \texttt{PSem.AStep}: side condition (\texttt{AStep.ok M}), transformer
(\texttt{AStep.apply M}), factor (\texttt{AStep.fac A}), and requirements (\texttt{AStep.reqs},
\texttt{AStep.qreqs M}).}\label{tab:lean-pathsum-asteps}\\
\toprule Constructor & Side condition / transformer & Factor & Requirements\\\midrule\endfirsthead
\toprule Constructor & Side condition / transformer & Factor & Requirements\\\midrule\endhead
\textnormal{\texttt{inner st}} & \textnormal{\texttt{stepOK M st t.\allowbreak{}s}} / \textnormal{\texttt{stepS}} on \textnormal{\texttt{t.\allowbreak{}s}} & \textnormal{\texttt{fac st}} & --\\
\textnormal{\texttt{hperm l}} & \textnormal{\texttt{decide (\allowbreak{}t.\allowbreak{}held.\allowbreak{}Perm l)}} / \textnormal{\texttt{hpermS l}} & $1$ & --\\
\textnormal{\texttt{release}} & \textnormal{\texttt{releaseOK}} / \textnormal{\texttt{releaseS}} & $1$ & --\\
\textnormal{\texttt{hold}} & \textnormal{\texttt{holdOK}} / \textnormal{\texttt{holdS}} & $1$ & --\\
\textnormal{\texttt{asub v}} & \textnormal{\texttt{asubOK v}} / \textnormal{\texttt{asubS v}} & $1$ & --\\
\textnormal{\texttt{asubB v B}} & \textnormal{\texttt{asubBOK v B}} / \textnormal{\texttt{asubBS v B}} & $1$ & --\\
\textnormal{\texttt{average y g j as bs}} & \textnormal{\texttt{averageOK y g}}, the alpha indices $=as$, the other indices $=bs$ / \textnormal{\texttt{averageS y g j}} & $2$ & $(2,j,[as,bs])$\\
\textnormal{\texttt{averageH y j0 j1 jn as bs}} & \textnormal{\texttt{averageHOK M y}}, indices as above with $g=1$ / \textnormal{\texttt{averageHS y j0 j1}} & $2$ & $(2,j_0,[as,bs])$, $(2,j_1,[as,bs\mathbin{+\!\!+}[jn]])$, $(0,jn,[[],[jn]])$\\
\textnormal{\texttt{azero j}} & \textnormal{\texttt{azeroOK j}} / \textnormal{\texttt{azeroS}} & $1$ & $(1,j,[])$\\
\textnormal{\texttt{aconst j}} & \textnormal{\texttt{aconstOK}}, \textnormal{\texttt{aconstIdx t = j}} / \textnormal{\texttt{aconstS}} & $A_j$ & --\\
\textnormal{\texttt{aone j}} & \textnormal{\texttt{aheadOK j}} / \textnormal{\texttt{aconstS}} & $1$ & $(1,j,[[]])$\\
\textnormal{\texttt{adrop}} & \textnormal{\texttt{true}} / \textnormal{\texttt{adropS}} & $1$ & --\\
\textnormal{\texttt{amerge j j1 j2}} & \textnormal{\texttt{amergeOK j1 j2}} / \textnormal{\texttt{amergeS j}} & $1$ & $(1,j,[[j_1,j_2]])$\\
\textnormal{\texttt{apick i}} & $i<|\text{atoms}|$ / \textnormal{\texttt{apickS i}} & $1$ & --\\
\textnormal{\texttt{anorm}} & \textnormal{\texttt{true}} / \textnormal{\texttt{anormS}} & $1$ & --\\
\textnormal{\texttt{quarter B0 B1 j}} & \textnormal{\texttt{quarterOK M B0 B1}} / \textnormal{\texttt{quarterS M B0 B1 j}} & $2$ & root $(M/4, j, 1)$\\
\textnormal{\texttt{acomp j}} & \textnormal{\texttt{acompOK j}} / \textnormal{\texttt{acompS}} & $A_j$ & --\\
\textnormal{\texttt{turn Bs ats}} & \textnormal{\texttt{turnOK M Bs ats}} / \textnormal{\texttt{turnS M Bs ats}} & $2$ & roots $(s\,(M/2^{|Bs|}), j, 1)$ per $(s,j)$\\
\textnormal{\texttt{psup a b}} & \textnormal{\texttt{psupOK a b t.\allowbreak{}s}} / \textnormal{\texttt{psupS a b}} on \textnormal{\texttt{t.\allowbreak{}s}} & $1$ & --\\
\textnormal{\texttt{aphase m e j}} & \textnormal{\texttt{aphaseOK m e}} / \textnormal{\texttt{aphaseS m e j}} & $1$ & root $(e,j,0)$\\
\textnormal{\texttt{aflip j j' o}} & \textnormal{\texttt{aflipOK j}} / \textnormal{\texttt{aflipS j'}} & $A_j$ & $(1,o,[[]])$, $(1,o,[[j,j']])$\\
\textnormal{\texttt{adropC}} & \textnormal{\texttt{decide (\allowbreak{}t.\allowbreak{}s.\allowbreak{}path = [])}} / \textnormal{\texttt{adropCS}} & $1$ & --\\
\textnormal{\texttt{acons j}} & \textnormal{\texttt{aconsOK}} / \textnormal{\texttt{aconsS j}} & $1$ & $(1,j,[])$\\
\textnormal{\texttt{avctx y Cs ats Qs js jn}} & \textnormal{\texttt{avctxOK M y Cs ats Qs}} / \textnormal{\texttt{avctxS y Cs js}} & $2$ & \textnormal{\texttt{avctxReqs |Cs| ats Qs js jn}}\\
\textnormal{\texttt{asplit}} & \textnormal{\texttt{asplitOK}} / \textnormal{\texttt{asplitS}} & $1$ & --\\
\bottomrule
\end{xltabular}

\noindent The requirement $(1,o,[[]])$ with $(1,o,[[j,j']])$ of \textnormal{\texttt{aflip}} states $A_o=1$ and
$A_o = A_jA_{j'}$; $(0,jn,[[],[jn]])$ states $0 = 1+A_{jn}$. A root requirement $(e,j,1)$ states
$2A_j = 1+\omega_M^e$, and $(e,j,0)$ states $A_j=\omega_M^e$.
\textnormal{\texttt{arunsOK}}, \textnormal{\texttt{aruns}}, and \textnormal{\texttt{afacs A}} compose the steps as \textnormal{\texttt{runOK}}, \textnormal{\texttt{run}},
and \textnormal{\texttt{facs}} do; \textnormal{\texttt{areqs sts}} and \textnormal{\texttt{aqreqs M sts}} concatenate the steps' requirements.

\begin{theorem}[\textnormal{\texttt{AStep.\allowbreak{}sound}}, \textnormal{\texttt{aruns\_\allowbreak{}sound}}]\label{thm:lean-pathsum-aruns}
For $2\le M$: if every requirement of a step holds (\textnormal{\texttt{\ensuremath{\forall} r \ensuremath{\in} st.\allowbreak{}reqs,\allowbreak{} reqHolds A r}} and
\textnormal{\texttt{\ensuremath{\forall} p \ensuremath{\in} st.\allowbreak{}qreqs M,\allowbreak{} qreqHolds M A p}}) and \textnormal{\texttt{st.\allowbreak{}ok M t}}, then
$\textnormal{\texttt{valA M c A t \ensuremath{\sigma}}} = \textnormal{\texttt{valA M (\allowbreak{}st.\allowbreak{}fac A * c) A (\allowbreak{}st.\allowbreak{}apply M t) \ensuremath{\sigma}}}$. For a list:
\textnormal{\texttt{(\allowbreak{}\ensuremath{\forall} r \ensuremath{\in} areqs sts,\allowbreak{} reqHolds A r) \ensuremath{\to} (\allowbreak{}\ensuremath{\forall} p \ensuremath{\in} aqreqs M sts,\allowbreak{} qreqHolds M A p) \ensuremath{\to} arunsOK M sts t = true}}
implies $\textnormal{\texttt{valA M c A t \ensuremath{\sigma}}} = \textnormal{\texttt{valA M (\allowbreak{}afacs A sts * c) A (\allowbreak{}aruns M sts t) \ensuremath{\sigma}}}$.
\end{theorem}

\section{Leftover supports: \texorpdfstring{\file{PathSumSupport}}{PathSumSupport}}\label{sec:lean-pathsum-support}

When Boolean elimination (\textnormal{\texttt{gather}}, \textnormal{\texttt{bcanon}}) leaves a support indicator other than
$1$, both reduced sums keep one equal constraint list. \textnormal{\texttt{finalSOK sL sR}} holds when both paths
are empty and the constraint lists, the row lists, and the stripped phases are equal.
\begin{theorem}[\textnormal{\texttt{finalS\_\allowbreak{}sound}}, \textnormal{\texttt{afinalS\_\allowbreak{}perm}}]
$0<M$, \textnormal{\texttt{finalSOK sL sR = true}}, $\textnormal{\texttt{constK sL.\allowbreak{}phase}} = k_L$, $\textnormal{\texttt{constK sR.\allowbreak{}phase}} = k_R$,
and $c_L\omega_M^{k_L} = c_R\omega_M^{k_R}$ give $\textnormal{\texttt{val M cL sL \ensuremath{\sigma}}} = \textnormal{\texttt{val M cR sR \ensuremath{\sigma}}}$.
\textnormal{\texttt{afinalS\_\allowbreak{}perm}} is the version with atoms: empty held lists and
\textnormal{\texttt{decide (\allowbreak{}tL.\allowbreak{}atoms.\allowbreak{}Perm tR.\allowbreak{}atoms)}}.
\end{theorem}

\section{Checkpointed step lists: \texorpdfstring{\file{PathSumChain}}{PathSumChain} and \texorpdfstring{\file{PathSumQuditChain}}{PathSumQuditChain}}\label{sec:lean-pathsum-chain}

A certificate needs \textnormal{\texttt{runOK M sts s = true}} and the literal value of \textnormal{\texttt{run M sts s}}.
Evaluating a long step list in one \textnormal{\texttt{decide}} makes the kernel hold every intermediate state at
once; these modules split the list into blocks, state each intermediate value as a literal
definition, and check each block in its own declaration.

\begin{theorem}[Cutting a step list]\label{thm:lean-pathsum-cut}
\textnormal{\texttt{runOK\_\allowbreak{}append}}: $\textnormal{\texttt{runOK M (\allowbreak{}a ++ b) s}} = \textnormal{\texttt{runOK M a s \&\& runOK M b (\allowbreak{}run M a s)}}$;
\textnormal{\texttt{run\_\allowbreak{}append}}: $\textnormal{\texttt{run M (\allowbreak{}a ++ b) s}} = \textnormal{\texttt{run M b (\allowbreak{}run M a s)}}$;
\textnormal{\texttt{run\_\allowbreak{}cut}}: if \textnormal{\texttt{(\allowbreak{}runOK M (\allowbreak{}sts.\allowbreak{}take k) s \&\& decide (\allowbreak{}run M (\allowbreak{}sts.\allowbreak{}take k) s = s')) = true}}
then $\textnormal{\texttt{runOK M sts s}} = \textnormal{\texttt{runOK M (\allowbreak{}sts.\allowbreak{}drop k) s'}}$ and
$\textnormal{\texttt{run M sts s}} = \textnormal{\texttt{run M (\allowbreak{}sts.\allowbreak{}drop k) s'}}$. \textnormal{\texttt{arunsOK\_\allowbreak{}append}}, \textnormal{\texttt{aruns\_\allowbreak{}append}},
and \textnormal{\texttt{aruns\_\allowbreak{}cut}} are the same for \textnormal{\texttt{AStep}} (with \textnormal{\texttt{AData}} given \textnormal{\texttt{DecidableEq}} by
\textnormal{\texttt{deriving instance}}), and \textnormal{\texttt{QSem.\allowbreak{}runOK\_\allowbreak{}append}}, \textnormal{\texttt{QSem.\allowbreak{}run\_\allowbreak{}append}},
\textnormal{\texttt{QSem.\allowbreak{}run\_\allowbreak{}cut}} for qudit steps.
\end{theorem}

\begin{definition}[State sizes]\label{def:lean-pathsum-size}
\textnormal{\texttt{Data.\allowbreak{}size d}} $= |\text{path}| + \sum_{(m,k)\in\text{phase}}(|m|+1) + \sum_{q\in\text{cons}\mathbin{+\!\!+}\text{rows}}\bigl(\sum_{m\in q}|m| + |q|\bigr)$;
\textnormal{\texttt{AData.\allowbreak{}size t}} $=\textnormal{\texttt{t.\allowbreak{}s.\allowbreak{}size}} + |\text{held}| + \sum_{(j,f)\in\text{atoms}}(\sum_{m\in f}|m| + |f| + 1)$;
\textnormal{\texttt{QSem.\allowbreak{}Data.\allowbreak{}size d}} $=|\text{path}| + \sum_{P\in\text{phase}::\text{cons}\mathbin{+\!\!+}\text{rows}}\sum_{(m,k)\in P}(2|m|+1)$.
\end{definition}

\begin{algorithm}[\textnormal{\texttt{blocks step size k w s0 l}}]\label{alg:lean-pathsum-blocks}
Input: a one-step function, a size function, a block length $k$, a size bound $w$, an initial
state, and a step list. Output: an array of pairs (state after the block, number of steps in the
block).
\begin{enumerate}
\item Set $cur := s_0$, $n := 0$, $out := [\,]$.
\item For each step $st$ of $l$ in order: let $nxt := \mathit{step}(st, cur)$.
  \begin{enumerate}
  \item If $\mathit{size}(cur)+\mathit{size}(nxt) > w$: if $n>0$ push $(cur,n)$; push $(nxt,1)$;
    set $n:=0$.
  \item Otherwise set $n := n+1$; if $n = k$, push $(nxt, n)$ and set $n := 0$.
  \item Set $cur := nxt$.
  \end{enumerate}
\item If $n>0$ push $(cur,n)$. Return $out$.
\end{enumerate}
\end{algorithm}

\paragraph{Literal conversion.} \textnormal{\texttt{listE}}, \textnormal{\texttt{natsE}}, \textnormal{\texttt{f2E}}, \textnormal{\texttt{dataE}}, \textnormal{\texttt{adataE}}
(and \textnormal{\texttt{QSem.\allowbreak{}polyE}}, \textnormal{\texttt{QSem.\allowbreak{}qdataE}}) build \textnormal{\texttt{Expr}} terms of raw numerals
(\textnormal{\texttt{mkRawNatLit}}) and \textnormal{\texttt{List.\allowbreak{}cons}}/\textnormal{\texttt{List.\allowbreak{}nil}}/\textnormal{\texttt{Prod.\allowbreak{}mk}}/\textnormal{\texttt{Data.\allowbreak{}mk}}/\textnormal{\texttt{AData.\allowbreak{}mk}}
applications. \textnormal{\texttt{knorm e}} reduces $e$ by the kernel's weak head normal form
(\textnormal{\texttt{Kernel.\allowbreak{}whnf}}) and recursively each constructor argument after the parameters, folding
\textnormal{\texttt{Nat.\allowbreak{}succ}} of a literal into a literal; it throws ``\texttt{knorm: the kernel cannot reduce}''
or ``\texttt{knorm: not a constructor}''. \textnormal{\texttt{decNat}}, \textnormal{\texttt{decList}}, \textnormal{\texttt{decPair}},
\textnormal{\texttt{decData}}, \textnormal{\texttt{decAData}}, and \textnormal{\texttt{QSem.\allowbreak{}decQData}} read such terms back into values, with
the errors ``\texttt{decNat: not a numeral}'', ``\texttt{decList: not a list}'',
``\texttt{decPair: not a pair}'', ``\texttt{decData: not a Data}'',
``\texttt{decAData: not an AData}'', and ``\texttt{decQData: not a Data}''.
\textnormal{\texttt{dataStates M k w s0 sts}} (\textnormal{\texttt{adataStates}}, \textnormal{\texttt{QSem.\allowbreak{}qdataStates}}) is an opaque
\textnormal{\texttt{MetaM}} function implemented by compiled code (\textnormal{\texttt{@[implemented\_\allowbreak{}by]}}): it reads the start
state, evaluates the step list with \textnormal{\texttt{evalExpr}}, runs \cref{alg:lean-pathsum-blocks} with
$\textnormal{\texttt{step st d := run M [st] d}}$ and the matching size, and returns the block states as literals.

\begin{algorithm}[\textnormal{\texttt{chain kind states n M sts start k w}}]\label{alg:lean-pathsum-chain}
\textnormal{\texttt{ChainKind}} names the step type, the state type, the checker, the runner, and the cut lemma
(\textnormal{\texttt{\ensuremath{\langle}Step,\allowbreak{} Data,\allowbreak{} runOK,\allowbreak{} run,\allowbreak{} run\_\allowbreak{}cut\ensuremath{\rangle}}}, \textnormal{\texttt{\ensuremath{\langle}AStep,\allowbreak{} AData,\allowbreak{} arunsOK,\allowbreak{} aruns,\allowbreak{} aruns\_\allowbreak{}cut\ensuremath{\rangle}}}, or the
\textnormal{\texttt{QSem}} names). With prefix $n$:
\begin{enumerate}
\item Evaluate $M$ to a numeral (else ``\texttt{chain: ... is not a numeral}''). Normalize the start
  by \textnormal{\texttt{knorm}} and add it as the abbreviation \textnormal{\texttt{n\_\allowbreak{}s0}}; add \textnormal{\texttt{n\_\allowbreak{}e0 : start = n\_\allowbreak{}s0}} proved by
  \textnormal{\texttt{decide}} (\textnormal{\texttt{addDecided}}: \textnormal{\texttt{mkDecideProof}}, elaborated synchronously).
\item Compute the blocks $(v_i, k_i)$ with \textnormal{\texttt{states M k w n\_\allowbreak{}s0 sts}}.
\item With $R := \text{sts}$, for each block $i$: add \textnormal{\texttt{n\_\allowbreak{}s\{i+1\} := v\_\allowbreak{}i}}; add
  \textnormal{\texttt{n\_\allowbreak{}h\{i\} : (\allowbreak{}ok M (\allowbreak{}R.\allowbreak{}take k\_\allowbreak{}i) n\_\allowbreak{}s\{i\} \&\& decide (\allowbreak{}run M (\allowbreak{}R.\allowbreak{}take k\_\allowbreak{}i) n\_\allowbreak{}s\{i\} = n\_\allowbreak{}s\{i+1\})) = true}}
  by \textnormal{\texttt{decide}}; chain the cut lemma's two conclusions onto the running equalities by
  transitivity; set $R := R.\textnormal{\texttt{drop}}\ k_i$.
\item Add \textnormal{\texttt{n\_\allowbreak{}hend : ok M R n\_\allowbreak{}s\{last\} = true}} and \textnormal{\texttt{n\_\allowbreak{}qend : run M R n\_\allowbreak{}s\{last\} = n\_\allowbreak{}s\{last\}}}
  by \textnormal{\texttt{decide}} (the remaining list is empty), and the abbreviation \textnormal{\texttt{n := n\_\allowbreak{}s\{last\}}}.
\item Add \textnormal{\texttt{n\_\allowbreak{}ok : ok M sts start = true}} and \textnormal{\texttt{n\_\allowbreak{}eq : run M sts start = n}}, each the
  congruence through \textnormal{\texttt{n\_\allowbreak{}e0}} composed with the chained equalities.
\end{enumerate}
\end{algorithm}

\paragraph{Commands.} The command syntax is
\begin{lstlisting}[style=grammar]
run_chain   <ident> <M> <sts> <s> <k> [<w>]    -- PSem.Step on Data
aruns_chain <ident> <M> <sts> <t> <k> [<w>]    -- PSem.AStep on AData
qrun_chain  <ident> <p> <sts> <s> <k>          -- QSem.Step on QSem.Data
\end{lstlisting}
Each elaborates its terms against the expected types, runs \cref{alg:lean-pathsum-chain} in the
current namespace with block length $k$ and size bound $w$ (default $2^{62}$, i.e.\ no size cut;
\textnormal{\texttt{qrun\_\allowbreak{}chain}} always uses $2^{62}$), and declares \textnormal{\texttt{<ident>\_\allowbreak{}ok}}, \textnormal{\texttt{<ident>\_\allowbreak{}eq}}, and
\textnormal{\texttt{<ident>}} together with the auxiliary \textnormal{\texttt{\_\allowbreak{}s}}$i$, \textnormal{\texttt{\_\allowbreak{}h}}$i$, \textnormal{\texttt{\_\allowbreak{}e0}}, \textnormal{\texttt{\_\allowbreak{}hend}},
\textnormal{\texttt{\_\allowbreak{}qend}} declarations.

\section{Qubit networks as path sums: \texorpdfstring{\file{PathSumNet}}{PathSumNet}}\label{sec:lean-pathsum-net}

\file{PathSumNet} (imports \file{PathSumRules}) connects path sums to the tensor networks of
\file{Net} (\cref{ch:lean-core}). There a network \textnormal{\texttt{Net}} has \textnormal{\texttt{nodes : List Node}}, summed
labels \textnormal{\texttt{int}}, and boundary labels \textnormal{\texttt{bnd}}; a node has a \textnormal{\texttt{kind}}, output and input
labels, a dimension index, and a phase index; an environment \textnormal{\texttt{Env}} gives \textnormal{\texttt{dimv}},
phase vectors \textnormal{\texttt{phase : \ensuremath{\mathbb{N}} \ensuremath{\to} \ensuremath{\mathbb{N}} \ensuremath{\to} \ensuremath{\mathbb{R}}}} in turns, the label table \textnormal{\texttt{labs}},
multiplicities, and coefficient vectors \textnormal{\texttt{coef}}. The entry of a node under a label assignment
$\sigma:\N\to\N$ is \textnormal{\texttt{ArchytasZX.\allowbreak{}val E \ensuremath{\sigma} n}}; for the kinds used here,
\begin{align*}
\text{Z:}\ & \textstyle\sum_{k<d}[\forall v\in vs,\ v=k]\operatorname{ex}(\varphi_k),\\
\text{X:}\ & \textstyle\sum_{k<d}\operatorname{ex}(\varphi_k)\prod_{o}\textnormal{\texttt{fou d o k}}\prod_i\overline{\textnormal{\texttt{fou d i k}}},\\
\text{F:}\ & \textnormal{\texttt{fou d o i}},\\
\text{W (input } i\text{):}\ & [i=0]\,[\forall o,\ o=0] + [0<i<d]\,[\text{the nonzero outputs are } [i]],\\
\text{B, S:}\ & [a<s][b<t][k=at+b],
\end{align*}
with B on outputs $[k]$ and inputs $[a,b]$, S on outputs $[a,b]$ and input $[k]$, $s$ the
dimension value of the node's \textnormal{\texttt{dim}} and $t$ that of its \textnormal{\texttt{ph}} index.
\textnormal{\texttt{ZW}}/\textnormal{\texttt{XW}} use the coefficient vector in place of $\operatorname{ex}(\varphi_k)$.
$\textnormal{\texttt{eval E N \ensuremath{\sigma}}}$ sums the product of the node entries over the summed labels, each label
over its dimension \textnormal{\texttt{labDim E l}}, and \textnormal{\texttt{Equal E s\ensuremath{_1} N\ensuremath{_1} s\ensuremath{_2} N\ensuremath{_2}}} states that the boundaries
are permutations of each other and $s_1\,\textnormal{\texttt{eval E N\ensuremath{_1} \ensuremath{\sigma}}} = s_2\,\textnormal{\texttt{eval E N\ensuremath{_2} \ensuremath{\sigma}}}$ at every
in-range boundary assignment.

\begin{definition}[Fragments]\label{def:lean-pathsum-frag}
\textnormal{\texttt{structure Frag}} has fields \textnormal{\texttt{path : List \ensuremath{\mathbb{N}}}}, \textnormal{\texttt{phase : Poly}}, \textnormal{\texttt{cons : List F2}},
and \textnormal{\texttt{c : \ensuremath{\mathbb{C}}}}: one generator's path sum. $\textnormal{\texttt{fval M f \ensuremath{\tau}}} = \textnormal{\texttt{val M f.\allowbreak{}c \ensuremath{\langle}f.\allowbreak{}path,\allowbreak{} f.\allowbreak{}phase,\allowbreak{} f.\allowbreak{}cons,\allowbreak{} []\ensuremath{\rangle} \ensuremath{\tau}}}$.
\textnormal{\texttt{Frag.\allowbreak{}join a b}} concatenates paths, phases, and constraints and multiplies the constants;
\textnormal{\texttt{Frag.\allowbreak{}unit}} $=\langle[\,],[\,],[\,],1\rangle$; \textnormal{\texttt{joinAll}} folds \textnormal{\texttt{join}} from the right
over a list, ending in \textnormal{\texttt{unit}}. \textnormal{\texttt{Disj a b}} holds when no path variable of $a$ is a path
variable, a phase variable, or a constraint variable of $b$ (decidable instance provided).
\end{definition}

\begin{theorem}[Products of fragments]
\textnormal{\texttt{fval\_\allowbreak{}join}}: \textnormal{\texttt{Disj a b}} and \textnormal{\texttt{Disj b a}} imply
$\textnormal{\texttt{fval M (\allowbreak{}a.\allowbreak{}join b) \ensuremath{\tau}}} = \textnormal{\texttt{fval M a \ensuremath{\tau}}} \cdot \textnormal{\texttt{fval M b \ensuremath{\tau}}}$.
\textnormal{\texttt{fval\_\allowbreak{}joinAll}}: for a list whose pairs are disjoint both ways
(\textnormal{\texttt{fs.\allowbreak{}Pairwise (\allowbreak{}fun a b => Disj a b \ensuremath{\wedge} Disj b a)}}),
$\textnormal{\texttt{fval M (\allowbreak{}joinAll fs) \ensuremath{\tau}}} = \prod_{f\in fs}\textnormal{\texttt{fval M f \ensuremath{\tau}}}$; \textnormal{\texttt{fval\_\allowbreak{}unit}} is $1$.
\end{theorem}
\noindent Supporting lemmas: \textnormal{\texttt{bsum\_\allowbreak{}inv}}, \textnormal{\texttt{summand\_\allowbreak{}join}} (the summand of a join is the
product of the summands), \textnormal{\texttt{summand\_\allowbreak{}frag\_\allowbreak{}inv}}, \textnormal{\texttt{pvars\_\allowbreak{}append}}, \textnormal{\texttt{fvars\_\allowbreak{}append}},
\textnormal{\texttt{disj\_\allowbreak{}join\_\allowbreak{}right}}, \textnormal{\texttt{disj\_\allowbreak{}join\_\allowbreak{}left}}, \textnormal{\texttt{disj\_\allowbreak{}unit\_\allowbreak{}left}}, \textnormal{\texttt{disj\_\allowbreak{}unit\_\allowbreak{}right}},
\textnormal{\texttt{disj\_\allowbreak{}joinAll\_\allowbreak{}right}}, \textnormal{\texttt{disj\_\allowbreak{}joinAll\_\allowbreak{}left}}.

\begin{definition}[Disjointness in one pass]\label{def:lean-pathsum-disjb}
\begin{itemize}
\item \textnormal{\texttt{nodupB xs acc}} walks $xs$ with a bit mask of the values seen: false at the first
  $x$ with bit $x$ of $acc$ set, else $acc:=acc\mathbin{|}2^x$. \lean!nodupB_sound!: true implies $xs$ has
  no repeats and sets no bit of $acc$. \textnormal{\texttt{nodupF xs}} is \textnormal{\texttt{nodupB xs 0}} when every value is
  below $2^{24}$, else \lean!decide xs.Nodup! (\lean!nodupF_nodup!).
\item \textnormal{\texttt{maskOf xs acc}} sets the bits of $xs$ in $acc$ (\lean!testBit_maskOf!);
  \textnormal{\texttt{fragVars f = f.path ++ pvars f.phase ++ fvars f.cons}}.
\item \textnormal{\texttt{disjB fs A B}} with $A$ the mask of the earlier fragments' variables and $B$ of their
  path variables: for $f::fs$, the path mask of $f$ meets no bit of $A$, the variable mask of $f$
  meets no bit of $B$, and \textnormal{\texttt{disjB fs (A ||| vars f) (B ||| path f)}} holds.
  \lean!disjB_sound!, \lean!land_zero_bit!, and \lean!disjB_pairwise!: \textnormal{\texttt{disjB fs 0 0 = true}}
  implies \textnormal{\texttt{fs.Pairwise (fun a b => Disj a b \ensuremath{\wedge} Disj b a)}}, the hypothesis of
  \lean!fval_joinAll!.
\item \textnormal{\texttt{allIn ys xs}}: every value of $xs$ has its bit set in \textnormal{\texttt{maskOf ys 0}}
  (\lean!allIn_sound!); \lean!cover_of!: \textnormal{\texttt{allIn (N.int ++ N.bnd) (labelsOf N.nodes)}} gives every
  node label summed or boundary; \lean!sub_of!: \textnormal{\texttt{allIn labs (N.int ++ N.bnd)}} gives every
  summed or boundary label in $labs$.
\end{itemize}
\end{definition}

\paragraph{Labels of dimension 2.} $\textnormal{\texttt{bits \ensuremath{\sigma} l}} = (\sigma(l) : \Zd2)$; \textnormal{\texttt{bits\_\allowbreak{}update}} commutes
it with updates. \textnormal{\texttt{sumOver\_\allowbreak{}bsum}}: if $f(\sigma) = g(\textnormal{\texttt{bits \ensuremath{\sigma}}})$ whenever every $\sigma(l)<2$, and
every label of $ls$ has dimension $2$, then $\textnormal{\texttt{sumOver D ls f \ensuremath{\sigma}}} = \textnormal{\texttt{bsum ls g (\allowbreak{}bits \ensuremath{\sigma})}}$ for
$\sigma$ in range. \textnormal{\texttt{eval\_\allowbreak{}frags}} states, for a network whose nodes are the first components of
a list $ps$ of (node, fragment) pairs with $\textnormal{\texttt{val E \ensuremath{\sigma} n}} = \textnormal{\texttt{fval M f (\allowbreak{}bits \ensuremath{\sigma})}}$ on in-range
assignments, pairwise disjoint fragments, and summed labels of dimension $2$:
$\textnormal{\texttt{eval E N \ensuremath{\sigma}}}$ equals \textnormal{\texttt{val M (\allowbreak{}joinAll fs).\allowbreak{}c d (\allowbreak{}bits \ensuremath{\sigma})}} with $fs$ the fragments of $ps$
and $d = \langle N.\text{int}\mathbin{+\!\!+}J.\text{path},\ J.\text{phase},\ J.\text{cons},\ [\,]\rangle$,
$J = \textnormal{\texttt{joinAll fs}}$.

\begin{definition}[Fragments of the generators at dimension 2]\label{def:lean-pathsum-frag2}
For a phase modulus $M$ and phase table entries $t_0 = M\varphi_0$, $t_1 = M\varphi_1$:
\begin{align*}
\textnormal{\texttt{zFrag M t0 t1 l0 ls}} &= \langle[\,],\ [([\,],t_0), ([l_0], t_1+(M-t_0))],\ [[[l_0],[l]] : l\in ls],\ 1\rangle,\\
\textnormal{\texttt{fFrag M o i}} &= \langle[\,],\ [([o,i], M/2)],\ [\,],\ 1/\sqrt2\rangle,\\
\textnormal{\texttt{xFrag M t0 t1 y ls}} &= \langle[y],\ [([\,],t_0),([y],t_1+(M-t_0))]\mathbin{+\!\!+}[([y,l],M/2): l\in ls],\\
 &\qquad [\,],\ (1/\sqrt2)^{|ls|}\rangle .
\end{align*}
The Z fragment is the constraints $l_0+l=0$ and the phase $t_0+(t_1-t_0)l_0$; the X fragment adds a
fresh variable $y$ with phase $t_0+(t_1-t_0)y+\sum_l yl/2$.
\end{definition}

\begin{theorem}[Generator fragments]
Let $\sigma(l)<2$ for every $l$, the node's dimension value be $2$, $0<M$, $t_0\le M$, and
$\varphi_0 = t_0/M$, $\varphi_1 = t_1/M$.
\begin{itemize}
\item \textnormal{\texttt{z\_\allowbreak{}frag}}: for a Z node with legs \textnormal{\texttt{outs ++ ins = l0 :: ls}},
  $\textnormal{\texttt{val E \ensuremath{\sigma} n}} = \textnormal{\texttt{fval M (\allowbreak{}zFrag M t0 t1 l0 ls) (\allowbreak{}bits \ensuremath{\sigma})}}$; \textnormal{\texttt{z\_\allowbreak{}frag\_\allowbreak{}at}}: the same with
  \textnormal{\texttt{zFrag M t0 t1 l0 (\allowbreak{}outs ++ ins)}} for any leg $l_0$ (\textnormal{\texttt{tZ\_\allowbreak{}congr}}, \textnormal{\texttt{tZ\_\allowbreak{}bits}}).
\item \textnormal{\texttt{f\_\allowbreak{}frag}}: for an F node with \textnormal{\texttt{outs = [o]}}, \textnormal{\texttt{ins = [i]}} and $2\mid M$,
  $\textnormal{\texttt{val E \ensuremath{\sigma} n}} = \textnormal{\texttt{fval M (\allowbreak{}fFrag M o i) (\allowbreak{}bits \ensuremath{\sigma})}}$ (\textnormal{\texttt{fou2}}:
  $\textnormal{\texttt{fou 2 j k}} = (-1)^{jk}/\sqrt2$).
\item \textnormal{\texttt{x\_\allowbreak{}frag}}: for an X node, $M=2r$, and $y$ not a leg,
  $\textnormal{\texttt{val E \ensuremath{\sigma} n}} = \textnormal{\texttt{fval M (\allowbreak{}xFrag M t0 t1 y (\allowbreak{}outs ++ ins)) (\allowbreak{}bits \ensuremath{\sigma})}}$ (\textnormal{\texttt{fou2\_\allowbreak{}ex}},
  \textnormal{\texttt{fou2\_\allowbreak{}conj}} (the qubit DFT entry is real), \textnormal{\texttt{ex\_\allowbreak{}half\_\allowbreak{}r}}, \textnormal{\texttt{prod\_\allowbreak{}ex}}).
\end{itemize}
\end{theorem}

\begin{definition}[Network path sums]\label{def:lean-pathsum-netps}
$\textnormal{\texttt{netData N fs}} = \langle N.\text{int}\mathbin{+\!\!+}(\textnormal{\texttt{joinAll fs}}).\text{path},\ (\textnormal{\texttt{joinAll fs}}).\text{phase},\ (\textnormal{\texttt{joinAll fs}}).\text{cons},\ [\,]\rangle$.
$\textnormal{\texttt{clamp bnd \ensuremath{\sigma} l}} = \sigma(l)$ for $l\in bnd$, else $0$.
\textnormal{\texttt{structure NetPS E N M ps : Prop}} has the fields
\textnormal{\texttt{nodes : N.\allowbreak{}nodes = ps.\allowbreak{}map Prod.\allowbreak{}fst}},
\textnormal{\texttt{node}} (each node's entry is its fragment's value at \textnormal{\texttt{bits \ensuremath{\sigma}}} for $\sigma<2$),
\textnormal{\texttt{disj}} (pairwise disjoint fragments, both ways),
\textnormal{\texttt{int}} and \textnormal{\texttt{bnd}} (every summed and every boundary label has dimension $2$), and
\textnormal{\texttt{cover}} (every label of the nodes is summed or on the boundary).
\end{definition}

\begin{theorem}[\textnormal{\texttt{NetPS.\allowbreak{}eval\_\allowbreak{}eq}}, \textnormal{\texttt{equal\_\allowbreak{}ps}}]
For \textnormal{\texttt{h : NetPS E N M ps}} and $\sigma$ in range on the boundary,
$\textnormal{\texttt{eval E N \ensuremath{\sigma}}} = \textnormal{\texttt{val M (\allowbreak{}joinAll fs).\allowbreak{}c (\allowbreak{}netData N fs) (\allowbreak{}bits (\allowbreak{}clamp N.\allowbreak{}bnd \ensuremath{\sigma}))}}$.
\textnormal{\texttt{equal\_\allowbreak{}ps}}: two networks with \textnormal{\texttt{NetPS}}, equal boundary lists, and
$s_1\,\textnormal{\texttt{val M c\ensuremath{_1} (\allowbreak{}netData N\ensuremath{_1} fs\ensuremath{_1}) \ensuremath{\tau}}} = s_2\,\textnormal{\texttt{val M c\ensuremath{_2} (\allowbreak{}netData N\ensuremath{_2} fs\ensuremath{_2}) \ensuremath{\tau}}}$ for every
$\tau:\N\to\Zd2$ satisfy \textnormal{\texttt{Equal E s\ensuremath{_1} N\ensuremath{_1} s\ensuremath{_2} N\ensuremath{_2}}}.
\end{theorem}

\section{Labels of dimension \texorpdfstring{$2^k$}{2\^{}k}: \texorpdfstring{\file{PathSumDigits}}{PathSumDigits}}\label{sec:lean-pathsum-digits}

\file{PathSumDigits} (imports \file{PathSumNet}, \file{Positional}) represents a label of
dimension $2^k$ by $k$ Boolean \term{digit variables}, low digit first. A function
$\mathit{dig}:\N\to$ \textnormal{\texttt{List \ensuremath{\mathbb{N}}}} gives each label's digit variables.

\begin{definition}[Digit assignments]\label{def:lean-pathsum-digit-assign}
\begin{itemize}
\item $\textnormal{\texttt{setBits [] j \ensuremath{\tau}}} = \tau$,
  $\textnormal{\texttt{setBits (\allowbreak{}x :: ds) j \ensuremath{\tau}}} = \textnormal{\texttt{setBits ds (\allowbreak{}j / 2)}}\,(\tau[x\mapsto j\bmod2])$: the
  variables $ds$ set to the binary digits of $j$. \textnormal{\texttt{sum\_\allowbreak{}setBits}}:
  $\sum_{j<2^{|ds|}}G(\textnormal{\texttt{setBits ds j \ensuremath{\tau}}}) = \textnormal{\texttt{bsum ds G \ensuremath{\tau}}}$; \textnormal{\texttt{setBits\_\allowbreak{}not}},
  \textnormal{\texttt{setBits\_\allowbreak{}idx}} (variable $i$ of a repeat-free $ds$ gets digit $i$ of $j$, \textnormal{\texttt{digit 2 i j}}).
\item \textnormal{\texttt{Own := List (\allowbreak{}\ensuremath{\mathbb{N}} \ensuremath{\times} \ensuremath{\mathbb{N}} \ensuremath{\times} \ensuremath{\mathbb{N}})}}, an \term{owner table}: entries (digit variable, label, digit
  index). $\textnormal{\texttt{dbits own \ensuremath{\sigma} x}}$ is \textnormal{\texttt{digit 2 i (\allowbreak{}\ensuremath{\sigma} l)}} when \textnormal{\texttt{lookup x own = some (\allowbreak{}l,\allowbreak{} i)}},
  else $0$.
\item \textnormal{\texttt{ownOK own labs dig}} holds when the owner variables have no repeats, $labs$ has no
  repeats, every entry $(x,(l,i))$ has $l\in labs$ and $\mathit{dig}(l)[i] = x$, and for every
  $l\in labs$ the list $\mathit{dig}(l)$ has no repeats and each of its positions $i$ looks up to
  $(l,i)$.
\end{itemize}
\end{definition}

\noindent Lemmas: \textnormal{\texttt{mem\_\allowbreak{}of\_\allowbreak{}lookup}}, \textnormal{\texttt{ownOK\_\allowbreak{}lookup}}, \textnormal{\texttt{ownOK\_\allowbreak{}get}},
\textnormal{\texttt{ownOK\_\allowbreak{}nodup}}, \textnormal{\texttt{dbits\_\allowbreak{}get}} (digit $i$ of label $l$ reads \textnormal{\texttt{digit 2 i (\allowbreak{}\ensuremath{\sigma} l)}}),
\textnormal{\texttt{dbits\_\allowbreak{}update}} ($\textnormal{\texttt{dbits own \ensuremath{\sigma}[l\ensuremath{\mapsto}j]}} = \textnormal{\texttt{setBits (\allowbreak{}dig l) j (\allowbreak{}dbits own \ensuremath{\sigma})}}$),
\textnormal{\texttt{eq\_\allowbreak{}of\_\allowbreak{}digits}} (values below $2^k$ with equal first $k$ digits are equal),
\textnormal{\texttt{digit\_\allowbreak{}bit}}, \textnormal{\texttt{bits\_\allowbreak{}eq'}}.

\begin{definition}[Owner and dimension tables in one pass]\label{def:lean-pathsum-owntab}
A \term{digit table} $T$ is a list of pairs (label, its digit variables, low digit first).
\begin{itemize}
\item \textnormal{\texttt{ownRow l i xs}} is $[(x_0,(l,i)),(x_1,(l,i+1)),\dots]$; \textnormal{\texttt{ownRows T}} concatenates
  \textnormal{\texttt{ownRow l 0 ds}} over $(l,ds)\in T$.
\item \textnormal{\texttt{ownT own labs T}}: $labs$ is the list of keys of $T$, $own=\textnormal{\texttt{ownRows T}}$, and
  \textnormal{\texttt{nodupF}} holds of the owner variables and of $labs$.
  \lean!ownOK_tab!: \textnormal{\texttt{ownT own labs T}} gives \textnormal{\texttt{ownOK own labs (fun l => (T.lookup l).getD [])}}
  (from \lean!ownRow_keys!, \lean!mem_ownRow!, \lean!getElem_mem_ownRow!, \lean!lookup_of_mem!).
\item \textnormal{\texttt{dimsM D L T}} checks, in one merged pass, that the keys of $T$ are a subsequence of the
  keys of $L$ in order, each with $D(k)=2^{|ds|}$ for its $L$ value $k$ (\lean!dimsM_sound!).
  \textnormal{\texttt{dimsT E labs T}}: $labs$ is the keys of $T$, \textnormal{\texttt{nodupF}} of the keys of $T$ and of
  \textnormal{\texttt{E.labs}}, and \textnormal{\texttt{dimsM E.dimv E.labs T}}. \lean!dims_tab!: every $l\in labs$ has
  $\textnormal{\texttt{labDim E l}} = 2^{|(T.\mathrm{lookup}\,l).\mathrm{getD}\,[\,]|}$.
\end{itemize}
\end{definition}

\begin{theorem}[\textnormal{\texttt{sumOver\_\allowbreak{}dbits}}]
For \textnormal{\texttt{ownOK own labs dig}}, labels $L\subseteq labs$ of dimension $2^{|\mathit{dig}(l)|}$, and
every $\sigma$: $\textnormal{\texttt{sumOver D L (\allowbreak{}fun \ensuremath{\sigma}' => g (\allowbreak{}dbits own \ensuremath{\sigma}')) \ensuremath{\sigma}}} = \textnormal{\texttt{bsum (\allowbreak{}L.\allowbreak{}flatMap dig) g (\allowbreak{}dbits own \ensuremath{\sigma})}}$.
\end{theorem}

\begin{definition}[Fragments over digits]\label{def:lean-pathsum-fragD}
\begin{itemize}
\item $\textnormal{\texttt{eqCons xs ys}} = [[[x],[y]] : (x,y)\in\mathrm{zip}(xs,ys)]$, the constraints $x+y=0$.
  \textnormal{\texttt{mem\_\allowbreak{}eqCons\_\allowbreak{}iff}}: for equal lengths, all hold iff $\tau(xs_i) = \tau(ys_i)$ for every $i$.
\item \textnormal{\texttt{lensOK dig k ls}}: every $\mathit{dig}(l)$, $l\in ls$, has length $k$.
\item $\textnormal{\texttt{anchorD dig as}} = [\mathit{dig}(as_j)[j] : j<|as|]$ (\textnormal{\texttt{getD j 0}}): the anchor
  digits, digit $j$ of label $as_j$.
\item \textnormal{\texttt{tableOK M ds P ts}}: for every $v<2^{|ds|}$,
  $\textnormal{\texttt{pn (\allowbreak{}setBits ds v 0) P}} - ts_v \equiv 0\pmod M$ ($ts_v$ = \textnormal{\texttt{ts.\allowbreak{}getD v 0}}).
\item $\textnormal{\texttt{zFragD dig as ls P}} = \langle[\,],\ P,\ \bigcup_{l\in ls}\textnormal{\texttt{eqCons (\allowbreak{}anchorD dig as) (\allowbreak{}dig l)}},\ 1\rangle$.
\item $\textnormal{\texttt{bFrag dig k a b}} = \langle[\,],[\,],\textnormal{\texttt{eqCons (\allowbreak{}dig k) (\allowbreak{}dig b ++ dig a)}},1\rangle$: the
  digits of $k$ equal the digits of the low part $b$ then those of the high part $a$.
\item $\textnormal{\texttt{prodTerms M K xs ys}} = [([xs_a, ys_b],\ 2^{a+b}(M/K)) : a<|xs|,\ b<|xs|-a]$, the terms
  $2^{a+b}x_ay_b/K$ with $a+b<k$; \textnormal{\texttt{pn\_\allowbreak{}prodTerms}}, \textnormal{\texttt{pn\_\allowbreak{}flatMap}}, \textnormal{\texttt{list\_\allowbreak{}sum\_\allowbreak{}range}},
  \textnormal{\texttt{val\_\allowbreak{}digits}} ($v = \sum_{a<k}2^a\,\textnormal{\texttt{digit 2 a v}}$ for $v<2^k$), \textnormal{\texttt{zq\_\allowbreak{}digit}}.
\item $\textnormal{\texttt{fFragD M K k dig o i}} = \langle[\,],\ \textnormal{\texttt{prodTerms M K (\allowbreak{}dig o) (\allowbreak{}dig i)}},\ [\,],\ (1/\sqrt2)^k\rangle$.
\item \textnormal{\texttt{xFragD M K k dig P ys os is}} has path $ys$ (fresh digit variables), phase
  $P\mathbin{+\!\!+}\bigcup_{o\in os}\textnormal{\texttt{prodTerms M K (\allowbreak{}dig o) ys}}\mathbin{+\!\!+}\bigcup_{i\in is}\textnormal{\texttt{negP M (\allowbreak{}prodTerms M K (\allowbreak{}dig i) ys)}}$,
  no constraints, and constant $(1/\sqrt2)^{k(|os|+|is|)}$.
\end{itemize}
\end{definition}

\begin{theorem}[\textnormal{\texttt{ex\_\allowbreak{}prod\_\allowbreak{}digits}}]
For $K=2^k$, $M=Kq$, $0<M$, $|xs|=|ys|=k$, $o,j<2^k$, and $\tau$ setting $xs_a$ to digit $a$ of $o$ and
$ys_b$ to digit $b$ of $j$:
$\operatorname{ex}(\textnormal{\texttt{pn \ensuremath{\tau} (\allowbreak{}prodTerms M K xs ys)}}/M) = \operatorname{ex}(oj/K)$.
\end{theorem}

\begin{theorem}[Generator fragments over digits]
Assume \textnormal{\texttt{ownOK own labs dig}}, $\sigma(l)<2^{|\mathit{dig}(l)|}$ for $l\in labs$, every leg in
$labs$ with $k$ digits (\textnormal{\texttt{lensOK}}), and the node's dimension value $2^k$.
\begin{itemize}
\item \textnormal{\texttt{z\_\allowbreak{}frag\_\allowbreak{}dig}} (Z): with anchors $as$ ($|as|=k$, each a leg), a leg $l_0$, $0<M$, phase
  entries $\varphi_v = ts_v/M$, \textnormal{\texttt{tableOK M (\allowbreak{}anchorD dig as) P ts}}, repeat-free anchor digits,
  and $P$ reading only anchor digits:
  $\textnormal{\texttt{val E \ensuremath{\sigma} n}} = \textnormal{\texttt{fval M (\allowbreak{}zFragD dig as (\allowbreak{}outs ++ ins) P) (\allowbreak{}dbits own \ensuremath{\sigma})}}$
  (\textnormal{\texttt{tZ\_\allowbreak{}val}}, \textnormal{\texttt{pn\_\allowbreak{}congr}}, \textnormal{\texttt{fz\_\allowbreak{}pair}}).
\item \textnormal{\texttt{b\_\allowbreak{}frag}} (B, \textnormal{\texttt{outs = [k]}}, \textnormal{\texttt{ins = [a,\allowbreak{} b]}}) and \textnormal{\texttt{s\_\allowbreak{}frag}} (S,
  \textnormal{\texttt{outs = [a,\allowbreak{} b]}}, \textnormal{\texttt{ins = [k]}}): with dimension values $2^p$ (\textnormal{\texttt{dim}}) and $2^q$
  (\textnormal{\texttt{ph}}), $|\mathit{dig}(k)| = q+p$, $|\mathit{dig}(a)| = p$, $|\mathit{dig}(b)| = q$:
  $\textnormal{\texttt{val E \ensuremath{\sigma} n}} = \textnormal{\texttt{fval M (\allowbreak{}bFrag dig k a b) (\allowbreak{}dbits own \ensuremath{\sigma})}}$. The core is \textnormal{\texttt{b\_\allowbreak{}cond}}:
  the constraints hold iff $\sigma(k) = \sigma(a)2^q+\sigma(b)$ (\textnormal{\texttt{digit\_\allowbreak{}hi}}).
\item \textnormal{\texttt{f\_\allowbreak{}frag\_\allowbreak{}dig}} (F, \textnormal{\texttt{outs = [o]}}, \textnormal{\texttt{ins = [i]}}, $K=2^k$, $M=Kq$, $0<M$):
  $\textnormal{\texttt{val E \ensuremath{\sigma} n}} = \textnormal{\texttt{fval M (\allowbreak{}fFragD M K k dig o i) (\allowbreak{}dbits own \ensuremath{\sigma})}}$.
\item \textnormal{\texttt{x\_\allowbreak{}frag\_\allowbreak{}dig}} (X): with fresh repeat-free $ys$ ($|ys|=k$, none in the owner table),
  $K=2^k$, $M=Kq$, \textnormal{\texttt{tableOK M ys P ts}}, and $P$ reading only $ys$:
  $\textnormal{\texttt{val E \ensuremath{\sigma} n}} = \textnormal{\texttt{fval M (\allowbreak{}xFragD M K k dig P ys outs ins) (\allowbreak{}dbits own \ensuremath{\sigma})}}$. Each term of
  the X sum is a summand of the fragment at \textnormal{\texttt{setBits ys j}} (\textnormal{\texttt{x\_\allowbreak{}term\_\allowbreak{}dig}}; \textnormal{\texttt{ex\_\allowbreak{}conj}},
  \textnormal{\texttt{ex\_\allowbreak{}neg\_\allowbreak{}of}}).
\end{itemize}
\end{theorem}

\begin{definition}[Network path sums over digits]\label{def:lean-pathsum-netpsd}
$\textnormal{\texttt{netDataD dig N fs}} = \langle\bigcup_{l\in N.\text{int}}\mathit{dig}(l)\mathbin{+\!\!+}(\textnormal{\texttt{joinAll fs}}).\text{path},\ \text{phase},\ \text{cons},\ [\,]\rangle$.
\textnormal{\texttt{structure NetPSD E N M own labs dig ps : Prop}} has fields \textnormal{\texttt{ownok}}, \textnormal{\texttt{nodes}},
\textnormal{\texttt{node}} (each entry is its fragment's value at \textnormal{\texttt{dbits own \ensuremath{\sigma}}} for in-range $\sigma$ on
$labs$), \textnormal{\texttt{disj}}, \textnormal{\texttt{dims}} (every label of $labs$ has dimension $2^{|\mathit{dig}(l)|}$),
\textnormal{\texttt{cover}}, and \textnormal{\texttt{sub}} (summed and boundary labels lie in $labs$).
\end{definition}

\begin{theorem}[Network theorems over digits]
\begin{itemize}
\item \textnormal{\texttt{NetPSD.\allowbreak{}eval\_\allowbreak{}eq}}: $\textnormal{\texttt{eval E N \ensuremath{\sigma}}} = \textnormal{\texttt{val M (\allowbreak{}joinAll fs).\allowbreak{}c (\allowbreak{}netDataD dig N fs) (\allowbreak{}dbits own (\allowbreak{}clamp N.\allowbreak{}bnd \ensuremath{\sigma}))}}$
  for $\sigma$ in range on the boundary.
\item \textnormal{\texttt{equal\_\allowbreak{}psd}}: one owner table for both sides, equal boundaries, and agreement of the
  scaled digit path sums at every $\tau$ give \textnormal{\texttt{Equal E s\ensuremath{_1} N\ensuremath{_1} s\ensuremath{_2} N\ensuremath{_2}}}.
\item \textnormal{\texttt{equal\_\allowbreak{}psd2}}: one owner table, label list, and digit map per side; the hypothesis is
  agreement at $\textnormal{\texttt{dbits own\ensuremath{_1} (\allowbreak{}clamp N\ensuremath{_1}.\allowbreak{}bnd \ensuremath{\sigma})}}$ and $\textnormal{\texttt{dbits own\ensuremath{_2} (\allowbreak{}clamp N\ensuremath{_2}.\allowbreak{}bnd \ensuremath{\sigma})}}$ for
  every $\sigma$.
\item \textnormal{\texttt{varsIn vs s}}: every variable of the phase, constraints, and rows lies in $vs$.
  \textnormal{\texttt{val\_\allowbreak{}agree}}: a path-free sum with \textnormal{\texttt{varsIn vs s}} has the same value at two assignments
  that agree on $vs$ (\textnormal{\texttt{mz\_\allowbreak{}agree}}, \textnormal{\texttt{fz\_\allowbreak{}agree}}). \textnormal{\texttt{dbits\_\allowbreak{}agree}}: if the two owner
  tables look up every $x\in vs$ alike, the two digit assignments agree on $vs$.
\end{itemize}
\end{theorem}

\section{The W node: \texorpdfstring{\file{PathSumW}}{PathSumW}}\label{sec:lean-pathsum-w}

\file{PathSumW} (imports \file{PathSumDigits}) gives the qubit W node as a fragment with no path
variable. $\textnormal{\texttt{pairsOf [a\_\allowbreak{}1,\allowbreak{}\textbackslash{}dots,\allowbreak{}a\_\allowbreak{}n]}} = [(a_i,a_{i'}) : i<i']$ (\textnormal{\texttt{pairsOf\_\allowbreak{}map}}).
\textnormal{\texttt{wFrag dig os i}} has no path variable, no phase, constant $1$, and the constraints
\[
\bigl[[x] : x\in\mathit{dig}(i)\mathbin{+\!\!+}\textstyle\bigcup_o\mathit{dig}(o)\bigr]\ ::\ \bigl[[[x,x']] : (x,x')\in\textnormal{\texttt{pairsOf}}(\textstyle\bigcup_o\mathit{dig}(o))\bigr]:
\]
the input digit plus all output digits is $0$, and the product of any two output digits is $0$.
$\textnormal{\texttt{ones vs}}$ counts the nonzero entries (\textnormal{\texttt{ones\_\allowbreak{}cons}}, \textnormal{\texttt{sum\_\allowbreak{}ones}}, \textnormal{\texttt{pairs\_\allowbreak{}ones}}:
the pairwise products vanish iff at most one entry is $1$; \textnormal{\texttt{cond0}}, \textnormal{\texttt{cond1}}).
\begin{theorem}[\textnormal{\texttt{tW\_\allowbreak{}bits}}, \textnormal{\texttt{wnode\_\allowbreak{}frag}}]
For values below $2$, $\textnormal{\texttt{tW 2 vs i}} = 1$ exactly when $i+\sum vs = 0$ in $\Zd2$ and all pairwise
products vanish, else $0$. \textnormal{\texttt{wnode\_\allowbreak{}frag}}: for a W node of dimension value $2$ with
\textnormal{\texttt{ins = [i]}} and every leg a one-digit label of $labs$,
$\textnormal{\texttt{val E \ensuremath{\sigma} n}} = \textnormal{\texttt{fval M (\allowbreak{}wFrag dig outs i) (\allowbreak{}dbits own \ensuremath{\sigma})}}$ (\textnormal{\texttt{dbits\_\allowbreak{}flat}}).
\end{theorem}

\section{Legless and weighted spiders: \texorpdfstring{\file{PathSumNetW}}{PathSumNetW}}\label{sec:lean-pathsum-netw}

\file{PathSumNetW} (imports \file{PathSumW}, \file{PathSumChain}) defines
\textnormal{\texttt{Node.\allowbreak{}withKind n k := \{ n with kind := k \}}} and \textnormal{\texttt{Frag.\allowbreak{}scale C f := \{ f with c := C * f.\allowbreak{}c \}}}
(\textnormal{\texttt{fval\_\allowbreak{}scale}}: the value scales by $C$).
\begin{itemize}
\item \textnormal{\texttt{val\_\allowbreak{}z\_\allowbreak{}legless}}: a Z node with no legs has the entry of the same node as X, so
  \textnormal{\texttt{x\_\allowbreak{}frag\_\allowbreak{}dig}} gives its fragment.
\item \textnormal{\texttt{val\_\allowbreak{}zw}} and \textnormal{\texttt{val\_\allowbreak{}xw}}: if $\textnormal{\texttt{coef ph v}} = C\operatorname{ex}(\textnormal{\texttt{phase ph v}})$
  for every $v$, a ZW (XW) node's entry is $C$ times the entry of the node as Z (X).
\item \textnormal{\texttt{w\_\allowbreak{}frag}}: if $\textnormal{\texttt{val E \ensuremath{\sigma} n}} = C\,\textnormal{\texttt{val E \ensuremath{\sigma} n'}}$ and $\textnormal{\texttt{val E \ensuremath{\sigma} n'}} = \textnormal{\texttt{fval M f \ensuremath{\tau}}}$,
  then $\textnormal{\texttt{val E \ensuremath{\sigma} n}} = \textnormal{\texttt{fval M (\allowbreak{}f.\allowbreak{}scale C) \ensuremath{\tau}}}$.
\item \textnormal{\texttt{zeroOK s}} holds when some constraint is \textnormal{\texttt{[[]]}} (the constant $1$, so the
  constraint $1=0$ fails everywhere); \textnormal{\texttt{val\_\allowbreak{}of\_\allowbreak{}zeroOK}}: then $\textnormal{\texttt{val M c s \ensuremath{\sigma}}} = 0$
  (\textnormal{\texttt{bsum\_\allowbreak{}zero}}).
\end{itemize}

\section{Networks with atoms: \texorpdfstring{\file{PathSumNetAmp}}{PathSumNetAmp}}\label{sec:lean-pathsum-netamp}

\file{PathSumNetAmp} (imports \file{PathSumNetW}, \file{PathSumAmp}) gives weighted spiders whose
weights are not phases as fragments with amplitude atoms.

\begin{definition}[Atom tables]\label{def:lean-pathsum-atab}
$\textnormal{\texttt{atomsAt as \ensuremath{\tau}}}$ lists the indices of the atoms whose polynomial is $1$ at $\tau$
(\textnormal{\texttt{aval\_\allowbreak{}atomsAt}}: \textnormal{\texttt{aval A as \ensuremath{\tau}}} is the product of $A$ over them; \textnormal{\texttt{aval\_\allowbreak{}agree}}).
\textnormal{\texttt{atabOK ds as tbl}} holds when the atoms read only $ds$ and, for every $v<2^{|ds|}$,
$\textnormal{\texttt{atomsAt as (\allowbreak{}setBits ds v 0)}} = tbl_v$ (\textnormal{\texttt{tbl.\allowbreak{}getD v []}}).
\end{definition}

\begin{theorem}[Weighted Z spiders]
\textnormal{\texttt{val\_\allowbreak{}zw\_\allowbreak{}at}}: if $\textnormal{\texttt{coef ph v}} = C\operatorname{ex}(\textnormal{\texttt{phase ph v}})\,W(v)$ and $l_0$ is a
leg, the ZW entry is $C\,W(\sigma(l_0))$ times the Z entry. \textnormal{\texttt{za\_\allowbreak{}frag}}: under the hypotheses of
\textnormal{\texttt{z\_\allowbreak{}frag\_\allowbreak{}dig}} for the anchors, $\textnormal{\texttt{coef ph v}} = C\operatorname{ex}(\textnormal{\texttt{phase ph v}})\prod_{i\in tbl_v}A_i$,
\textnormal{\texttt{atabOK (\allowbreak{}anchorD dig as) ats tbl}}, and a fragment $f$ of the node as Z,
$\textnormal{\texttt{val E \ensuremath{\sigma} n}} = \textnormal{\texttt{fval M (\allowbreak{}f.\allowbreak{}scale C) (\allowbreak{}dbits own \ensuremath{\sigma})}} \cdot \textnormal{\texttt{aval A ats (\allowbreak{}dbits own \ensuremath{\sigma})}}$.
\textnormal{\texttt{atoms\_\allowbreak{}nil}}: an entry equal to a fragment value is that value times no atoms.
\end{theorem}

\begin{definition}[Fragments with atoms]
$\textnormal{\texttt{fvalA M A f as \ensuremath{\tau}}} = \textnormal{\texttt{bsum f.\allowbreak{}path}}\,(\rho\mapsto\textnormal{\texttt{aval A as \ensuremath{\rho}}}\cdot\textnormal{\texttt{summand M f.\allowbreak{}c \ensuremath{\langle}f.\allowbreak{}path,\allowbreak{} f.\allowbreak{}phase,\allowbreak{} f.\allowbreak{}cons,\allowbreak{} []\ensuremath{\rangle} \ensuremath{\rho}}})\,\tau$.
\textnormal{\texttt{DisjA (\allowbreak{}f,\allowbreak{} as) (\allowbreak{}g,\allowbreak{} bs)}} is \textnormal{\texttt{Disj f g}} plus: no path variable of $f$ is read by $bs$.
\end{definition}
\noindent \textnormal{\texttt{fvalA\_\allowbreak{}of}}: atoms reading no path variable of $f$ factor out
($x = \textnormal{\texttt{fval M f \ensuremath{\tau}}}\cdot\textnormal{\texttt{aval A as \ensuremath{\tau}}}\Rightarrow x = \textnormal{\texttt{fvalA M A f as \ensuremath{\tau}}}$).
\textnormal{\texttt{summand\_\allowbreak{}scale}}, \textnormal{\texttt{fvalA\_\allowbreak{}join}} (product for \textnormal{\texttt{DisjA}} both ways), \textnormal{\texttt{fvalA\_\allowbreak{}unit}},
\textnormal{\texttt{mem\_\allowbreak{}avars\_\allowbreak{}flatten}}, \textnormal{\texttt{mem\_\allowbreak{}joinAll\_\allowbreak{}path}}, and \textnormal{\texttt{fvalA\_\allowbreak{}joinAll}}:
$\textnormal{\texttt{fvalA M A (\allowbreak{}joinAll fs) (\allowbreak{}concatenated atoms) \ensuremath{\tau}}} = \prod\textnormal{\texttt{fvalA M A f as \ensuremath{\tau}}}$ for pairwise
\textnormal{\texttt{DisjA}} pairs.

\begin{theorem}[Weighted X spiders]
\textnormal{\texttt{val\_\allowbreak{}zw\_\allowbreak{}legless}}: a legless ZW node has the XW entry. \textnormal{\texttt{setBits\_\allowbreak{}agree}}.
\textnormal{\texttt{xa\_\allowbreak{}frag}}: under the hypotheses of \textnormal{\texttt{x\_\allowbreak{}frag\_\allowbreak{}dig}}, with
$\textnormal{\texttt{coef ph v}} = C\operatorname{ex}(\textnormal{\texttt{phase ph v}})\prod_{i\in tbl_v}A_i$ and
\textnormal{\texttt{atabOK ys ats tbl}},
$\textnormal{\texttt{val E \ensuremath{\sigma} n}} = \textnormal{\texttt{fvalA M A (\allowbreak{}(\allowbreak{}xFragD M K k dig P ys outs ins).\allowbreak{}scale C) ats (\allowbreak{}dbits own \ensuremath{\sigma})}}$.
\end{theorem}

\begin{definition}[Networks with atoms]\label{def:lean-pathsum-netpsda}
$\textnormal{\texttt{netDataA dig N fs ats}} = \bigl\langle\langle[\,], \text{phase}, \text{cons}, [\,]\rangle,\ \bigcup_{l\in N.\text{int}}\mathit{dig}(l)\mathbin{+\!\!+}(\textnormal{\texttt{joinAll fs}}).\text{path},\ \text{the concatenated atoms}\bigr\rangle$:
every summed digit and every fragment path variable is held, the inner path is empty.
\textnormal{\texttt{structure NetPSDA E N M A own labs dig ps}} is \textnormal{\texttt{NetPSD}} for triples (node, fragment,
atoms), with \textnormal{\texttt{node}} stating $\textnormal{\texttt{val E \ensuremath{\sigma} n}} = \textnormal{\texttt{fvalA M A f as (\allowbreak{}dbits own \ensuremath{\sigma})}}$ and
\textnormal{\texttt{disj}} using \textnormal{\texttt{DisjA}}.
\end{definition}

\begin{theorem}[\textnormal{\texttt{NetPSDA.\allowbreak{}eval\_\allowbreak{}eq}}, \textnormal{\texttt{equal\_\allowbreak{}psda2}}]
$\textnormal{\texttt{eval E N \ensuremath{\sigma}}} = \textnormal{\texttt{valA M (\allowbreak{}joinAll fs).\allowbreak{}c A (\allowbreak{}netDataA dig N fs ats) (\allowbreak{}dbits own (\allowbreak{}clamp N.\allowbreak{}bnd \ensuremath{\sigma}))}}$.
\textnormal{\texttt{equal\_\allowbreak{}psda2}} is \textnormal{\texttt{equal\_\allowbreak{}psd2}} with \textnormal{\texttt{valA}}. \textnormal{\texttt{valA\_\allowbreak{}smul}} scales the constant;
\textnormal{\texttt{valA\_\allowbreak{}agree}}: a sum with no held and no path variable, \textnormal{\texttt{varsIn vs t.\allowbreak{}s}}, and atoms reading
only $vs$ has equal values at assignments agreeing on $vs$.
\end{theorem}

\section{Qudit path sums: \texorpdfstring{\file{PathSumQudit}}{PathSumQudit}}\label{sec:lean-pathsum-qudit}

\file{PathSumQudit} (imports \file{PathSum} and Mathlib's quadratic characters) defines path sums
over $\Zd p$ for a prime $p$ in namespace \textnormal{\texttt{QSem}}: variables take values in $\Zd p$, the phase is
a polynomial over $\mathbb F_p$ worth $\psi(\cdot)$, and the steps are \textnormal{\texttt{perm}}, \textnormal{\texttt{norm}},
\textnormal{\texttt{elim}}, \textnormal{\texttt{hhCons}}, \textnormal{\texttt{hhPivot}}, \textnormal{\texttt{gauss}}, \textnormal{\texttt{canon}}, and \textnormal{\texttt{pin}}. Theorems that
need a field carry \textnormal{\texttt{[Fact p.\allowbreak{}Prime]}}; the sums carry \textnormal{\texttt{[NeZero p]}}.

\subsection{Polynomials over \texorpdfstring{$\mathbb F_p$}{F\_p}}

\begin{definition}[Qudit polynomials]\label{def:lean-pathsum-qpoly}
\textnormal{\texttt{QSem.\allowbreak{}Mono := List (\allowbreak{}\ensuremath{\mathbb{N}} \ensuremath{\times} \ensuremath{\mathbb{N}})}} (variable, exponent) pairs; \textnormal{\texttt{QSem.\allowbreak{}Poly := List (\allowbreak{}Mono \ensuremath{\times} \ensuremath{\mathbb{N}})}}
with coefficients read mod $p$. For $\sigma:\N\to\Zd p$:
$\textnormal{\texttt{qm \ensuremath{\sigma} m}} = \prod_{(x,e)\in m}\sigma(x)^e$ and $\textnormal{\texttt{qv \ensuremath{\sigma} P}} = \sum_{(m,k)\in P} k\,\textnormal{\texttt{qm \ensuremath{\sigma} m}}\in\Zd p$.
$\textnormal{\texttt{yexp y m}} = \sum_{(x,e)\in m,\ x=y} e$ is the exponent of $y$; \textnormal{\texttt{strip y m}} removes $y$;
\textnormal{\texttt{qm\_\allowbreak{}update}}: $\textnormal{\texttt{qm \ensuremath{\sigma}[y\ensuremath{\mapsto}b] m}} = b^{\textnormal{\texttt{yexp y m}}}\,\textnormal{\texttt{qm \ensuremath{\sigma} (\allowbreak{}strip y m)}}$.
\textnormal{\texttt{pvars P}}, \textnormal{\texttt{cvars C}} list variables; \textnormal{\texttt{qm\_\allowbreak{}update\_\allowbreak{}not}}, \textnormal{\texttt{qv\_\allowbreak{}update\_\allowbreak{}not}},
\textnormal{\texttt{holds\_\allowbreak{}update\_\allowbreak{}not}} give invariance off the variables. \textnormal{\texttt{qm\_\allowbreak{}nil}}, \textnormal{\texttt{qm\_\allowbreak{}cons}},
\textnormal{\texttt{qm\_\allowbreak{}append}}, \textnormal{\texttt{qv\_\allowbreak{}nil}}, \textnormal{\texttt{qv\_\allowbreak{}cons}}, \textnormal{\texttt{qv\_\allowbreak{}append}}, \textnormal{\texttt{qv\_\allowbreak{}perm}}.
\end{definition}

\begin{definition}[Splitting and operations]
\begin{itemize}
\item $\textnormal{\texttt{yk y k P}} = [(\textnormal{\texttt{strip y m}}, c) : (m,c)\in P,\ \textnormal{\texttt{yexp y m}} = k]$, the coefficient
  of $y^k$; \textnormal{\texttt{qv\_\allowbreak{}split}}: when every exponent of $y$ is at most $2$,
  $\textnormal{\texttt{qv \ensuremath{\sigma}[y\ensuremath{\mapsto}b] P}} = \textnormal{\texttt{qv \ensuremath{\sigma} (\allowbreak{}yk y 0 P)}} + b\,\textnormal{\texttt{qv \ensuremath{\sigma} (\allowbreak{}yk y 1 P)}} + b^2\,\textnormal{\texttt{qv \ensuremath{\sigma} (\allowbreak{}yk y 2 P)}}$;
  \textnormal{\texttt{qv\_\allowbreak{}split1}} is the linear case; \textnormal{\texttt{yk\_\allowbreak{}two\_\allowbreak{}nil}}, \textnormal{\texttt{not\_\allowbreak{}mem\_\allowbreak{}yk}}, \textnormal{\texttt{qv\_\allowbreak{}yk\_\allowbreak{}update}}.
\item \textnormal{\texttt{csum P}} is the sum of the coefficients; \textnormal{\texttt{qv\_\allowbreak{}const}}: for constant terms,
  $\textnormal{\texttt{qv \ensuremath{\sigma} P}} = \textnormal{\texttt{csum P}} \bmod p$.
\item $\textnormal{\texttt{scaleQ k P}}$ multiplies every coefficient by $k$; $\textnormal{\texttt{mulQ A B}}$ is the product
  ($[(s_1\mathbin{+\!\!+}t_1, s_2t_2)]$); $\textnormal{\texttt{powQ B 0}} = [([\,],1)]$,
  $\textnormal{\texttt{powQ B (\allowbreak{}n+1)}} = \textnormal{\texttt{mulQ (\allowbreak{}powQ B n) B}}$;
  $\textnormal{\texttt{substQ z B P}} = \bigcup_{(m,c)}\textnormal{\texttt{mulQ [(\allowbreak{}strip z m,\allowbreak{} c)] (\allowbreak{}powQ B (\allowbreak{}yexp z m))}}$.
  \textnormal{\texttt{qv\_\allowbreak{}scaleQ}}, \textnormal{\texttt{qv\_\allowbreak{}mapmul}}, \textnormal{\texttt{qv\_\allowbreak{}mulQ}}, \textnormal{\texttt{qv\_\allowbreak{}powQ}}, and
  \textnormal{\texttt{qv\_\allowbreak{}substQ}} ($\textnormal{\texttt{qv \ensuremath{\sigma} (\allowbreak{}substQ z B P)}} = \textnormal{\texttt{qv \ensuremath{\sigma}[z\ensuremath{\mapsto}qv \ensuremath{\sigma} B] P}}$) give the values.
\end{itemize}
\end{definition}

\begin{definition}[Qudit normal form]\label{def:lean-pathsum-qnorm}
\textnormal{\texttt{mergeE t m}} merges $t=(x,e)$ into the head of $m$ when the head has variable $x$ (exponents
added), else conses. $\textnormal{\texttt{redE p e}} = ((e-1)\bmod(p-1))+1$ reduces an exponent by $x^p=x$
(\textnormal{\texttt{pow\_\allowbreak{}shift}}, \textnormal{\texttt{pow\_\allowbreak{}redE}}: $x^{\textnormal{\texttt{redE p e}}} = x^e$ for $e\ne0$, $p$ prime).
\[
\textnormal{\texttt{mnormQ p m}} = [(x,\textnormal{\texttt{redE p e}}) : (x,e)\in\textnormal{\texttt{foldr mergeE []}}(\textnormal{\texttt{insertionSort}}_{x\le x'}\,m),\ e\ne0].
\]
\textnormal{\texttt{insQ p t P}} adds $t$ to the first term of $P$ with the same monomial mod $p$, dropping zero
terms (\textnormal{\texttt{cast\_\allowbreak{}mod\_\allowbreak{}zero}}, \textnormal{\texttt{qv\_\allowbreak{}insQ}}). \textnormal{\texttt{mleQ}} is the lexicographic order on monomials
comparing variables, then exponents. \textnormal{\texttt{normQ p P}} inserts every term with its monomial
\textnormal{\texttt{mnormQ}}-normalized by \textnormal{\texttt{insQ}} (a right fold) and sorts the result by
\textnormal{\texttt{insertionSort}} on \textnormal{\texttt{mleQ}}. \textnormal{\texttt{qm\_\allowbreak{}mnormQ}}, \textnormal{\texttt{qm\_\allowbreak{}foldr\_\allowbreak{}mergeE}},
\textnormal{\texttt{qm\_\allowbreak{}filter\_\allowbreak{}ne}}, and \textnormal{\texttt{qv\_\allowbreak{}normQ}} ($p$ prime) state value preservation.
\end{definition}

\subsection{Roots of unity indexed by \texorpdfstring{$\Zd p$}{Z\_p}}

\begin{definition}[Characters and Gauss sums]\label{def:lean-pathsum-psi}
$\psi_p(x) = \omega_p^{x.\mathrm{val}}$ for $x\in\Zd p$ (\textnormal{\texttt{\ensuremath{\psi} p x}}). \textnormal{\texttt{\ensuremath{\psi}\_\allowbreak{}add}}, \textnormal{\texttt{\ensuremath{\psi}\_\allowbreak{}zero}},
\textnormal{\texttt{\ensuremath{\psi}\_\allowbreak{}nat}} ($\psi_p(k) = \omega_p^k$), \textnormal{\texttt{zsum\_\allowbreak{}range}}, \textnormal{\texttt{sum\_\allowbreak{}\ensuremath{\psi}}} ($\sum_x\psi_p(x)=0$ for
$p$ prime), and \textnormal{\texttt{char\_\allowbreak{}sum}}: $\sum_{b\in\Zd p}\psi_p(bQ) = p\,[Q=0]$.
The \term{quadratic Gauss sum} is $G_p(a) = \textnormal{\texttt{qg p a}} = \sum_{y\in\Zd p}\psi_p(ay^2)$
\cite{berndt1998gauss}. The Legendre symbol by Euler's criterion is
\[
\textnormal{\texttt{lg p a}} = \begin{cases}0 & a\equiv0\\ 1 & a^{\lfloor p/2\rfloor}\equiv1\\ -1 & \text{else}\end{cases}\pmod p.
\]
\end{definition}

\begin{theorem}[Legendre form of the Gauss sum]
\textnormal{\texttt{lg\_\allowbreak{}eq}} ($p\ne2$ prime): \textnormal{\texttt{lg p a}} equals Mathlib's \textnormal{\texttt{quadraticChar (\allowbreak{}ZMod p) a}}.
\textnormal{\texttt{qg\_\allowbreak{}char}}: for $a\ne0$, $\sum_y\psi_p(ay^2) = \chi(a)\sum_x\chi(x)\psi_p(x)$, $\chi$ the
quadratic character. \textnormal{\texttt{qg\_\allowbreak{}legendre}}: for $p\ne2$ prime and $a\not\equiv0$,
\[
G_p(a) = \Bigl(\frac ap\Bigr)\sum_{k<p}\Bigl(\frac kp\Bigr)\omega_p^k .
\]
\end{theorem}

\begin{definition}[Sums over $\Zd p$]
$\textnormal{\texttt{qsum [] f \ensuremath{\sigma}}} = f(\sigma)$, $\textnormal{\texttt{qsum (\allowbreak{}y :: ys) f \ensuremath{\sigma}}} = \sum_{b\in\Zd p}\textnormal{\texttt{qsum ys f \ensuremath{\sigma}[y\ensuremath{\mapsto}b]}}$.
Lemmas: \textnormal{\texttt{qsum\_\allowbreak{}congr}}, \textnormal{\texttt{qsum\_\allowbreak{}fsum}} (a finite sum commutes), \textnormal{\texttt{qsum\_\allowbreak{}out}},
\textnormal{\texttt{qsum\_\allowbreak{}perm}}, \textnormal{\texttt{qsum\_\allowbreak{}head}}; in \file{PathSumQuditNet} also \textnormal{\texttt{qsum\_\allowbreak{}append}},
\textnormal{\texttt{qsum\_\allowbreak{}mul\_\allowbreak{}left}}, \textnormal{\texttt{qsum\_\allowbreak{}inv}}.
\end{definition}

\begin{definition}[Qudit path sum]\label{def:lean-pathsum-qdata}
\textnormal{\texttt{QSem.\allowbreak{}Data}} has fields \textnormal{\texttt{path}}, \textnormal{\texttt{phase : Poly}}, \textnormal{\texttt{cons : List Poly}},
\textnormal{\texttt{rows : List Poly}} (derives \textnormal{\texttt{DecidableEq}}).
\begin{align*}
\textnormal{\texttt{summand p c s \ensuremath{\sigma}}} &= \Bigl[\textstyle\bigwedge_{q\in\text{cons}\mathbin{+\!\!+}\text{rows}}\textnormal{\texttt{qv \ensuremath{\sigma} q}} = 0\Bigr]\,c\,\psi_p(\textnormal{\texttt{qv \ensuremath{\sigma} s.\allowbreak{}phase}}),\\
\textnormal{\texttt{val p c s \ensuremath{\sigma}}} &= \textnormal{\texttt{qsum s.\allowbreak{}path (\allowbreak{}summand p c s) \ensuremath{\sigma}}}.
\end{align*}
\textnormal{\texttt{summand\_\allowbreak{}update\_\allowbreak{}not}} gives invariance off the variables.
\end{definition}

\subsection{Qudit steps}

\paragraph{Permutation and normalization.} \textnormal{\texttt{permS l}} with \textnormal{\texttt{decide (\allowbreak{}s.\allowbreak{}path.\allowbreak{}Perm l)}}
(\textnormal{\texttt{perm\_\allowbreak{}sound}}); \textnormal{\texttt{normS p}} applies \textnormal{\texttt{normQ p}} to the phase, constraints, and rows
(\textnormal{\texttt{norm\_\allowbreak{}sound}}, $p$ prime; \textnormal{\texttt{holds\_\allowbreak{}map}}).

\paragraph{[Elim].} \textnormal{\texttt{elimOK s}}: path $=y::ys$, $y\notin ys$, $y$ not in the phase, a
constraint, or a row; \textnormal{\texttt{elimS}} drops it. \textnormal{\texttt{elim\_\allowbreak{}sound}}: factor $p$.

\paragraph{[HH] on a constraint.} \textnormal{\texttt{hhConsOK s}}: path $=y::ys$, $y\notin ys$, every term has
exponent of $y$ at most $1$, $y$ in no constraint or row.
$\textnormal{\texttt{hhConsS s}} = \{ys,\ \textnormal{\texttt{yk y 0 phase}},\ \textnormal{\texttt{yk y 1 phase}} :: \text{cons},\ \text{rows}\}$.
\textnormal{\texttt{hhCons\_\allowbreak{}sound}} ($p$ prime): factor $p$ (\textnormal{\texttt{char\_\allowbreak{}sum}}).

\paragraph{Pivots.} $\textnormal{\texttt{inv p a}} = a^{p-2}\bmod p$. \textnormal{\texttt{pivOK p v q}}: every term of $q$ has
exponent of $v$ at most $1$, the coefficient of $v$ is a constant (every monomial of
\textnormal{\texttt{yk v 1 q}} is empty), and $c\cdot\textnormal{\texttt{inv p c}}\equiv1$ for $c = \textnormal{\texttt{csum (\allowbreak{}yk v 1 q)}}$.
$\textnormal{\texttt{pivVal p v q}} = \textnormal{\texttt{scaleQ (\allowbreak{}(\allowbreak{}p -{} 1) * inv p c) (\allowbreak{}yk v 0 q)}}$, the value $-c^{-1}q_0$ of $v$
where $q = cv+q_0$ vanishes. \textnormal{\texttt{cast\_\allowbreak{}pred\_\allowbreak{}one}} ($p-1\equiv-1$), \textnormal{\texttt{piv\_\allowbreak{}iff}}:
$\textnormal{\texttt{qv \ensuremath{\tau}[v\ensuremath{\mapsto}w] q}} = 0\iff w = \textnormal{\texttt{qv \ensuremath{\tau} (\allowbreak{}pivVal p v q)}}$; \textnormal{\texttt{sum\_\allowbreak{}pin}}:
$\sum_w[G(w)]\,p\,F(w) = p\,F(B)$ when $G(w)\iff w=B$.

\paragraph{[HH] on a path pivot.} \textnormal{\texttt{hhPivotOK p s}}: path $= y::z::ys$, $y\ne z$, $y,z\notin ys$,
$y$ linear in every term, $y$ in no constraint or row, and \textnormal{\texttt{pivOK p z (\allowbreak{}yk y 1 phase)}}. With
$B = \textnormal{\texttt{pivVal p z (\allowbreak{}yk y 1 phase)}}$, \textnormal{\texttt{hhPivotS p s}} removes $y,z$ from the path and
substitutes $z:=B$ (\textnormal{\texttt{normQ p (\allowbreak{}substQ z B \ensuremath{\cdot})}}) in \textnormal{\texttt{yk y 0 phase}}, the constraints, and the
rows. \textnormal{\texttt{hhPivot\_\allowbreak{}sound}}: factor $p$.

\paragraph{[Gauss].} \textnormal{\texttt{gaussOK p A s}}: path $=y::ys$, $y\notin ys$, exponents of $y$ at most
$2$, $y$ in no constraint or row, the coefficient of $y^2$ is constant with
\textnormal{\texttt{csum (\allowbreak{}yk y 2 phase) = A}}, and $4A\cdot\textnormal{\texttt{inv p (\allowbreak{}4A)}}\equiv1$.
\textnormal{\texttt{gaussS p A s}} has path $ys$, the same constraints and rows, and phase
\[
\textnormal{\texttt{normQ p}}\bigl(\textnormal{\texttt{yk y 0 phase}} \mathbin{+\!\!+}\textnormal{\texttt{scaleQ}}((p-1)\,\textnormal{\texttt{inv p (\allowbreak{}4A)}})\,(\textnormal{\texttt{mulQ Q Q}})\bigr),
\]
with $Q = \textnormal{\texttt{yk y 1 phase}}$: the phase $Ay^2+Qy+R$ becomes $R-Q^2/(4A)$.
\begin{theorem}[\textnormal{\texttt{complete\_\allowbreak{}square}}, \textnormal{\texttt{gauss\_\allowbreak{}sound}}]
For $4Ai=1$ in $\Zd p$: $\sum_b\psi_p(R+bQ+b^2A) = \psi_p(R-iQ^2)\sum_b\psi_p(Ab^2)$.
\textnormal{\texttt{gauss\_\allowbreak{}sound}} ($p$ prime): $\textnormal{\texttt{val p c s \ensuremath{\sigma}}} = \textnormal{\texttt{val p (\allowbreak{}qg p A * c) (\allowbreak{}gaussS p A s) \ensuremath{\sigma}}}$.
\end{theorem}

\paragraph{Canonical pivot.} \textnormal{\texttt{canonOK p v s}}: some constraint satisfies \textnormal{\texttt{pivOK p v}} and
$v\notin\text{path}$. With $q$ the first one, $c$ its coefficient of $v$, and $B = \textnormal{\texttt{pivVal p v q}}$,
\textnormal{\texttt{canonS p v s}} substitutes $v:=B$ in the phase, the other constraints, and the rows, and
appends the row \textnormal{\texttt{normQ p (\allowbreak{}scaleQ (\allowbreak{}inv p c) q)}} (coefficient of $v$ normalized to $1$).
\textnormal{\texttt{canon\_\allowbreak{}sound}}: factor $1$.

\paragraph{The final comparison.} \textnormal{\texttt{constK P}}, \textnormal{\texttt{stripC P}}, \textnormal{\texttt{qv\_\allowbreak{}stripC}}
($\textnormal{\texttt{qv \ensuremath{\tau} P}} = \textnormal{\texttt{constK P}} + \textnormal{\texttt{qv \ensuremath{\tau} (\allowbreak{}stripC P)}}$), and \textnormal{\texttt{finalOK sL sR}} as in the
qubit case. \textnormal{\texttt{final\_\allowbreak{}sound}}: if \textnormal{\texttt{finalOK}}, $\textnormal{\texttt{constK sL.\allowbreak{}phase}} = k_L$,
$\textnormal{\texttt{constK sR.\allowbreak{}phase}} = k_R$, and $c_L\omega_p^{k_L} = c_R\omega_p^{k_R}$, the values agree.

\paragraph{Pin.} \textnormal{\texttt{pinTo p y x q}}: \textnormal{\texttt{pivOK p y q}} and \textnormal{\texttt{normQ p (\allowbreak{}pivVal p y q) = [(\allowbreak{}[(\allowbreak{}x,\allowbreak{} 1)],\allowbreak{} 1)]}}
(the constraint sets $y$ to $x$). \textnormal{\texttt{pinOK p x s}}: path $=y::ys$, $y\notin ys$, some constraint
satisfies \textnormal{\texttt{pinTo p y x}}. \textnormal{\texttt{pinS p x s}} removes $y$ and that constraint and substitutes
$y := \textnormal{\texttt{pivVal p y q}}$ elsewhere. \textnormal{\texttt{pin\_\allowbreak{}sound}}: factor $1$.

\begin{definition}[\textnormal{\texttt{QSem.\allowbreak{}Step}}]
\textnormal{\texttt{inductive Step}} has \textnormal{\texttt{perm l}}, \textnormal{\texttt{norm}}, \textnormal{\texttt{elim}}, \textnormal{\texttt{hhCons}}, \textnormal{\texttt{hhPivot}},
\textnormal{\texttt{gauss A}}, \textnormal{\texttt{canon v}}, and \textnormal{\texttt{pin x}}, with \textnormal{\texttt{stepOK p}} the side conditions above,
\textnormal{\texttt{stepS p}} the transformers, and factors $\textnormal{\texttt{fac p}}$: $p$ for \textnormal{\texttt{elim}}, \textnormal{\texttt{hhCons}},
\textnormal{\texttt{hhPivot}}; \textnormal{\texttt{qg p A}} for \textnormal{\texttt{gauss A}}; $1$ otherwise. \textnormal{\texttt{runOK p}}, \textnormal{\texttt{run p}},
\textnormal{\texttt{facs p}} compose them.
\end{definition}
\begin{theorem}[\textnormal{\texttt{QSem.\allowbreak{}step\_\allowbreak{}sound}}, \textnormal{\texttt{QSem.\allowbreak{}run\_\allowbreak{}sound}}]
For $p$ prime: $\textnormal{\texttt{stepOK p st s}}\Rightarrow\textnormal{\texttt{val p c s \ensuremath{\sigma}}} = \textnormal{\texttt{val p (\allowbreak{}fac p st * c) (\allowbreak{}stepS p st s) \ensuremath{\sigma}}}$,
and $\textnormal{\texttt{runOK p sts s}}\Rightarrow\textnormal{\texttt{val p c s \ensuremath{\sigma}}} = \textnormal{\texttt{val p (\allowbreak{}facs p sts * c) (\allowbreak{}run p sts s) \ensuremath{\sigma}}}$.
\end{theorem}

\section{Qudit networks: \texorpdfstring{\file{PathSumQuditNet}}{PathSumQuditNet} and \texorpdfstring{\file{PathSumQuditDigits}}{PathSumQuditDigits}}\label{sec:lean-pathsum-qnet}

\subsection{Labels of dimension \texorpdfstring{$p$}{p}}

\file{PathSumQuditNet} (imports \file{PathSumQudit}, \file{PathSumDigits}, \file{GaussSign},
\file{PathSumQuditChain}) mirrors \file{PathSumNet} over $\Zd p$.
\textnormal{\texttt{QFrag}} has \textnormal{\texttt{path}}, \textnormal{\texttt{phase}}, \textnormal{\texttt{cons : List Poly}}, \textnormal{\texttt{c}};
$\textnormal{\texttt{fval p f \ensuremath{\tau}}} = \textnormal{\texttt{val p f.\allowbreak{}c \ensuremath{\langle}f.\allowbreak{}path,\allowbreak{} f.\allowbreak{}phase,\allowbreak{} f.\allowbreak{}cons,\allowbreak{} []\ensuremath{\rangle} \ensuremath{\tau}}}$; \textnormal{\texttt{QFrag.\allowbreak{}join}},
\textnormal{\texttt{QFrag.\allowbreak{}unit}}, \textnormal{\texttt{qjoinAll}}, and \textnormal{\texttt{QDisj}} are as for qubits, with the theorems
\textnormal{\texttt{fval\_\allowbreak{}join}}, \textnormal{\texttt{fval\_\allowbreak{}joinAll}}, \textnormal{\texttt{fval\_\allowbreak{}unit}}, \textnormal{\texttt{qdisj\_\allowbreak{}join\_\allowbreak{}right}},
\textnormal{\texttt{qdisj\_\allowbreak{}join\_\allowbreak{}left}}, \textnormal{\texttt{qdisj\_\allowbreak{}unit\_\allowbreak{}left}}, \textnormal{\texttt{qdisj\_\allowbreak{}unit\_\allowbreak{}right}},
\textnormal{\texttt{qdisj\_\allowbreak{}joinAll\_\allowbreak{}right}}, \textnormal{\texttt{qdisj\_\allowbreak{}joinAll\_\allowbreak{}left}}, \textnormal{\texttt{summand\_\allowbreak{}join}},
\textnormal{\texttt{summand\_\allowbreak{}frag\_\allowbreak{}inv}}, \textnormal{\texttt{pvars\_\allowbreak{}append}}, \textnormal{\texttt{cvars\_\allowbreak{}append}}, and \textnormal{\texttt{val\_\allowbreak{}smul}}.
$\textnormal{\texttt{zm p \ensuremath{\sigma} l}} = (\sigma(l):\Zd p)$ (\textnormal{\texttt{zm\_\allowbreak{}update}}); \textnormal{\texttt{sumOver\_\allowbreak{}qsum}} turns the network sum
over labels of dimension $p$ into \textnormal{\texttt{qsum}}; \textnormal{\texttt{qv\_\allowbreak{}agree}}.

\begin{definition}[Fragments on dimension $p$]
\textnormal{\texttt{tabOK p y P ts}}: for every $k<p$, $\textnormal{\texttt{qv (\allowbreak{}fun \_\allowbreak{} => k) P}} = ts_k \bmod p$, and $P$ has no
variable but $y$ (\textnormal{\texttt{tab\_\allowbreak{}val}}: then $\psi_p(\textnormal{\texttt{qv \ensuremath{\tau} P}}) = \operatorname{ex}(ts_k/p)$ at
$\tau(y)=k$).
\begin{align*}
\textnormal{\texttt{zFragQ p l0 ls P}} &= \langle[\,],\ P,\ [[([(l_0,1)],1),([(l,1)],p-1)] : l\in ls],\ 1\rangle,\\
\textnormal{\texttt{fFragQ p o i}} &= \langle[\,],\ [([(o,1),(i,1)],1)],\ [\,],\ 1/\sqrt p\rangle,\\
\textnormal{\texttt{xFragQ p y P os is}} &= \langle[y],\ P\mathbin{+\!\!+}[([(y,1),(o,1)],1)]_{o}\mathbin{+\!\!+}[([(y,1),(i,1)],p-1)]_{i},\\
 &\qquad [\,],\ (1/\sqrt p)^{|os|+|is|}\rangle .
\end{align*}
$\textnormal{\texttt{qnetData N fs}} = \langle N.\text{int}\mathbin{+\!\!+}(\textnormal{\texttt{qjoinAll fs}}).\text{path},\ \text{phase},\ \text{cons},\ [\,]\rangle$.
\end{definition}

\begin{theorem}[Generators on dimension $p$]
For $\sigma<p$ everywhere and dimension value $p$:
\textnormal{\texttt{z\_\allowbreak{}frag\_\allowbreak{}q}} (a leg $l_0$, $\varphi_v = ts_v/p$, \textnormal{\texttt{tabOK p l0 P ts}}):
$\textnormal{\texttt{val E \ensuremath{\sigma} n}} = \textnormal{\texttt{fval p (\allowbreak{}zFragQ p l0 (\allowbreak{}outs ++ ins) P) (\allowbreak{}zm p \ensuremath{\sigma})}}$;
\textnormal{\texttt{x\_\allowbreak{}frag\_\allowbreak{}q}} ($y$ not a leg, \textnormal{\texttt{tabOK p y P ts}}): with \textnormal{\texttt{xFragQ p y P outs ins}};
\textnormal{\texttt{f\_\allowbreak{}frag\_\allowbreak{}q}} (\textnormal{\texttt{outs = [o]}}, \textnormal{\texttt{ins = [i]}}, any $\sigma$): with \textnormal{\texttt{fFragQ p o i}}.
Auxiliary: \textnormal{\texttt{fou\_\allowbreak{}q}} ($\textnormal{\texttt{fou p j k}} = p^{-1/2}\psi_p(kj)$), \textnormal{\texttt{conj\_\allowbreak{}fou\_\allowbreak{}q}}
($p^{-1/2}\psi_p((p-1)kj)$), \textnormal{\texttt{prod\_\allowbreak{}terms}}.
\end{theorem}

\begin{definition}[\textnormal{\texttt{QNetPS}}]
\textnormal{\texttt{structure QNetPS E N p ps}} has fields \textnormal{\texttt{nodes}}, \textnormal{\texttt{node}} (entries as fragment values at
\textnormal{\texttt{zm p \ensuremath{\sigma}}} for $\sigma<p$), \textnormal{\texttt{disj}}, \textnormal{\texttt{dims}} (every summed or boundary label has dimension
$p$), and \textnormal{\texttt{cover}}.
\end{definition}
\begin{theorem}[\textnormal{\texttt{QNetPS.\allowbreak{}eval\_\allowbreak{}eq}}, \textnormal{\texttt{equal\_\allowbreak{}qps}}]
$\textnormal{\texttt{eval E N \ensuremath{\sigma}}} = \textnormal{\texttt{val p (\allowbreak{}qjoinAll fs).\allowbreak{}c (\allowbreak{}qnetData N fs) (\allowbreak{}zm p (\allowbreak{}clamp N.\allowbreak{}bnd \ensuremath{\sigma}))}}$ for $\sigma$ in
range on the boundary; two such networks with equal boundaries and agreeing scaled path sums at
every $\tau$ satisfy \textnormal{\texttt{Equal}}.
\end{theorem}

\subsection{Labels of dimension \texorpdfstring{$p^k$}{p\^{}k}}

\file{PathSumQuditDigits} (imports \file{PathSumQuditNet}) represents a label of dimension $p^k$ by
$k$ digit variables over $\Zd p$, low digit first, reusing \textnormal{\texttt{PSem.\allowbreak{}Own}} and \textnormal{\texttt{PSem.\allowbreak{}ownOK}}
(\textnormal{\texttt{PSem.\allowbreak{}ownOK\_\allowbreak{}get?}}: \textnormal{\texttt{ownOK}} and $\mathit{dig}(l)[i]? = x$ give the owner entry).

\begin{definition}[Qudit digit assignments]
$\textnormal{\texttt{setDigs p [] j \ensuremath{\tau}}} = \tau$, $\textnormal{\texttt{setDigs p (\allowbreak{}x :: ds) j \ensuremath{\tau}}} = \textnormal{\texttt{setDigs p ds (\allowbreak{}j / p)}}\,(\tau[x\mapsto j\bmod p])$;
\textnormal{\texttt{qsum\_\allowbreak{}setDigs}}: $\sum_{j<p^{|ds|}}G(\textnormal{\texttt{setDigs p ds j \ensuremath{\tau}}}) = \textnormal{\texttt{qsum ds G \ensuremath{\tau}}}$;
\textnormal{\texttt{setDigs\_\allowbreak{}not}}, \textnormal{\texttt{setDigs\_\allowbreak{}idx}} (variable $i$ gets \textnormal{\texttt{digit p i j}}).
$\textnormal{\texttt{qdig p own \ensuremath{\sigma} x}} = \textnormal{\texttt{digit p i (\allowbreak{}\ensuremath{\sigma} l)}}$ when $x$ looks up to $(l,i)$, else $0$
(\textnormal{\texttt{qdig\_\allowbreak{}of\_\allowbreak{}lookup}}, \textnormal{\texttt{qdig\_\allowbreak{}get}}, \textnormal{\texttt{qdig\_\allowbreak{}update}}).
\textnormal{\texttt{sumOver\_\allowbreak{}qdig}}: for \textnormal{\texttt{ownOK}}, labels of dimension $p^{|\mathit{dig}(l)|}$,
$\textnormal{\texttt{sumOver D L (\allowbreak{}fun \ensuremath{\sigma}' => g (\allowbreak{}qdig p own \ensuremath{\sigma}')) \ensuremath{\sigma}}} = \textnormal{\texttt{qsum (\allowbreak{}L.\allowbreak{}flatMap dig) g (\allowbreak{}qdig p own \ensuremath{\sigma})}}$.
\textnormal{\texttt{eq\_\allowbreak{}of\_\allowbreak{}digitsP}}, \textnormal{\texttt{natCast\_\allowbreak{}inj\_\allowbreak{}lt}}, \textnormal{\texttt{digit\_\allowbreak{}hiP}}.
\end{definition}

\begin{definition}[Agreement]
\textnormal{\texttt{qsum\_\allowbreak{}agree}}, \textnormal{\texttt{summand\_\allowbreak{}agree}}, \textnormal{\texttt{val\_\allowbreak{}agreeQ}} (values agree at assignments agreeing
off the path on the phase and constraint variables). $\textnormal{\texttt{fvarsQ f}}$ lists the variables a
fragment reads off its path (\textnormal{\texttt{fval\_\allowbreak{}agree}}). \textnormal{\texttt{freeIn vs s}}: every variable of the phase,
constraints, and rows is on the path or in $vs$ (\textnormal{\texttt{val\_\allowbreak{}agree\_\allowbreak{}free}}). \textnormal{\texttt{qdig\_\allowbreak{}agree}}: owner
tables that look up $vs$ alike give agreeing digit assignments on $vs$.
\end{definition}

\begin{definition}[Fragments over qudit digits]
$\textnormal{\texttt{eqQ p x y}} = [([(x,1)],1),([(y,1)],p-1)]$ ($x-y$, \textnormal{\texttt{qv\_\allowbreak{}eqQ}}); \textnormal{\texttt{eqConsQ p xs ys}}
zips it (\textnormal{\texttt{mem\_\allowbreak{}eqConsQ\_\allowbreak{}iff}}). $\textnormal{\texttt{tabGo p ds P v ts}}$ checks
$\textnormal{\texttt{qv (\allowbreak{}setDigs p ds (\allowbreak{}v+i) 0) P}} = ts_i$ for each entry of $ts$ (\textnormal{\texttt{tabGo\_\allowbreak{}spec}});
\textnormal{\texttt{tabQD p ds P ts}} is \textnormal{\texttt{tabGo p ds P 0 ts}} together with $\textnormal{\texttt{qv}} = 0$ at every
$v\in[|ts|, p^{|ds|})$, so entries past the end of $ts$ read $0$ (\textnormal{\texttt{tabQD\_\allowbreak{}spec}}).
$\textnormal{\texttt{zFragQD p dig as ls P}} = \langle[\,],\ P,\ \bigcup_{l\in ls}\textnormal{\texttt{eqConsQ p (\allowbreak{}anchorD dig as) (\allowbreak{}dig l)}},\ 1\rangle$;
$\textnormal{\texttt{bFragQ p dig k a b}} = \langle[\,],[\,],\textnormal{\texttt{eqConsQ p (\allowbreak{}dig k) (\allowbreak{}dig b ++ dig a)}},1\rangle$.
\end{definition}

\begin{theorem}[Generators over qudit digits]
\begin{itemize}
\item \textnormal{\texttt{z\_\allowbreak{}frag\_\allowbreak{}qd}}: a Z node on $p^k$ with anchors, $\varphi_v = ts_v/p$,
  \textnormal{\texttt{tabQD p (\allowbreak{}anchorD dig as) P ts}}, repeat-free anchors, and $P$ over the anchors:
  $\textnormal{\texttt{val E \ensuremath{\sigma} n}} = \textnormal{\texttt{fval p (\allowbreak{}zFragQD p dig as (\allowbreak{}outs ++ ins) P) (\allowbreak{}qdig p own \ensuremath{\sigma})}}$.
\item \textnormal{\texttt{b\_\allowbreak{}cond\_\allowbreak{}q}}, \textnormal{\texttt{b\_\allowbreak{}frag\_\allowbreak{}q}}, \textnormal{\texttt{s\_\allowbreak{}frag\_\allowbreak{}q}}: binder and splitter on $p^r$ with
  \textnormal{\texttt{ph}} dimension $p^q$: the constraints hold iff $\sigma(k) = \sigma(a)p^q+\sigma(b)$; the
  entry is the value of \textnormal{\texttt{bFragQ p dig k a b}}.
\item \textnormal{\texttt{lift\_\allowbreak{}q}}: a node every leg of which is a one-digit label with $\mathit{dig}(l) = [l]$,
  and a fragment reading off its path only legs and equal to the node entry at every
  $\sigma<p$, has that fragment's value at \textnormal{\texttt{qdig p own \ensuremath{\sigma}}}. This moves \textnormal{\texttt{z\_\allowbreak{}frag\_\allowbreak{}q}},
  \textnormal{\texttt{x\_\allowbreak{}frag\_\allowbreak{}q}}, and \textnormal{\texttt{f\_\allowbreak{}frag\_\allowbreak{}q}} to digit variables.
\end{itemize}
\end{theorem}

\begin{definition}[\textnormal{\texttt{QNetPSD}}]
$\textnormal{\texttt{qnetDataD dig N fs}}$ puts the summed labels' digits before the fragments' paths.
\textnormal{\texttt{structure QNetPSD E N p own labs dig ps}} has the fields of \textnormal{\texttt{NetPSD}} over $\Zd p$ with
dimensions $p^{|\mathit{dig}(l)|}$.
\end{definition}
\begin{theorem}[\textnormal{\texttt{QNetPSD.\allowbreak{}eval\_\allowbreak{}eq}}, \textnormal{\texttt{equal\_\allowbreak{}qpsd2}}]
$\textnormal{\texttt{eval E N \ensuremath{\sigma}}} = \textnormal{\texttt{val p (\allowbreak{}qjoinAll fs).\allowbreak{}c (\allowbreak{}qnetDataD dig N fs) (\allowbreak{}qdig p own (\allowbreak{}clamp N.\allowbreak{}bnd \ensuremath{\sigma}))}}$.
\textnormal{\texttt{equal\_\allowbreak{}qpsd2}}: if the second side's sum reads off its path only variables of $vs$
(\textnormal{\texttt{freeIn}}), both owner tables look up $vs$ alike, the boundaries are equal, and the scaled
sums agree at every $\tau$, then \textnormal{\texttt{Equal E s\ensuremath{_1} N\ensuremath{_1} s\ensuremath{_2} N\ensuremath{_2}}}.
\end{theorem}

\section{Mixed path sums on dimensions 2, 3, and 6: \texorpdfstring{\file{PathSumMixed}}{PathSumMixed}}\label{sec:lean-pathsum-mixed}

\file{PathSumMixed} (imports \file{PathSumQudit}, \file{PathSumAmp}; namespace
\textnormal{\texttt{ArchytasZX.MSem}}) defines path sums with a qubit part over $\Zd2$ and a qutrit part over
$\Zd3$, coupled through \term{indicator entries}: $\mathbb F_2$ variables that read the indicator
$[a z+b=k]$ of a $\Zd3$ variable $z$. A label of dimension $6$ is the digit pair $(s,f)$ of its
value $2s+f$, $s\in\Zd3$, $f\in\Zd2$.

\begin{definition}[Indicators]\label{def:lean-pathsum-ind}
$\textnormal{\texttt{ind k s}} = [s = k]\in\Zd2$ for $k\in\N$, $s\in\Zd3$. The identities are
\begin{gather*}
\textnormal{\texttt{ind\_one\_two}}:\ \textnormal{\texttt{ind 1 s}}\cdot\textnormal{\texttt{ind 2 s}}=0,\\
\textnormal{\texttt{ind\_zero}}:\ \lfloor\textnormal{\texttt{ind 0 s}}\rfloor = 1-\lfloor\textnormal{\texttt{ind 1 s}}\rfloor-\lfloor\textnormal{\texttt{ind 2 s}}\rfloor,\\
\textnormal{\texttt{ind\_affine}}\ (a\ne0):\ \textnormal{\texttt{ind k (a z + b)}} = \textnormal{\texttt{ind}}\bigl(((k-b)a^{-1}).\mathrm{val}\bigr)\,z,\\
\textnormal{\texttt{ind\_const}}:\ \textnormal{\texttt{ind k b}} = [b=k].
\end{gather*}
\end{definition}

\begin{theorem}[Sums over $\Zd3$ and $6$]\label{thm:lean-pathsum-elim3}
\begin{itemize}
\item \textnormal{\texttt{sum\_six}}: $\sum_{j<6}G(j) = \sum_{s\in\Zd3}\sum_{f\in\Zd2}G(2\,s.\mathrm{val}+f.\mathrm{val})$.
\item \textnormal{\texttt{elim3}} ([Elim3]): $\sum_{s\in\Zd3}\psi_3(sL)\,H = 3[L=0]\,H$, with $\psi_3$ the additive
  character of \cref{def:lean-pathsum-psi}.
\item \textnormal{\texttt{count3}} ([Count3]): for $F:\Zd2\to\Zd2\to\C$,
  $\sum_{s\in\Zd3}F(\textnormal{\texttt{ind 1 s}},\textnormal{\texttt{ind 2 s}}) = \sum_{u,v\in\Zd2}[uv=0]\,F(u,v)$.
\end{itemize}
\end{theorem}

\begin{definition}[Lifts of assignments]\label{def:lean-pathsum-lifta}
\textnormal{\texttt{lift es \ensuremath{\rho} \ensuremath{\tau}}} sets, for each entry $(u,s,k)$ of $es$ (the first entry set last), the
$\mathbb F_2$ variable $u$ to $\textnormal{\texttt{ind k (\ensuremath{\rho} s)}}$. An \term{indicator entry}
\lean!Ent := ℕ × ℕ × ℕ × ℕ × ℕ!, $(u,z,a,b,k)$, makes $u$ the value
$\textnormal{\texttt{entVal \ensuremath{\rho} e}} = \textnormal{\texttt{ind k (a \ensuremath{\rho}(z) + b)}}$; \textnormal{\texttt{liftA es \ensuremath{\rho} \ensuremath{\tau}}} sets every entry's
variable likewise. \textnormal{\texttt{pinE x z' a' b' e}} rewrites an entry on $x$ through
$x=a'z'+b'$: $(u,z',\ a a'\bmod3,\ (a b'+b)\bmod3,\ k)$; other entries stay.
\end{definition}
\noindent Lemmas: \textnormal{\texttt{lift\_out}}, \textnormal{\texttt{lift\_rho}} (a lift reads $\rho$ only at the entries'
variables), \textnormal{\texttt{entVal\_pinE}} and \textnormal{\texttt{liftA\_pinE}} (given $\rho(x)=a'\rho(z')+b'$, \textnormal{\texttt{pinE}}
keeps the values), \textnormal{\texttt{liftA\_rho}}, \textnormal{\texttt{entVal\_norm}} (an entry with a unit coefficient reads
$\textnormal{\texttt{ind}}(((k-b)a^{-1}).\mathrm{val})(\rho z)$). \textnormal{\texttt{lift\_count}} is [Count3] for one $\Zd3$ variable
$s$ read only by the first two entries $(u_1,s,1)$, $(u_2,s,2)$: the sum over $s$ is the sum over
$u_1,u_2\in\Zd2$ with $[u_1u_2=0]$; \textnormal{\texttt{lift\_elim}} is [Elim3] for a variable $s$ that $H$ and $L$
do not read.

\begin{definition}[Mixed value]\label{def:lean-pathsum-mval}
For a qubit path sum with atoms $q$ (\lean!PSem.AData!, \cref{sec:lean-pathsum-amp}), a qudit path sum
$t$ (\lean!QSem.Data! over $p=3$), entries $es$, $\tau:\N\to\Zd2$ and $\rho:\N\to\Zd3$,
\begin{gather*}
\textnormal{\texttt{wt3 t \ensuremath{\rho}}} = [\forall r\in t.\mathrm{cons}\mathbin{+\!\!+}t.\mathrm{rows},\ \textnormal{\texttt{qv \ensuremath{\rho} r}}=0]\ \psi_3(\textnormal{\texttt{qv \ensuremath{\rho} t.phase}}),\\
\textnormal{\texttt{mval M c A q t es \ensuremath{\tau} \ensuremath{\rho}}} = \textnormal{\texttt{qsum}}\ t.\mathrm{path}\ \bigl(\rho'\mapsto\textnormal{\texttt{wt3 t \ensuremath{\rho}'}}\cdot\textnormal{\texttt{valA M c A q (liftA es \ensuremath{\rho}' \ensuremath{\tau})}}\bigr)\ \rho .
\end{gather*}
\end{definition}

\begin{theorem}[Steps on mixed path sums]\label{thm:lean-pathsum-mixed-steps}
\begin{itemize}
\item \textnormal{\texttt{valA\_zero\_atom}}: two held variables $u_1\ne u_2$, both not held in $q$, with the atom
  $(j,[[u_1,u_2]])$ and $A(j)=0$, give $\sum_{x,y\in\Zd2}[xy=0]$ of the value of $q$ at
  $\tau[u_2\mapsto y][u_1\mapsto x]$ (the zero atom $0^{u_1u_2}$). \textnormal{\texttt{bsum\_zero}} sums the zero
  function to $0$.
\item \textnormal{\texttt{mval\_count3}} ([Count3]): if $t.\mathrm{path}=s::rest$, $s\notin rest$, $s$ is off the $\Zd3$ phase
  and constraints, no other entry reads $s$, $A(j)=0$, $u_1\ne u_2$, both not held in $q$, then
  \textnormal{\texttt{mval}} with the entries $(u_1,s,1,0,1)::(u_2,s,1,0,2)::es$ equals \textnormal{\texttt{mval}} with
  $u_1::u_2::q.\mathrm{held}$, the atom $(j,[[u_1,u_2]])$ prepended, the path $rest$, and entries $es$.
\item \textnormal{\texttt{mval\_hhPivot}}: if $t.\mathrm{path}=y::z::ys$, \textnormal{\texttt{hhPivotOK 3 t}}, no entry reads $y$, and
  $es'$ reads at $\rho$ what $es$ reads at $\rho[z\mapsto\textnormal{\texttt{qv \ensuremath{\rho} (pivVal 3 z (yk y 1 t.phase))}}]$, then
  \textnormal{\texttt{mval \ldots t es}} $=3\cdot$\textnormal{\texttt{mval \ldots (hhPivotS 3 t) es'}}.
\item \textnormal{\texttt{mval\_hhCons}}, \textnormal{\texttt{mval\_elim}}: with the head $y$ read by no entry,
  \textnormal{\texttt{hhConsOK t}} (resp.\ \textnormal{\texttt{elimOK t}}) gives the factor $3$ and \textnormal{\texttt{hhConsS t}}
  (resp.\ \textnormal{\texttt{elimS t}}).
\item \textnormal{\texttt{mval\_norm}}: \textnormal{\texttt{normS 3}} of the $\Zd3$ part keeps the value.
\item \textnormal{\texttt{mval\_pin}}: \textnormal{\texttt{pinOK 3 x t}}, $t.\mathrm{path}=y::ys$, $x\ne y$ give
  \textnormal{\texttt{pinS 3 x t}} with the entries mapped by \textnormal{\texttt{pinE y x 1 0}}.
\end{itemize}
\textnormal{\texttt{wt3\_update}} (a variable off the $\Zd3$ phase and constraints does not change the weight),
\textnormal{\texttt{qsum\_const\_mul}}, and the instance \textnormal{\texttt{fact3 : Fact (Nat.Prime 3)}} support the
proofs; each step unfolds \textnormal{\texttt{mval}} and applies the corresponding qudit soundness theorem of
\cref{sec:lean-pathsum-qudit} under \textnormal{\texttt{liftA\_rho}}.
\end{theorem}

\section{Mixed networks: \texorpdfstring{\file{PathSumMixedNet}}{PathSumMixedNet}}\label{sec:lean-pathsum-mixednet}

\file{PathSumMixedNet} (imports \file{PathSumMixed}, \file{PathSumNetAmp}, \file{PathSumQuditNet},
\file{PathSumQuditDigits}; namespace \textnormal{\texttt{ArchytasZX.MSem}}) translates a network whose labels have
dimension $2$, $3$, or $6$ into a mixed path sum. \cref{ch:lean-export} specifies its caller
\lean!export_mixed_net!.

\begin{definition}[Mixed labels]\label{def:lean-pathsum-mlab}
\lean!MLab := ℕ × ℕ × ℕ × ℕ!, $(l,k,b,t)$: the label $l$ of dimension $k\in\{2,3,6\}$ with the
$\mathbb F_2$ variable $b$ (used for $k=2,6$) and the $\Zd3$ variable $t$ (used for $k=3,6$).
\textnormal{\texttt{mbit j}} $= j\bmod 2$ in $\Zd2$; \textnormal{\texttt{mtrit k j}} is $\lfloor j/2\rfloor$ in $\Zd3$ for $k=6$ and $j$
otherwise. \textnormal{\texttt{bvars ms}} lists the $b$ of the labels with $k\ne3$, \textnormal{\texttt{tvars ms}} the $t$ of the
labels with $k\ne2$. \textnormal{\texttt{mbits ms \ensuremath{\sigma}}} and \textnormal{\texttt{mtrits ms \ensuremath{\sigma}}} read the digit assignments from
label values (the first label set last, $0$ elsewhere). \textnormal{\texttt{mOK ms}}: distinct labels, distinct
$\mathbb F_2$ variables, distinct $\Zd3$ variables, every dimension $2$, $3$, or $6$.
\textnormal{\texttt{MIn ms \ensuremath{\sigma}}}: $\sigma(l)<k$ for every $(l,k,\_,\_)\in ms$.
\end{definition}
\noindent Lemmas: \textnormal{\texttt{mbits\_agree}}, \textnormal{\texttt{mtrits\_agree}}, \textnormal{\texttt{mem\_bvars}}, \textnormal{\texttt{mem\_tvars}},
\textnormal{\texttt{mbits\_update}}, \textnormal{\texttt{mtrits\_update}} (setting a label updates its digit variables),
\textnormal{\texttt{mbits\_get}}, \textnormal{\texttt{mtrits\_get}}, \textnormal{\texttt{bsum\_sum3}}, \textnormal{\texttt{sum\_range\_three}}, \textnormal{\texttt{sum\_two}}.

\begin{theorem}[\textnormal{\texttt{sumOver\_mixed}}]
If \textnormal{\texttt{mOK ms}}, every $m\in L$ lies in $ms$ and has $D(l)=k$, then for every $\sigma$
\begin{multline*}
\textnormal{\texttt{sumOver D (L.map (\ensuremath{\cdot}.1))}}\ \bigl(\sigma'\mapsto g(\textnormal{\texttt{mbits ms \ensuremath{\sigma}'}},\textnormal{\texttt{mtrits ms \ensuremath{\sigma}'}})\bigr)\ \sigma\\
 = \textnormal{\texttt{bsum (bvars L)}}\bigl(\tau\mapsto\textnormal{\texttt{qsum (tvars L) (g \ensuremath{\tau}) (mtrits ms \ensuremath{\sigma})}}\bigr)\,(\textnormal{\texttt{mbits ms \ensuremath{\sigma}}}).
\end{multline*}
The proof is induction on $L$ with \textnormal{\texttt{sum\_six}} for dimension $6$.
\end{theorem}

\begin{definition}[\texttt{NetPSM}]\label{def:lean-pathsum-netpsm}
\textnormal{\texttt{NetPSM E N M A ms mi es ps}} is a structure of proofs, with
$ps$ a list of $((n, J), F, \mathit{ats})$: a node, its $\Zd3$ fragment, its $\mathbb F_2$ fragment, and
its atoms. Fields: \textnormal{\texttt{ok}} (\textnormal{\texttt{mOK ms}}); \textnormal{\texttt{nodes}} ($N$'s nodes are the $n$ of $ps$);
\textnormal{\texttt{node}} (for \textnormal{\texttt{MIn ms \ensuremath{\sigma}}},
$\textnormal{\texttt{val E \ensuremath{\sigma} n}} = \textnormal{\texttt{QSem.fval 3 J (mtrits ms \ensuremath{\sigma})}}\cdot\textnormal{\texttt{fvalA M A F ats (liftA es (mtrits ms \ensuremath{\sigma}) (mbits ms \ensuremath{\sigma}))}}$);
\textnormal{\texttt{disj3}} (the $\Zd3$ fragments pairwise \textnormal{\texttt{QDisj}} both ways); \textnormal{\texttt{disj2}} (the $\mathbb F_2$
fragments with atoms pairwise \textnormal{\texttt{DisjA}} both ways); \textnormal{\texttt{dims}} (\textnormal{\texttt{labDim E l = k}} for
every label of $ms$); \textnormal{\texttt{sub}} ($mi\subseteq ms$); \textnormal{\texttt{int}} ($N.\mathrm{int}$ is the labels of $mi$);
\textnormal{\texttt{bnd}} (every boundary label is a label of $ms$); \textnormal{\texttt{cover}} (every node label is summed or
boundary).
\end{definition}

\begin{theorem}[\textnormal{\texttt{NetPSM.eval\_eq}}, \textnormal{\texttt{NetPSM.eval\_mval}}]
Under \textnormal{\texttt{NetPSM}} and $\sigma$ in range on $N.\mathrm{bnd}$, \textnormal{\texttt{eval E N \ensuremath{\sigma}}} is
\textnormal{\texttt{bsum (bvars mi)}} of \textnormal{\texttt{qsum (tvars mi)}} of
$\textnormal{\texttt{QSem.fval 3 (qjoinAll J's)}}\cdot\textnormal{\texttt{fvalA M A (joinAll F's) (atoms) (liftA es \ensuremath{\rho} \ensuremath{\tau})}}$, at
\textnormal{\texttt{mtrits ms (clamp N.bnd \ensuremath{\sigma})}} and \textnormal{\texttt{mbits ms (clamp N.bnd \ensuremath{\sigma})}}. If moreover no summed
$\mathbb F_2$ variable is an entry variable and no entry reads a path variable of the joined $\Zd3$
fragment, \textnormal{\texttt{eval E N \ensuremath{\sigma}}} $= J.c\cdot\textnormal{\texttt{mval M F.c A q t es}}$ with
$q = \langle\langle[\,],F.\mathrm{phase},F.\mathrm{cons},[\,]\rangle,\ \textnormal{\texttt{bvars mi}}\mathbin{+\!\!+}F.\mathrm{path},\ \text{atoms}\rangle$ and
$t=\langle\textnormal{\texttt{tvars mi}}\mathbin{+\!\!+}J.\mathrm{path},\ J.\mathrm{phase},\ J.\mathrm{cons},\ [\,]\rangle$ for the joined
fragments $J$, $F$ (\textnormal{\texttt{sum\_mval}}, from \textnormal{\texttt{bsum\_qsum}}, \textnormal{\texttt{liftA\_update}},
\textnormal{\texttt{bsum\_liftA}}, \textnormal{\texttt{fvalA\_valA}}).
\end{theorem}

\paragraph{Node lemmas.} Each gives the \textnormal{\texttt{node}} field for one kind of generator.
\begin{itemize}
\item \textnormal{\texttt{node\_q3}}: a node on dimension $3$ whose labels $l$ are the mixed labels $(l,3,l,l)$, with a
  $\Zd3$ fragment $f$ over its labels, has entry \textnormal{\texttt{QSem.fval 3 f (mtrits ms \ensuremath{\sigma})}} times the unit
  $\mathbb F_2$ fragment.
\item \textnormal{\texttt{node\_q2}}: a node on dimension $2$ whose labels are $(l,2,l,l)$, $\mathit{dig}(l)=[l]$, not entry
  variables, with an $\mathbb F_2$ fragment with atoms over \textnormal{\texttt{dbits own}} (variables in $vs$, each
  on the fragment's path or a label), has the unit $\Zd3$ fragment times that fragment at
  \textnormal{\texttt{liftA es \ensuremath{\rho} (mbits ms \ensuremath{\sigma})}} (\textnormal{\texttt{fvalA\_agree}}, \textnormal{\texttt{bsum\_agree}}, \textnormal{\texttt{liftA\_out}}).
\item \textnormal{\texttt{node\_b}}: a $(3,2)$ binder or splitter on labels $k$ ($6$), $a$ ($3$), $b$ ($2$), each its own
  digit variable and $k,b$ not entry variables, is
  $\textnormal{\texttt{bFrag3 k a}}=\langle[\,],[\,],[\textnormal{\texttt{eqQ 3 k a}}],1\rangle$ times
  $\textnormal{\texttt{bFrag2 k b}}=\langle[\,],[\,],[[[k],[b]]],1\rangle$ (\textnormal{\texttt{tB\_six}}:
  $\mathrm{tB}(3,2,x,y,z) = [\textnormal{\texttt{mtrit 6 x}}=\textnormal{\texttt{mtrit 3 y}}]\,[\textnormal{\texttt{mbit x}}=\textnormal{\texttt{mbit z}}]$ in range;
  \textnormal{\texttt{bFrag\_val}}; \textnormal{\texttt{val\_binder}}, \textnormal{\texttt{val\_splitter}} give the entries of \lean!B! and \lean!S! with
  $\mathrm{dimv}(n.\mathrm{dim})=3$, $\mathrm{dimv}(n.\mathrm{ph})=2$).
\item \textnormal{\texttt{node\_z6}}: a Z spider on dimension $6$ with anchor $l_0$, indicator variables $u_1$, $u_2$ and
  entries $(u_1,l_0,1,0,1)$, $(u_2,l_0,1,0,2)$, phases $\varphi_v = ts_v/M$, a table
  \textnormal{\texttt{tableOK M [l0, u1, u2] P ts8}} (\cref{sec:lean-pathsum-digits}) and
  \textnormal{\texttt{sixOK ts8 ts}}, is $\textnormal{\texttt{z6Frag3 l0 legs}}=\langle[\,],[\,],[\textnormal{\texttt{eqQ 3 l0 l}}]_l,1\rangle$ times
  $\textnormal{\texttt{z6Frag2 l0 legs P}}=\langle[\,],P,[[[l_0],[l]]]_l,1\rangle$. Here
  $\textnormal{\texttt{sixIdx x}} = x\bmod2+2[\lfloor x/2\rfloor=1]+4[\lfloor x/2\rfloor=2]$ indexes the value $x=2s+f$ on the
  variables $(f,u_1,u_2)$, and \textnormal{\texttt{sixOK ts8 ts}} states $ts8[\textnormal{\texttt{sixIdx x}}]=ts[x]$ for
  $x<6$ (\textnormal{\texttt{six\_eq}}, \textnormal{\texttt{six\_bits}}, \textnormal{\texttt{sixIdx\_lt}}, \textnormal{\texttt{liftA\_get}}).
\end{itemize}

\paragraph{Steps of the $\Zd3$ replay.}
\textnormal{\texttt{mval\_perm}} permutes the $\Zd3$ path. \textnormal{\texttt{substA M z B q}} replaces $z$ by $B$ in the phase
(\textnormal{\texttt{substP}}), constraints, rows, and atoms (\textnormal{\texttt{substF}}) of $q$; \textnormal{\texttt{valA\_subst}} ($2\le M$; $z$
and the variables of $B$ neither held nor on the path; no path variable read by the atoms before or
after): the value of $q$ at $\tau[z\mapsto\textnormal{\texttt{fz \ensuremath{\tau} B}}]$ is the value of \textnormal{\texttt{substA M z B q}} at $\tau$
(\textnormal{\texttt{bsum\_update\_dep}}). \textnormal{\texttt{liftA\_dup}}: two entries with one reading take one value;
\textnormal{\texttt{mval\_dup}} drops the entry $e'$ whose reading equals that of $e$ and substitutes $[[e.1]]$ for
$e'.1$ in $q$. \textnormal{\texttt{entN e}} puts an entry with a unit coefficient in the form
$(u,z,1,0,((k+2b)a)\bmod 3)$ (\textnormal{\texttt{ind\_unit}}, \textnormal{\texttt{entVal\_entN}}, \textnormal{\texttt{liftA\_entN}},
\textnormal{\texttt{liftA\_perm}}); \textnormal{\texttt{mval\_entN}} maps every entry by it.

\paragraph{The end of the replay.} \textnormal{\texttt{cB e}} is $[[\,]]$ when $b\equiv k\pmod 3$ and $[\,]$ otherwise, the
constant of an entry with coefficient $a\equiv0$ (\textnormal{\texttt{entVal\_const}}, \textnormal{\texttt{cB\_nil}});
\textnormal{\texttt{substEs M es q}} substitutes every entry variable by its constant, first entry first
(\textnormal{\texttt{mem\_substF\_const}}, \textnormal{\texttt{avars\_substA\_const}}, \textnormal{\texttt{valA\_substEs}}).
\textnormal{\texttt{mval\_fin}}: with no $\Zd3$ path left and constant entries,
$\textnormal{\texttt{mval}} = \textnormal{\texttt{wt3 t \ensuremath{\rho}}}\cdot\textnormal{\texttt{valA M c A (substEs M es q) \ensuremath{\tau}}}$; \textnormal{\texttt{wt3\_one}}: zero phase and
zero constraints give weight $1$; \textnormal{\texttt{mval\_end}}: then, with no entries, $\textnormal{\texttt{mval}} = \textnormal{\texttt{valA M c A q \ensuremath{\tau}}}$.

\paragraph{Disjointness in one pass.} \textnormal{\texttt{disjG P V fs A B}} generalizes \textnormal{\texttt{disjB}}
(\cref{def:lean-pathsum-disjb}) to any item type with projections $P$ (path variables) and $V$
(variables); \textnormal{\texttt{disjG\_sound}}. \textnormal{\texttt{disj3\_of}} and \textnormal{\texttt{disj2\_of}} give the \textnormal{\texttt{disj3}} and
\textnormal{\texttt{disj2}} fields of \textnormal{\texttt{NetPSM}} from \textnormal{\texttt{disjG \ldots{} fs 0 0 = true}}. \textnormal{\texttt{mbits\_dbits}}: for
labels $(x,2,x,x)$ with $\mathit{dig}(x)=[x]$ under \textnormal{\texttt{ownOK}}, \textnormal{\texttt{mbits ms \ensuremath{\sigma} x}} $=$
\textnormal{\texttt{dbits own \ensuremath{\sigma} x}}.

\section{Auxiliary lemmas}\label{sec:lean-pathsum-aux}

The following lemmas serve the proofs above; a reimplementation needs them only to rebuild the
proofs.
\begin{itemize}
\item \file{PathSumSem}: \textnormal{\texttt{mem\_\allowbreak{}insU}} ($w\in\textnormal{\texttt{insU v l}}\iff w=v\vee w\in l$),
  \textnormal{\texttt{mem\_\allowbreak{}foldr\_\allowbreak{}insU}}, \textnormal{\texttt{pnM\_\allowbreak{}nil}}, \textnormal{\texttt{pnM\_\allowbreak{}cons}}, \textnormal{\texttt{pnM\_\allowbreak{}append}} (the mask phase value
  is additive), \textnormal{\texttt{mz\_\allowbreak{}squash}} (\textnormal{\texttt{squash}} keeps the value), \textnormal{\texttt{pn\_\allowbreak{}mapmul}}
  ($\textnormal{\texttt{pn \ensuremath{\sigma} (\allowbreak{}L.\allowbreak{}map fun t => (\allowbreak{}m ++ t.\allowbreak{}1,\allowbreak{} c * t.\allowbreak{}2))}} = c\lfloor\textnormal{\texttt{mz \ensuremath{\sigma} m}}\rfloor\,\textnormal{\texttt{pn \ensuremath{\sigma} L}}$),
  \textnormal{\texttt{pnM\_\allowbreak{}mapor}} (the mask version), \textnormal{\texttt{pn\_\allowbreak{}liftNM}} (\textnormal{\texttt{liftNM}} is $c$ times the lift mod
  $M$), \textnormal{\texttt{fz\_\allowbreak{}mapprefix}} ($\textnormal{\texttt{fz \ensuremath{\sigma} (\allowbreak{}B.\allowbreak{}map (\allowbreak{}m ++ \ensuremath{\cdot}))}} = \textnormal{\texttt{mz \ensuremath{\sigma} m}}\cdot\textnormal{\texttt{fz \ensuremath{\sigma} B}}$),
  \textnormal{\texttt{zsum\_\allowbreak{}cast}} (the integer sum of $\lfloor x_i\rfloor$ reduces mod $2$ to $\sum x_i$),
  \textnormal{\texttt{fz\_\allowbreak{}f2mulM}}, \textnormal{\texttt{fz\_\allowbreak{}f2mulN}} (products of $\mathbb F_2$ polynomials, on masks),
  \textnormal{\texttt{fz\_\allowbreak{}gatherM}} ($\textnormal{\texttt{fz}} = 1$ of the mask indicator iff all constraints hold), and
  \textnormal{\texttt{z2\_\allowbreak{}add\_\allowbreak{}self}} ($x+x=0$).
\item \file{PathSumRules}: \textnormal{\texttt{quot\_\allowbreak{}free}} (the quotient of $y$ omits $y$), \textnormal{\texttt{fz\_\allowbreak{}free}},
  \textnormal{\texttt{splitAt\_\allowbreak{}free1}}, \textnormal{\texttt{splitAt\_\allowbreak{}free2}} (both parts of \textnormal{\texttt{splitAt z}} omit $z$),
  \textnormal{\texttt{splitAt\_\allowbreak{}sub1}}, \textnormal{\texttt{splitAt\_\allowbreak{}sub2}} (both parts omit any $y$ that $Q$ omits),
  \textnormal{\texttt{om8\_\allowbreak{}ex}} ($\omega_8^k = \operatorname{ex}(k/8)$), \textnormal{\texttt{zq\_\allowbreak{}le}} ($\lfloor x\rfloor\in\{0,1\}$),
  \textnormal{\texttt{pn\_\allowbreak{}filter\_\allowbreak{}split}} (a phase splits into the terms failing and passing a filter),
  \textnormal{\texttt{not\_\allowbreak{}hasAny}}, \textnormal{\texttt{pn\_\allowbreak{}update\_\allowbreak{}free}}, and \textnormal{\texttt{fvars\_\allowbreak{}mem}}.
\item \file{PathSumAmp}: \textnormal{\texttt{aval\_\allowbreak{}drop}} (atoms on $0$ contribute $1$), \textnormal{\texttt{z2\_\allowbreak{}one\_\allowbreak{}add\_\allowbreak{}one}}
  ($1+1=0$), \textnormal{\texttt{fz\_\allowbreak{}oneY}} ($\textnormal{\texttt{fz \ensuremath{\tau} (\allowbreak{}oneY y g)}} = (1+\tau(y))\textnormal{\texttt{fz \ensuremath{\tau} g}}$), \textnormal{\texttt{fz\_\allowbreak{}onlyY}}
  ($=\tau(y)\textnormal{\texttt{fz \ensuremath{\tau} g}}$), \textnormal{\texttt{avars\_\allowbreak{}not\_\allowbreak{}reads}} (the atoms not reading $y$ omit $y$),
  \textnormal{\texttt{fz\_\allowbreak{}of\_\allowbreak{}normF}} (equal normal forms have equal values), \textnormal{\texttt{aval\_\allowbreak{}map}} (a value-preserving
  map of the polynomials keeps \textnormal{\texttt{aval}}), \textnormal{\texttt{offPath\_\allowbreak{}not}}, \textnormal{\texttt{mz\_\allowbreak{}of\_\allowbreak{}mHas}},
  \textnormal{\texttt{fz\_\allowbreak{}flatMap}}, \textnormal{\texttt{sum\_\allowbreak{}ind}} (the sum of indicators $[N=n]$ over a repeat-free list is
  $[N\in L]$), \textnormal{\texttt{fz\_\allowbreak{}at}} (an atom after substituting $y:=B$ reads \textnormal{\texttt{tSel}}), \textnormal{\texttt{aval\_\allowbreak{}ctx}}
  (the atoms at $y=0$ and $y=1$ are the products over \textnormal{\texttt{ctx0}} and \textnormal{\texttt{ctx1}}),
  \textnormal{\texttt{prod\_\allowbreak{}range\_\allowbreak{}hit}} ($\prod_{n<K}f(n)^{[N=n]} = f(N)$ for $N<K$), and \textnormal{\texttt{fz\_\allowbreak{}zero\_\allowbreak{}of}}.
\item \file{PathSumNet}: \textnormal{\texttt{bit\_\allowbreak{}cases}} ($n<2\Rightarrow n\in\{0,1\}$) and \textnormal{\texttt{zq\_\allowbreak{}bit}}
  ($\lfloor(n:\Zd2)\rfloor = n$ for $n<2$).
\item \file{PathSumQudit}: \textnormal{\texttt{yexp\_\allowbreak{}nil}}, \textnormal{\texttt{yexp\_\allowbreak{}cons}}, \textnormal{\texttt{strip\_\allowbreak{}cons}}, \textnormal{\texttt{yk\_\allowbreak{}cons}},
  \textnormal{\texttt{not\_\allowbreak{}mem\_\allowbreak{}strip}}, and \textnormal{\texttt{qm\_\allowbreak{}mergeE}} ($\textnormal{\texttt{qm \ensuremath{\sigma} (\allowbreak{}mergeE t l)}} = \sigma(t_1)^{t_2}\textnormal{\texttt{qm \ensuremath{\sigma} l}}$).
\item \file{PathSumChain} and \file{PathSumQuditChain}: \textnormal{\texttt{dataStatesImpl}},
  \textnormal{\texttt{adataStatesImpl}}, and \textnormal{\texttt{qdataStatesImpl}} are the \textnormal{\texttt{unsafe}} compiled implementations of
  \textnormal{\texttt{dataStates}}, \textnormal{\texttt{adataStates}}, and \textnormal{\texttt{qdataStates}}; \textnormal{\texttt{addLit n type value}} adds an
  abbreviation (reducibility hint \textnormal{\texttt{.\allowbreak{}abbrev}}, safe), and \textnormal{\texttt{addThm n type value}} adds a
  theorem with a given proof term, elaborated synchronously.
\end{itemize}

\section{Proof techniques}\label{sec:lean-pathsum-proofs}

Every soundness theorem follows one pattern. The step's side condition is unfolded with
\textnormal{\texttt{Bool.\allowbreak{}and\_\allowbreak{}eq\_\allowbreak{}true}}, \textnormal{\texttt{decide\_\allowbreak{}eq\_\allowbreak{}true\_\allowbreak{}eq}}, and \textnormal{\texttt{Bool.\allowbreak{}not\_\allowbreak{}eq\_\allowbreak{}true'}} into propositions.
The summed head variable $y$ is moved inside the remaining sum by \textnormal{\texttt{bsum\_\allowbreak{}head}} (qudit:
\textnormal{\texttt{qsum\_\allowbreak{}head}}), which needs $y\notin ys$; the summand at $y=b$ is rewritten with the update lemmas
into a factor depending on $b$ and a summand that no longer reads $y$; the pointwise identity in
$\Zd2$ or $\Zd p$ is closed by \textnormal{\texttt{decide}} (qubit identities over the four values of two bits) or
by a ring identity, and \textnormal{\texttt{bsum\_\allowbreak{}congr}} lifts it to the sum. Steps that substitute a variable
($\textnormal{\texttt{canon}}$, $\textnormal{\texttt{bcanon}}$, $\textnormal{\texttt{pin}}$, $\textnormal{\texttt{hhPivot}}$, [HHgen]) use the substitution lemmas
of \cref{prop:lean-pathsum-subst} under the support indicator: on the assignments where the
constraint holds, the variable equals its value, so replacing it changes no summand that the
indicator keeps. Normalizations are value-preserving by \cref{prop:lean-pathsum-norm} and never by
evaluation.

The concrete instances are where \textnormal{\texttt{decide}} does the work: the certificate states every side
condition, every literal intermediate state (\cref{sec:lean-pathsum-chain}), the constant-term
values, and the equality of the final rows and phases as closed Boolean terms, and the Lean
kernel reduces them. The bounded recursion of \textnormal{\texttt{mergeN}}, \textnormal{\texttt{msortN}}, and \textnormal{\texttt{sortB}}, and
the mask paths of \textnormal{\texttt{normF}} (for 32 or more monomials) and \textnormal{\texttt{liftN}}, keep these reductions
within the kernel's limits; \textnormal{\texttt{normFM}} chooses the bit-set path when the largest mask is below
$4|Q|+64$, so the bit set it scans has at most that many bits.

\section{Worked examples}\label{sec:lean-pathsum-examples}

\begin{example}[The Hadamard squared]\label{ex:lean-pathsum-hh}
\begin{center}
\begin{tikzpicture}[zxd]
\node[fbox] (a) at (0,1) {F};
\node[fbox] (b) at (0,2) {F};
\draw[wire] (0,0) -- (a) -- (b) -- (0,3);
\node[dimlabel,right] at (0,0.4) {1};
\node[dimlabel,right] at (0,1.5) {$2^{16}$};
\node[dimlabel,right] at (0,2.6) {0};
\node at (1.4,1.5) {$=$};
\draw[wire] (2.4,0) -- (2.4,3);
\node[dimlabel,right] at (2.4,0.4) {1};
\node[dimlabel,right] at (2.4,2.6) {0};
\end{tikzpicture}
\end{center}
The two qubit Fourier boxes $F$ composed (input port $1$ at the bottom, output port $0$ at the top,
the summed wire labeled $2^{16} = 65536$ in the network export) equal the bare wire. The engine's
path sums are $\frac12\sum_{y}(-1)^{y x_0 + y x_1}$ (path variable $2$, constant $(1/\sqrt2)^2$) and
the wire, a constraint $x_0+x_1=0$. With $M=4$ (the minimum of \cref{alg:lean-pathsum-plan}), the
left phase is $[([0,2],2),([1,2],2)]$: $2/4$ turn on each of $x_0y$ and $x_1y$. The step
\textnormal{\texttt{.\allowbreak{}hhCons}} sums $y$ into the constraint $x_0+x_1$ with factor $2$, and \textnormal{\texttt{.\allowbreak{}canon 0}} turns
the constraint into the row $x_0+x_1$ on both sides. \textnormal{\texttt{export\_\allowbreak{}proof}} writes:
\begin{lstlisting}[style=lean]
def sL : Data := { path := [2], phase := [([0, 2], 2), ([1, 2], 2)], cons := [], rows := [] }
def sR : Data := { path := [], phase := [], cons := [[[0], [1]]], rows := [] }
def stepsL : List Step := [.front 2, .hhCons, .canon 0]
def stepsR : List Step := [.canon 0]

theorem facL : facs stepsL = (2 : ℂ) := by
  simp only [stepsL, facs, fac]
  ring

theorem facR : facs stepsR = 1 := by
  simp only [stepsR, facs, fac]
  ring

theorem consts : ((2 : ℂ)) * (((rt 2))⁻¹ ^ 2) * om 4 ^ 0 =
    (1) * 1 * om 4 ^ 0 := by
  have key : ((2 : ℂ) * (((rt 2))⁻¹ ^ 2)) = 1 :=
    by have r0 := (mul_inv_cancel₀ (rt_ne 2 (by norm_num))); have r1 := (rt_sq' 2);
       linear_combination ((1 + (((rt 2))⁻¹ * (rt 2)))) * r0
         + (((-1 : ℂ) * (((rt 2))⁻¹ ^ 2))) * r1
  linear_combination key

run_chain runL 4 stepsL sL 4
run_chain runR 4 stepsR sR 4

theorem main (σ : ℕ → ZMod 2) :
    PSem.val 4 (((rt 2))⁻¹ ^ 2) sL σ = PSem.val 4 1 sR σ := by
  rw [run_sound 4 (by norm_num) stepsL sL _ σ runL_ok, runL_eq,
    run_sound 4 (by norm_num) stepsR sR _ σ runR_ok, runR_eq]
  apply final_sound 4 (by norm_num) _ _ _ _ σ 0 0 (by decide) (by decide)
    (by decide)
  simpa only [facL, facR] using consts
\end{lstlisting}
(the \textnormal{\texttt{have key}} proof is one line in the file; it is wrapped here). \textnormal{\texttt{export\_\allowbreak{}pair}} for the
same pair writes the comment \textnormal{\texttt{-{}-{} L step 0: [hh\_\allowbreak{}constraint] var 2,\allowbreak{} value [[0],\allowbreak{} [1]]}}, the line
\textnormal{\texttt{example := @ps\_\allowbreak{}hh}}, and the two equal records
\textnormal{\texttt{\{ rows := [(\allowbreak{}0,\allowbreak{} [[1]])],\allowbreak{} phase := [],\allowbreak{} const := [(\allowbreak{}1 : \ensuremath{\mathbb{Q}})],\allowbreak{} n := 1 \}}}.
\end{example}

\begin{example}[The same pair as networks]
\textnormal{\texttt{export\_\allowbreak{}net}} writes the environment \textnormal{\texttt{E}} with \textnormal{\texttt{labs := [(\allowbreak{}0,\allowbreak{} 1),\allowbreak{} (\allowbreak{}1,\allowbreak{} 1),\allowbreak{} (\allowbreak{}65536,\allowbreak{} 1)]}}
(one digit per label), the owner table \textnormal{\texttt{ownL := [(\allowbreak{}0,\allowbreak{} (\allowbreak{}0,\allowbreak{} 0)),\allowbreak{} (\allowbreak{}1,\allowbreak{} (\allowbreak{}1,\allowbreak{} 0)),\allowbreak{} (\allowbreak{}2,\allowbreak{} (\allowbreak{}65536,\allowbreak{} 0))]}}
(the summed wire's digit is the engine's path variable $2$), and
\begin{lstlisting}[style=lean]
def NL : Net := ⟨[⟨.F, [65536], [1], 1, 0⟩, ⟨.F, [0], [65536], 1, 0⟩], [65536], [0, 1]⟩
noncomputable def psL : List (Node × Frag) :=
  [(⟨.F, [65536], [1], 1, 0⟩, fFragD 4 2 1 digL 65536 1),
   (⟨.F, [0], [65536], 1, 0⟩, fFragD 4 2 1 digL 0 65536)]
theorem netL : NetPSD E NL 4 ownL labsL digL psL where
  ownok := ownLOK
  nodes := rfl
  node := by
    intro p hp σ hσ
    have hown := ownLOK
    simp only [psL, List.mem_cons, List.not_mem_nil, or_false] at hp
    rcases hp with rfl | rfl
    · exact f_frag_dig E σ ownL labsL digL hown hσ _ rfl 1 rfl 65536 1 rfl rfl
        (by decide) (by decide) rfl rfl 4 2 2 (by norm_num) (by norm_num) (by norm_num)
    · exact f_frag_dig E σ ownL labsL digL hown hσ _ rfl 1 rfl 0 65536 rfl rfl
        (by decide) (by decide) rfl rfl 4 2 2 (by norm_num) (by norm_num) (by norm_num)
  disj := by decide
  dims := by decide
  cover := by decide
  sub := by decide
\end{lstlisting}
The right side is \textnormal{\texttt{NR := \ensuremath{\langle}[\ensuremath{\langle}.\allowbreak{}Z,\allowbreak{} [0],\allowbreak{} [1],\allowbreak{} 1,\allowbreak{} 0\ensuremath{\rangle}],\allowbreak{} [],\allowbreak{} [0,\allowbreak{} 1]\ensuremath{\rangle}}} with the fragment
\textnormal{\texttt{zFragD digR [0] [0,\allowbreak{} 1] []}} proved by \textnormal{\texttt{z\_\allowbreak{}frag\_\allowbreak{}dig}}. The step lists are
\textnormal{\texttt{[.\allowbreak{}norm,\allowbreak{} .\allowbreak{}perm [2],\allowbreak{} .\allowbreak{}perm [2],\allowbreak{} .\allowbreak{}hhCons,\allowbreak{} .\allowbreak{}canon 0]}} and \textnormal{\texttt{[.\allowbreak{}norm,\allowbreak{} .\allowbreak{}canon 0]}} (the network
sum starts with a \textnormal{\texttt{.\allowbreak{}norm}} step and its own path order), and the theorem
\begin{lstlisting}[style=lean]
theorem hh : Equal E (1) NL (1) NR := by
  apply equal_psd2 _ _ netL netR rfl
  intro σ
  rw [cL, cR, ← val_smul, ← val_smul,
    run_sound 4 (by norm_num) stepsL (netDataD digL NL (psL.map Prod.snd)) _ _ runL_ok, runL_eq,
    run_sound 4 (by norm_num) stepsR (netDataD digR NR (psR.map Prod.snd)) _ _ runR_ok, runR_eq,
    val_agree 4 _ runL [0, 1] _ _ (by decide) (by decide)
      (dbits_agree ownL ownR _ [0, 1] (by decide))]
  apply final_sound 4 (by norm_num) _ _ _ _ _ 0 0 (by decide) (by decide)
    (by decide)
  first | (rw [facL, facR]; exact consts) | rw [facL, facR]
\end{lstlisting}
reads the two final sums at the two owner tables' digit assignments, which agree on the boundary
digits $[0,1]$.
\end{example}

\begin{example}[Four qutrit Fourier boxes]
For $p=3$, $F^4$ equals the wire. The left path sum has path $[2,3,4]$ and phase
$x_0x_4 + x_1x_2 + x_2x_3 + x_3x_4$ (each coefficient $1$, worth $1/3$ turn), the right the
constraint $x_0 + 2x_1 = 0$ ($x_0 - x_1$). \textnormal{\texttt{export\_\allowbreak{}qudit\_\allowbreak{}proof}} writes
\textnormal{\texttt{stepsL := [.\allowbreak{}perm [2,\allowbreak{} 3,\allowbreak{} 4],\allowbreak{} .\allowbreak{}hhPivot,\allowbreak{} .\allowbreak{}perm [4],\allowbreak{} .\allowbreak{}hhCons,\allowbreak{} .\allowbreak{}canon 0,\allowbreak{} .\allowbreak{}norm]}} and
\textnormal{\texttt{stepsR := [.\allowbreak{}canon 0,\allowbreak{} .\allowbreak{}norm]}}, the Legendre values \textnormal{\texttt{lg 3 0 = 0}}, \textnormal{\texttt{lg 3 1 = 1}},
\textnormal{\texttt{lg 3 2 = -{}1}} by \textnormal{\texttt{decide}}, \textnormal{\texttt{facL : QSem.\allowbreak{}facs 3 stepsL = (\allowbreak{}3 : \ensuremath{\mathbb{C}}) * (\allowbreak{}3 : \ensuremath{\mathbb{C}})}}, and the
constants' identity over $\omega_{12}$ (the lcm of $p$ and the constant's order $4$, from
$(1/\sqrt3)^4 = 1/9$ written as $\frac{2}{27}-\frac1{27}\omega_{12}^4-\frac1{27}\omega_{12}^8$),
proved after \textnormal{\texttt{om\_\allowbreak{}sub 12 3 4}} from $1+\omega_{12}^4+\omega_{12}^8=0$ (\textnormal{\texttt{om\_\allowbreak{}geom}}). The main
theorem is \textnormal{\texttt{QSem.\allowbreak{}val 3 c sL \ensuremath{\sigma} = QSem.\allowbreak{}val 3 (\allowbreak{}1) sR \ensuremath{\sigma}}} for $\sigma:\N\to\Zd3$, closed by
\textnormal{\texttt{QSem.\allowbreak{}run\_\allowbreak{}sound}} twice and \textnormal{\texttt{QSem.\allowbreak{}final\_\allowbreak{}sound}}.
\end{example}

\section{The Python exporters}\label{sec:lean-pathsum-python}

Four Python modules write Lean files against the modules of this chapter. Every exporter reads
the engine's path sums (\textnormal{\texttt{from\_\allowbreak{}diagram}}, \textnormal{\texttt{reduce}}, \textnormal{\texttt{canonical}}, \textnormal{\texttt{reduce\_\allowbreak{}qudit}},
\textnormal{\texttt{canonical\_\allowbreak{}qudit}}, \textnormal{\texttt{eliminate}}, \textnormal{\texttt{boolean\_\allowbreak{}eliminate}} of
\module{archytaszx.algebra.pathsum}, \cref{ch:semantics}) and raises \textnormal{\texttt{PathSumUnsupported}} when a
pair falls outside the class it covers. Constants are rendered and proved by the scalar renderer
of \module{archytaszx.semantics.lean_scalar} (\textnormal{\texttt{Scalars}} with atoms \textnormal{\texttt{Om(\allowbreak{}n)}} for $\omega_n$,
\textnormal{\texttt{Rt(\allowbreak{}2)}} for $\sqrt2$, and \textnormal{\texttt{Inv}}), configured by \textnormal{\texttt{\_\allowbreak{}scalars(\allowbreak{})}} with natural symbols printed as
themselves, other symbols typed \textnormal{\texttt{\ensuremath{\mathbb{C}}}}, and positivity side goals \textnormal{\texttt{(\allowbreak{}by norm\_\allowbreak{}num)}}; its
\textnormal{\texttt{prove(\allowbreak{}lhs,\allowbreak{} rhs)}} returns a Lean proof text of the identity and \textnormal{\texttt{text(\allowbreak{}e)}} a Lean term.

\subsection{\texorpdfstring{\module{archytaszx.semantics.lean_pathsum}}{lean\_pathsum}}

\paragraph{Constants.} \textnormal{\texttt{PIN\_\allowbreak{}OFFSET = 1 << 20}} (first variable for pinned digits of qubit
networks), \textnormal{\texttt{QPIN\_\allowbreak{}OFFSET = 1 << 18}} (first label for pinned qudit labels), \textnormal{\texttt{CHAIN\_\allowbreak{}BLOCK = 20}}
(block length of the pin-step chains of [Quarter] network proofs), \textnormal{\texttt{RUN\_\allowbreak{}BLOCK = 4}} (block length of
every other \textnormal{\texttt{run\_\allowbreak{}chain}}, \textnormal{\texttt{aruns\_\allowbreak{}chain}}, and \textnormal{\texttt{qrun\_\allowbreak{}chain}}), \textnormal{\texttt{PROOF\_\allowbreak{}OPTIONS = [set\_\allowbreak{}option maxRecDepth 100000,\allowbreak{} set\_\allowbreak{}option maxHeartbeats 0]}},
\textnormal{\texttt{NET\_\allowbreak{}OPTIONS}} (\textnormal{\texttt{maxRecDepth 100000}}, \textnormal{\texttt{maxHeartbeats 4000000}},
\textnormal{\texttt{linter.\allowbreak{}unusedSimpArgs false}}, \textnormal{\texttt{linter.\allowbreak{}unusedTactic false}}), and \textnormal{\texttt{QNET\_\allowbreak{}OPTIONS}}
(\textnormal{\texttt{NET\_\allowbreak{}OPTIONS}} plus \textnormal{\texttt{linter.\allowbreak{}unreachableTactic false}} and \textnormal{\texttt{linter.\allowbreak{}unusedVariables false}}).
The exported names are \textnormal{\texttt{NET\_\allowbreak{}OPTIONS}}, \textnormal{\texttt{QNET\_\allowbreak{}OPTIONS}}, \textnormal{\texttt{NetBody}}, \textnormal{\texttt{export\_\allowbreak{}net}},
\textnormal{\texttt{export\_\allowbreak{}pair}}, \textnormal{\texttt{export\_\allowbreak{}proof}}, \textnormal{\texttt{export\_\allowbreak{}qudit\_\allowbreak{}net}}, \textnormal{\texttt{export\_\allowbreak{}qudit\_\allowbreak{}proof}},
\textnormal{\texttt{net\_\allowbreak{}body}}, \textnormal{\texttt{net\_\allowbreak{}env}}, and \textnormal{\texttt{qnet\_\allowbreak{}body}}.

\paragraph{Literals.} \textnormal{\texttt{\_\allowbreak{}data\_\allowbreak{}term(\allowbreak{}ps,\allowbreak{} m)}} renders a path sum as
\textnormal{\texttt{\{ path := [.\allowbreak{}.\allowbreak{}.\allowbreak{}],\allowbreak{} phase := [.\allowbreak{}.\allowbreak{}.\allowbreak{}],\allowbreak{} cons := [.\allowbreak{}.\allowbreak{}.\allowbreak{}],\allowbreak{} rows := [] \}}}: the path sorted; the phase terms
sorted by (monomial size, sorted monomial), each $(\text{sorted monomial}, \lfloor cm\rfloor\bmod m)$
for a coefficient $c$ in turns; each constraint a list of sorted monomials sorted by their sorted
lists. \textnormal{\texttt{\_\allowbreak{}data\_\allowbreak{}text(\allowbreak{}ps,\allowbreak{} m,\allowbreak{} path)}} is the same with a given path order.

\begin{signature}
\textnormal{\texttt{export\_\allowbreak{}pair(\allowbreak{}left: Diagram,\allowbreak{} right: Diagram,\allowbreak{} name: str = "pair") -{}> str}}
\end{signature}
\noindent Writes a file importing \file{ArchytasZX.PathSum} in namespace
\textnormal{\texttt{ArchytasZX.\allowbreak{}PathSumCert.\allowbreak{}<name>}}. Both diagrams' path sums are reduced and canonicalized
(qubit or qudit; mixed bases raise ``\texttt{the two path sums have different bases}''); when the
canonical forms differ, are \textnormal{\texttt{none}}, or are zero it raises ``\texttt{the path sums do not reduce
to one canonical form}'' (qudit: ``\texttt{the qudit path sums ...}''). The file has, for each step
$k$ of the left trace, a comment \textnormal{\texttt{-{}-{} L step k: [rule] var v,\allowbreak{} pivot z,\allowbreak{} value .\allowbreak{}.\allowbreak{}.\allowbreak{},\allowbreak{} digits .\allowbreak{}.\allowbreak{}.\allowbreak{},\allowbreak{} fresh .\allowbreak{}.\allowbreak{}.\allowbreak{},\allowbreak{} bits .\allowbreak{}.\allowbreak{}.\allowbreak{},\allowbreak{} poly .\allowbreak{}.\allowbreak{}.\allowbreak{}}}
(present fields only) and a line \textnormal{\texttt{example := @<lemma>}} naming the rule lemma of
\cref{sec:lean-pathsum-rules0} by \textnormal{\texttt{\_\allowbreak{}LEMMA}} (\textnormal{\texttt{elim}}$\mapsto$\textnormal{\texttt{ps\_\allowbreak{}elim}},
\textnormal{\texttt{hh\_\allowbreak{}pivot}}$\mapsto$\textnormal{\texttt{ps\_\allowbreak{}hh\_\allowbreak{}pin}}, \textnormal{\texttt{hh\_\allowbreak{}gen}}$\mapsto$\textnormal{\texttt{ps\_\allowbreak{}hh\_\allowbreak{}gen}}, \textnormal{\texttt{rem}}$\mapsto$\textnormal{\texttt{ps\_\allowbreak{}rem}},
\textnormal{\texttt{hh\_\allowbreak{}constraint}}$\mapsto$\textnormal{\texttt{ps\_\allowbreak{}hh}}, \textnormal{\texttt{hh\_\allowbreak{}nl}}$\mapsto$\textnormal{\texttt{ps\_\allowbreak{}hh\_\allowbreak{}nl}},
\textnormal{\texttt{q\_\allowbreak{}elim}}$\mapsto$\textnormal{\texttt{q\_\allowbreak{}elim}}, \textnormal{\texttt{q\_\allowbreak{}hh\_\allowbreak{}pivot}} and \textnormal{\texttt{q\_\allowbreak{}hh\_\allowbreak{}constraint}}$\mapsto$\textnormal{\texttt{q\_\allowbreak{}hh}},
\textnormal{\texttt{q\_\allowbreak{}gauss}}$\mapsto$\textnormal{\texttt{q\_\allowbreak{}gauss}}); [omega] writes \textnormal{\texttt{example := And.\allowbreak{}intro ps\_\allowbreak{}omega0 ps\_\allowbreak{}omega1}}
(conjugate: \textnormal{\texttt{ps\_\allowbreak{}omega\_\allowbreak{}conj0}}, \textnormal{\texttt{ps\_\allowbreak{}omega\_\allowbreak{}conj1}}), [Dyadic] writes
\textnormal{\texttt{example (\allowbreak{}S : \ensuremath{\mathbb{N}}) := ps\_\allowbreak{}dyadic <k> S}}, and qudit rules pass $p$ (and the pivot for
\textnormal{\texttt{q\_\allowbreak{}gauss}}) as explicit arguments. Then \textnormal{\texttt{def canonL}}, \textnormal{\texttt{def canonR}} (as
\textnormal{\texttt{PSCanon}} or \textnormal{\texttt{QCanon}} literals over $N = \operatorname{lcm}$ of the two constants'
orders, constants as reduced rational coefficient lists), \textnormal{\texttt{theorem canon\_\allowbreak{}eq : canonL = canonR := rfl}},
and the right trace. The file is a record of the trace; it proves only the equality of the two
records.

\begin{signature}
\textnormal{\texttt{export\_\allowbreak{}proof(\allowbreak{}left: Diagram,\allowbreak{} right: Diagram,\allowbreak{} name: str = "pair") -{}> str}}
\end{signature}
\noindent Proves that the two qubit path sums have one value at every boundary assignment.
Non-qubit sums raise ``\texttt{export\_proof covers qubit path sums only}''. Sums with amplitude
atoms, or whose reduction takes a \textnormal{\texttt{quarter}}, \textnormal{\texttt{turn}}, \textnormal{\texttt{turn\_\allowbreak{}held}}, or \textnormal{\texttt{average\_\allowbreak{}ctx}}
step (\textnormal{\texttt{\_\allowbreak{}quartered}}), go to \textnormal{\texttt{\_\allowbreak{}export\_\allowbreak{}amp\_\allowbreak{}proof}}. Otherwise:

\begin{algorithm}[\textnormal{\texttt{\_\allowbreak{}plan(\allowbreak{}raw\_\allowbreak{}l,\allowbreak{} raw\_\allowbreak{}r,\allowbreak{} m\_\allowbreak{}min=1,\allowbreak{} support=False)}}]\label{alg:lean-pathsum-plan}
\begin{enumerate}
\item Raise when a side carries atoms (``\texttt{export\_proof has no amplitude atoms}''). Reduce
  both sides; raise unless both canonical forms exist, are equal (\textnormal{\texttt{same}}), and are nonzero.
\item The phase modulus $M$ is $\operatorname{lcm}(4, m_{\min})$, then the lcm with the
  denominator of every phase coefficient of both raw sums, then with $8$ when a trace takes
  [omega] or its conjugate.
\item Boolean elimination is needed when \textnormal{\texttt{eliminate}} fails on either reduced sum. Then both
  sides' \textnormal{\texttt{boolean\_\allowbreak{}eliminate}} must give the same leftover support $g$ and the same pivot order,
  and $g$ must be $1$ unless \textnormal{\texttt{support}} (else ``\texttt{the Boolean eliminations differ or leave a
  support}''). Otherwise the pivot order is \textnormal{\texttt{\_\allowbreak{}elimination\_\allowbreak{}order}} of the left constraints
  (\cref{alg:lean-pathsum-elimorder}).
\item Each side's step list is \textnormal{\texttt{\_\allowbreak{}steps}}: the trace steps translated by \textnormal{\texttt{trace\_\allowbreak{}step}}
  (\cref{tab:lean-pathsum-trace}) from the sorted raw path, followed by
  \textnormal{\texttt{.\allowbreak{}gather}}, \textnormal{\texttt{.\allowbreak{}bcanon v}} per pivot, and \textnormal{\texttt{.\allowbreak{}pdrop}} (Boolean) or \textnormal{\texttt{.\allowbreak{}canon v}} per
  pivot.
\end{enumerate}
\end{algorithm}

\begin{algorithm}[\textnormal{\texttt{\_\allowbreak{}elimination\_\allowbreak{}order(\allowbreak{}constraints,\allowbreak{} boundary)}}]\label{alg:lean-pathsum-elimorder}
Keep the nonzero constraints. Repeat while constraints remain and the previous pass made
progress: for each boundary variable $v$ not yet ordered, ascending, find the first constraint
containing the monomial $\{v\}$ and no other monomial containing $v$; when found, remove it,
substitute $v := (\text{constraint} - \{v\})$ in the remaining constraints, drop the ones that
become $0$, and append $v$ to the order.
\end{algorithm}

\begin{xltabular}{\textwidth}{@{}>{\raggedright\arraybackslash}p{3.2cm}>{\raggedright\arraybackslash}X>{\raggedright\arraybackslash}p{4.2cm}@{}}
\caption{\texttt{trace\_step}: the \texttt{PSem.Step} texts of one engine trace step on the current
path \texttt{cur}, and the path after it ($y$ the step's variable, $z$ its pivot).}\label{tab:lean-pathsum-trace}\\
\toprule Engine rule & Step texts & Path after\\\midrule\endfirsthead
\toprule Engine rule & Step texts & Path after\\\midrule\endhead
\textnormal{\texttt{elim}}, \textnormal{\texttt{hh\_\allowbreak{}constraint}}, \textnormal{\texttt{omega}}, \textnormal{\texttt{omega\_\allowbreak{}conj}} & \textnormal{\texttt{.\allowbreak{}perm [y,\allowbreak{} rest]}}, then \textnormal{\texttt{.\allowbreak{}elim}}, \textnormal{\texttt{.\allowbreak{}hhCons}}, \textnormal{\texttt{.\allowbreak{}omega}}, \textnormal{\texttt{.\allowbreak{}omegaConj}} & rest ($y$ removed)\\
\textnormal{\texttt{rem}} & \textnormal{\texttt{.\allowbreak{}perm [y,\allowbreak{} rest]}}, \textnormal{\texttt{.\allowbreak{}rem}} & $[y,\text{rest}]$\\
\textnormal{\texttt{bsub}} & \textnormal{\texttt{.\allowbreak{}bsub y}} & unchanged\\
\textnormal{\texttt{dyadic}} & \textnormal{\texttt{.\allowbreak{}perm [digits,\allowbreak{} rest]}}, \textnormal{\texttt{.\allowbreak{}dyadicB digits [poly] fresh bits}} & $[\text{fresh},\text{rest}]$\\
\textnormal{\texttt{hh\_\allowbreak{}pivot}} & \textnormal{\texttt{.\allowbreak{}perm [y,\allowbreak{} z,\allowbreak{} rest]}}, \textnormal{\texttt{.\allowbreak{}hhPivot}} & rest\\
\textnormal{\texttt{hh\_\allowbreak{}gen}}, \textnormal{\texttt{hh\_\allowbreak{}nl}} & \textnormal{\texttt{.\allowbreak{}perm [y,\allowbreak{} z,\allowbreak{} rest]}}, \textnormal{\texttt{.\allowbreak{}hhGen}} / \textnormal{\texttt{.\allowbreak{}hhNl}} & $[y,\text{rest}]$\\
other & raises ``\texttt{the rule ... has no PathSumRules step}'' & --\\
\bottomrule
\end{xltabular}

\noindent The [Dyadic] polynomial is printed as \textnormal{\texttt{(\allowbreak{}variables of the monomial,\allowbreak{} coefficient)}}
pairs, the digits \textnormal{\texttt{bits}} as $\mathbb F_2$ polynomials. The constant data are:
\textnormal{\texttt{\_\allowbreak{}fac\_\allowbreak{}model(\allowbreak{}steps,\allowbreak{} m)}} multiplies $2$ for \textnormal{\texttt{.\allowbreak{}elim}}, \textnormal{\texttt{.\allowbreak{}hhCons}}, \textnormal{\texttt{.\allowbreak{}hhPivot}},
\textnormal{\texttt{.\allowbreak{}hhNl}}, $\sqrt2\,\omega_m^{m/8}$ for \textnormal{\texttt{.\allowbreak{}omega}}, and $\sqrt2\,\omega_m^{7m/8}$ for
\textnormal{\texttt{.\allowbreak{}omegaConj}}; \textnormal{\texttt{\_\allowbreak{}const\_\allowbreak{}model(\allowbreak{}c,\allowbreak{} k)}} is the raw constant over $\omega_n$ times
$\sqrt2^{k}$ ($k$ the raw sum's power of $\sqrt2$); \textnormal{\texttt{\_\allowbreak{}const\_\allowbreak{}turn(\allowbreak{}ps,\allowbreak{} m,\allowbreak{} boolean)}} substitutes
the elimination rows (Boolean or plain) into the reduced phase and returns its constant term
times $m$ mod $m$; \textnormal{\texttt{\_\allowbreak{}leftover(\allowbreak{}red)}} is true when plain elimination fails and the Boolean
leftover support is not $1$. The emitted file is:
\begin{lstlisting}[style=lean]
import ArchytasZX.PathSumChain
namespace ArchytasZX.PathSumProof.<name>
open ArchytasZX ArchytasZX.PSem
set_option maxRecDepth 100000
set_option maxHeartbeats 0
def sL : Data := { path := ..., phase := ..., cons := ..., rows := [] }
def sR : Data := ...
def stepsL : List Step := [.front ..., .elim, ...]
def stepsR : List Step := [...]
theorem facL : facs stepsL = <fL> := by
  simp only [stepsL, facs, fac]
  rw [om_sub M 8 (M/8) (by norm_num) (by norm_num)]   -- only when M != 8 and omega occurs
  ring
theorem facR : ...
theorem consts : (<fL>) * <cL> * om M ^ kL = (<fR>) * <cR> * om M ^ kR := by
  have key : <lhs> = <rhs> :=
    <proof of the scalar renderer>
  linear_combination key
run_chain runL M stepsL sL 4
run_chain runR M stepsR sR 4
theorem main (σ : ℕ → ZMod 2) :
    PSem.val M <cL> sL σ = PSem.val M <cR> sR σ := by
  rw [run_sound M (by norm_num) stepsL sL _ σ runL_ok, runL_eq,
    run_sound M (by norm_num) stepsR sR _ σ runR_ok, runR_eq]
  apply final_sound M (by norm_num) _ _ _ _ σ kL kR (by decide) (by decide)
    (by decide)
  simpa only [facL, facR] using consts
end ArchytasZX.PathSumProof.<name>
\end{lstlisting}
with \textnormal{\texttt{finalS\_\allowbreak{}sound}} in place of \textnormal{\texttt{final\_\allowbreak{}sound}} when \textnormal{\texttt{\_\allowbreak{}leftover}} holds for the left
side, a module docstring between the import and the namespace, and blank lines between
declarations.

\paragraph{Amplitude-atom proofs.} \textnormal{\texttt{\_\allowbreak{}export\_\allowbreak{}amp\_\allowbreak{}proof(\allowbreak{}raw\_\allowbreak{}l,\allowbreak{} raw\_\allowbreak{}r,\allowbreak{} name)}} takes
$M = \textnormal{\texttt{amp\_\allowbreak{}modulus}}$ (a multiple of $8$ and of every phase denominator), a shared
\textnormal{\texttt{Coefs(\allowbreak{}M)}}, and each side's \textnormal{\texttt{initial}} run (\cref{sec:lean-pathsum-python-amp}), and calls
\textnormal{\texttt{amp\_\allowbreak{}plan}}:
\begin{algorithm}[\textnormal{\texttt{amp\_\allowbreak{}plan(\allowbreak{}sides)}}]\label{alg:lean-pathsum-ampplan}
\begin{enumerate}
\item Reduce both sides; raise unless both canonical forms exist and are nonzero. When they
  differ (\emph{atom form}), their rows, supports, and boundaries must agree.
\item Compute the elimination: plain (\textnormal{\texttt{\_\allowbreak{}elimination\_\allowbreak{}order}}, steps \textnormal{\texttt{.\allowbreak{}canon v}}) or Boolean
  (equal supports and pivot orders required, steps \textnormal{\texttt{.\allowbreak{}gather}}, \textnormal{\texttt{.\allowbreak{}bcanon v}}, \textnormal{\texttt{.\allowbreak{}pdrop}}).
\item Try the modes (drop, split) in the order $(\mathrm{F},\mathrm{F})$, and, for Boolean
  elimination with support $g\ne1$, $(\mathrm{T},\mathrm{F})$, $(\mathrm{T},\mathrm{T})$. In each
  mode and for each side: set the run's support to $g$ when dropping; record the start (inner
  path, held, atoms); \textnormal{\texttt{reduce\_\allowbreak{}run}} along the trace; add the canon steps as inner steps;
  \textnormal{\texttt{substitute}} the rows into the atoms; in atom form \textnormal{\texttt{normalize}} the canonical phase with
  splitting per the mode. Stop at the first mode whose two final atom multisets (\textnormal{\texttt{final\_\allowbreak{}key}})
  are equal; between modes restore both runs and truncate the coefficient list. When no mode
  matches, raise ``\texttt{the reduced sums carry different amplitude atoms}''.
\end{enumerate}
\end{algorithm}
\noindent The file defines \textnormal{\texttt{A}} (\textnormal{\texttt{Coefs.\allowbreak{}lean\_\allowbreak{}def}}: \textnormal{\texttt{| i => <base i over om M>}}, then
\textnormal{\texttt{| \_\allowbreak{} => 0}}), \textnormal{\texttt{sL}}, \textnormal{\texttt{tL : AData := \ensuremath{\langle}sL,\allowbreak{} held,\allowbreak{} atoms\ensuremath{\rangle}}}, \textnormal{\texttt{stepsL : List AStep}}, and the
same for R; the theorems \textnormal{\texttt{facL}}/\textnormal{\texttt{facR}} (\textnormal{\texttt{afac\_\allowbreak{}lines}}:
\textnormal{\texttt{simp only [stepsL,\allowbreak{} afacs,\allowbreak{} AStep.\allowbreak{}fac,\allowbreak{} fac,\allowbreak{} A]}}, the \textnormal{\texttt{om\_\allowbreak{}sub}} rewrite when an [omega]
step occurs and $M\ne8$, \textnormal{\texttt{all\_\allowbreak{}goals ring}}); \textnormal{\texttt{reqL<i>}} and \textnormal{\texttt{reqsL}} (\textnormal{\texttt{reqs\_\allowbreak{}lines}}),
\textnormal{\texttt{qreqL<j>e<e>b<b>}} and \textnormal{\texttt{qreqsL}} (\textnormal{\texttt{qreqs\_\allowbreak{}lines}}) and their R versions; \textnormal{\texttt{consts}};
\textnormal{\texttt{aruns\_\allowbreak{}chain arunsL M stepsL tL 4}} and the R command; and
\begin{lstlisting}[style=lean]
theorem main (σ : ℕ → ZMod 2) :
    valA M <cL> A tL σ = valA M <cR> A tR σ := by
  rw [aruns_sound M (by norm_num) A stepsL tL _ σ reqsL qreqsL arunsL_ok, arunsL_eq,
    aruns_sound M (by norm_num) A stepsR tR _ σ reqsR qreqsR arunsR_ok, arunsR_eq]
  apply afinal_perm M (by norm_num) _ _ A _ _ σ kL kR (by decide) (by decide)
    (by decide) (by decide) (by decide) (by decide)
  simpa only [facL, facR] using consts
\end{lstlisting}
(\textnormal{\texttt{afinalS\_\allowbreak{}perm}} when a support other than $1$ is left). Each \textnormal{\texttt{reqL<i>}} states
\textnormal{\texttt{reqHolds A (\allowbreak{}k,\allowbreak{} j,\allowbreak{} [[.\allowbreak{}.\allowbreak{}.\allowbreak{}],\allowbreak{} .\allowbreak{}.\allowbreak{}.\allowbreak{}])}} and is proved by the scalar identity
\textnormal{\texttt{have key : <k A\_\allowbreak{}j> = <\ensuremath{\Sigma} \ensuremath{\Pi} A\_\allowbreak{}i>}}, \textnormal{\texttt{simp only [reqHolds,\allowbreak{} A,\allowbreak{} List.\allowbreak{}map\_\allowbreak{}cons,\allowbreak{} .\allowbreak{}.\allowbreak{}.\allowbreak{}]}},
\textnormal{\texttt{all\_\allowbreak{}goals push\_\allowbreak{}cast}}, \textnormal{\texttt{all\_\allowbreak{}goals linear\_\allowbreak{}combination key}}; \textnormal{\texttt{reqsL}} proves
\textnormal{\texttt{\ensuremath{\forall} r \ensuremath{\in} areqs stepsL,\allowbreak{} reqHolds A r}} by computing \textnormal{\texttt{areqs stepsL}} with \textnormal{\texttt{decide}},
splitting membership, and \textnormal{\texttt{exacts}} the individual lemmas. \textnormal{\texttt{qreqsL}} does the same for
\textnormal{\texttt{aqreqs M stepsL}} and the root requirements, with one lemma per distinct triple.

\begin{signature}
\textnormal{\texttt{export\_\allowbreak{}qudit\_\allowbreak{}proof(\allowbreak{}left: Diagram,\allowbreak{} right: Diagram,\allowbreak{} name: str = "pair") -{}> str}}
\end{signature}
\noindent Requires two \textnormal{\texttt{QuditPathSum}}s over one odd prime $p$ with equal boundaries and no
atoms; reduces and canonicalizes; requires equal \textnormal{\texttt{\_\allowbreak{}qudit\_\allowbreak{}order}}s (the pivots of
\textnormal{\texttt{eliminate\_\allowbreak{}qudit}}, computed as \cref{alg:lean-pathsum-elimorder} over $\mathbb F_p$).
\textnormal{\texttt{\_\allowbreak{}qudit\_\allowbreak{}steps}} translates \textnormal{\texttt{q\_\allowbreak{}elim}} and \textnormal{\texttt{q\_\allowbreak{}hh\_\allowbreak{}constraint}} to
\textnormal{\texttt{.\allowbreak{}perm [y,\allowbreak{} rest]}}, \textnormal{\texttt{.\allowbreak{}elim}}/\textnormal{\texttt{.\allowbreak{}hhCons}}; \textnormal{\texttt{q\_\allowbreak{}gauss}} to \textnormal{\texttt{.\allowbreak{}perm [y,\allowbreak{} rest]}},
\textnormal{\texttt{.\allowbreak{}gauss a}} ($a$ the step's pivot field); \textnormal{\texttt{q\_\allowbreak{}hh\_\allowbreak{}pivot}} to \textnormal{\texttt{.\allowbreak{}perm [y,\allowbreak{} z,\allowbreak{} rest]}},
\textnormal{\texttt{.\allowbreak{}hhPivot}}; then \textnormal{\texttt{.\allowbreak{}canon v}} per pivot and \textnormal{\texttt{.\allowbreak{}norm}}. The constant model multiplies
$p$ per \textnormal{\texttt{.\allowbreak{}elim}}/\textnormal{\texttt{.\allowbreak{}hhCons}}/\textnormal{\texttt{.\allowbreak{}hhPivot}} and, per \textnormal{\texttt{.\allowbreak{}gauss a}}, the Legendre form
$\left(\frac ap\right)\sum_{k=1}^{p-1}\left(\frac kp\right)\omega_n^{kn/p}$ over
$n=\operatorname{lcm}(p,\text{the constants' orders})$. The file imports \file{ArchytasZX.GaussSign}
and \file{ArchytasZX.PathSumQuditChain}, opens \textnormal{\texttt{ArchytasZX.\allowbreak{}QSem}}, declares
\textnormal{\texttt{instance : Fact (\allowbreak{}Nat.\allowbreak{}Prime p) := \ensuremath{\langle}by norm\_\allowbreak{}num\ensuremath{\rangle}}}, the literals \textnormal{\texttt{sL}}, \textnormal{\texttt{sR}}, \textnormal{\texttt{stepsL}},
\textnormal{\texttt{stepsR}}, one \textnormal{\texttt{theorem lg<k> : lg p k = v := by decide}} per residue $k<p$, one
\textnormal{\texttt{theorem qg<a> : qg p a = <Legendre form>}} per Gauss coefficient (by \textnormal{\texttt{qg\_\allowbreak{}legendre}},
\textnormal{\texttt{Finset.\allowbreak{}sum\_\allowbreak{}range\_\allowbreak{}succ}}, the \textnormal{\texttt{lg}} values, the \textnormal{\texttt{om\_\allowbreak{}sub}} rewrite when $n\ne p$,
\textnormal{\texttt{push\_\allowbreak{}cast}}, \textnormal{\texttt{ring}}), \textnormal{\texttt{facL}}, \textnormal{\texttt{facR}}, \textnormal{\texttt{consts}}, the two
\textnormal{\texttt{qrun\_\allowbreak{}chain qrunL p stepsL sL 4}} commands, and \textnormal{\texttt{main (\allowbreak{}\ensuremath{\sigma} : \ensuremath{\mathbb{N}} \ensuremath{\to} ZMod p)}} closing with
\textnormal{\texttt{QSem.\allowbreak{}final\_\allowbreak{}sound}}.

\paragraph{Networks.} \textnormal{\texttt{export\_\allowbreak{}net(\allowbreak{}left,\allowbreak{} right,\allowbreak{} name="main")}} proves \textnormal{\texttt{Equal E sL NL sR NR}}
for two diagrams of Z, X, F, B, and S nodes on dimensions $2^k$ and W nodes on dimension $2$.
It dispatches sums with atoms to \textnormal{\texttt{\_\allowbreak{}export\_\allowbreak{}amp\_\allowbreak{}net}} and quartered pairs to
\textnormal{\texttt{\_\allowbreak{}export\_\allowbreak{}quarter\_\allowbreak{}net}}; otherwise it builds \textnormal{\texttt{net\_\allowbreak{}body}} and writes a file importing
\file{ArchytasZX.PathSumNetW} (and \file{ArchytasZX.PathSumSupport} when the body uses
\textnormal{\texttt{finalS\_\allowbreak{}sound}}), in namespace \textnormal{\texttt{ArchytasZX.\allowbreak{}PathSumNetProof}}, with \textnormal{\texttt{NET\_\allowbreak{}OPTIONS}},
\textnormal{\texttt{net\_\allowbreak{}env}}, and the body.

\begin{algorithm}[\textnormal{\texttt{net\_\allowbreak{}body}}]
Signature \textnormal{\texttt{net\_\allowbreak{}body(\allowbreak{}left,\allowbreak{} right,\allowbreak{} name,\allowbreak{} phases,\allowbreak{} m\_\allowbreak{}min=1,\allowbreak{} coefs=None,\allowbreak{} prefer=None) -{}> NetBody}}.
\begin{enumerate}
\item Compute both path sums with \textnormal{\texttt{from\_\allowbreak{}diagram(\allowbreak{}\ensuremath{\cdot},\allowbreak{} w\_\allowbreak{}nodes=True)}}; require qubit sums with equal
  boundaries; raise $m_{\min}$ to the lcm with \textnormal{\texttt{\_\allowbreak{}table\_\allowbreak{}modulus}} of both diagrams (the lcm of the
  denominators of the Z and X phase tables in turns on dimensions $2^k$).
\item Run \cref{alg:lean-pathsum-plan} with \textnormal{\texttt{support=True}}; when it raises, try
  \textnormal{\texttt{\_\allowbreak{}zero\_\allowbreak{}plan}} (each side zero by its scalar or reducing to a constraint $1=0$) and emit
  \textnormal{\texttt{\_\allowbreak{}zero\_\allowbreak{}body}}, whose sides end with \textnormal{\texttt{val\_\allowbreak{}of\_\allowbreak{}zeroOK}}; else re-raise.
\item Build each side's network by \textnormal{\texttt{\_\allowbreak{}net\_\allowbreak{}side}} with label bases $2^{16}$ (left) and $2^{17}$
  (right): boundary port $i$ is label $i$, summed wire $w$ is label $\text{base}+w$, digit $j$ of a
  summed wire is the variable of its class when that is a path variable not yet taken, else a
  fresh variable from \textnormal{\texttt{PIN\_\allowbreak{}OFFSET}}; each Z spider's digit $j$ is anchored at the leg through
  which a search from the class's variable first reaches it; pin steps substitute the other
  summed digits deepest first. Phase vectors are collected into \textnormal{\texttt{phases}}.
\item Require equal boundary digit counts; append each side's reduction steps to its pin steps;
  compute the factor models, constant turns, and the scalar identity
  $f_L\,s_L\,w_L\,(1/\sqrt2)^{e_L}\,\omega_M^{k_L} = f_R\,s_R\,w_R\,(1/\sqrt2)^{e_R}\,\omega_M^{k_R}$
  ($s$ the diagram scalars, $w$ the weighted spiders' base coefficients, $e$ the fragments'
  powers of $1/\sqrt2$).
\item Emit per side (\textnormal{\texttt{\_\allowbreak{}net\_\allowbreak{}lines}}): \textnormal{\texttt{ownL : Own}}, \textnormal{\texttt{labsL}}, \textnormal{\texttt{digL}}, the theorem
  \textnormal{\texttt{ownLOK}} by \textnormal{\texttt{decide}}, \textnormal{\texttt{NL : Net}}, \textnormal{\texttt{psL : List (\allowbreak{}Node \ensuremath{\times} Frag)}}, the theorem
  \textnormal{\texttt{netL : NetPSD E NL M ownL labsL digL psL}} (fields \textnormal{\texttt{ownok}}, \textnormal{\texttt{nodes := rfl}},
  \textnormal{\texttt{node}} by case split on the pairs with one fragment theorem per node, \textnormal{\texttt{disj}},
  \textnormal{\texttt{dims}}, \textnormal{\texttt{cover}}, \textnormal{\texttt{sub}} by \textnormal{\texttt{decide}}), and \textnormal{\texttt{cL}} (the joined constant by
  \textnormal{\texttt{simp only [.\allowbreak{}.\allowbreak{}.\allowbreak{}]}} and \textnormal{\texttt{ring}}); then \textnormal{\texttt{stepsL}}, \textnormal{\texttt{stepsR}}, \textnormal{\texttt{facL}},
  \textnormal{\texttt{facR}}, \textnormal{\texttt{consts}}, two \textnormal{\texttt{run\_\allowbreak{}chain}} commands on
  \textnormal{\texttt{(\allowbreak{}netDataD digL NL (\allowbreak{}psL.\allowbreak{}map Prod.\allowbreak{}snd))}}, and the theorem \textnormal{\texttt{<name>}} proved by
  \textnormal{\texttt{equal\_\allowbreak{}psd2 \_\allowbreak{} \_\allowbreak{} netL netR rfl}}, \textnormal{\texttt{cL}}, \textnormal{\texttt{cR}}, \textnormal{\texttt{val\_\allowbreak{}smul}}, \textnormal{\texttt{run\_\allowbreak{}sound}},
  \textnormal{\texttt{val\_\allowbreak{}agree M \_\allowbreak{} runL [boundary digits] \_\allowbreak{} \_\allowbreak{} (\allowbreak{}by decide) (\allowbreak{}by decide) (\allowbreak{}dbits\_\allowbreak{}agree ownL ownR \_\allowbreak{} [.\allowbreak{}.\allowbreak{}.\allowbreak{}] (\allowbreak{}by decide))}},
  and \textnormal{\texttt{final\_\allowbreak{}sound}} or \textnormal{\texttt{finalS\_\allowbreak{}sound}}.
\end{enumerate}
\end{algorithm}

\noindent \textnormal{\texttt{net\_\allowbreak{}env(\allowbreak{}phases,\allowbreak{} labs,\allowbreak{} dims,\allowbreak{} m,\allowbreak{} coefs=None,\allowbreak{} tabs=None)}} emits the shared environment:
\textnormal{\texttt{def \ensuremath{\Delta} (\allowbreak{}k : \ensuremath{\mathbb{N}}) : \ensuremath{\mathbb{N}} := <dims>}} (\textnormal{\texttt{2 \^{} k}} for qubit digits), \textnormal{\texttt{\ensuremath{\Phi} : \ensuremath{\mathbb{N}} \ensuremath{\to} \ensuremath{\mathbb{N}} \ensuremath{\to} \ensuremath{\mathbb{R}}}} with
\textnormal{\texttt{| i+1 => fun v => <body> / (\allowbreak{}(\allowbreak{}m : \ensuremath{\mathbb{N}}) : \ensuremath{\mathbb{R}})}} per collected phase vector and the zero vector
otherwise, \textnormal{\texttt{labs : List (\allowbreak{}\ensuremath{\mathbb{N}} \ensuremath{\times} \ensuremath{\mathbb{N}})}} (label, digit count), and, when \textnormal{\texttt{coefs}} is nonempty,
\textnormal{\texttt{\ensuremath{\Psi} : \ensuremath{\mathbb{N}} \ensuremath{\to} \ensuremath{\mathbb{N}} \ensuremath{\to} \ensuremath{\mathbb{C}}}} with \textnormal{\texttt{| i => fun v => C * ex (\allowbreak{}\ensuremath{\Phi} i v)}} (times
\textnormal{\texttt{(\allowbreak{}(\allowbreak{}tbl.\allowbreak{}getD v []).\allowbreak{}map A).\allowbreak{}prod}} at an index of \textnormal{\texttt{tabs}}) and \textnormal{\texttt{| \_\allowbreak{} => fun \_\allowbreak{} => 1}}; finally
\textnormal{\texttt{E : Env := \{ dimv := \ensuremath{\Delta},\allowbreak{} phase := \ensuremath{\Phi},\allowbreak{} labs := labs,\allowbreak{} coef := \ensuremath{\Psi} \}}} (without \textnormal{\texttt{coef}} when
there are no coefficients). \textnormal{\texttt{NetBody}} carries the lines, the labels' digit counts, the two
scalar texts, the modulus, an optional head, and the atom tables.

\textnormal{\texttt{net\_\allowbreak{}amp\_\allowbreak{}body}} and \textnormal{\texttt{\_\allowbreak{}export\_\allowbreak{}amp\_\allowbreak{}net}} are the versions with atoms: weighted spiders whose
weights are $0$ or $\sqrt2$ times a common magnitude take atoms on their anchor digits (Z with
legs, by \textnormal{\texttt{za\_\allowbreak{}frag}}) or their own digits (X and legless Z, by \textnormal{\texttt{xa\_\allowbreak{}frag}}); every summed
digit and fragment path variable is held (\textnormal{\texttt{NetPSDA}}, \textnormal{\texttt{netDataA}}), \textnormal{\texttt{\_\allowbreak{}net\_\allowbreak{}run}} replays the
pins as atom substitutions, and the main theorem uses \textnormal{\texttt{equal\_\allowbreak{}psda2}}, \textnormal{\texttt{aruns\_\allowbreak{}sound}},
\textnormal{\texttt{valA\_\allowbreak{}agree}}, and \textnormal{\texttt{afinal\_\allowbreak{}perm}}. \textnormal{\texttt{net\_\allowbreak{}quarter\_\allowbreak{}body}} and \textnormal{\texttt{\_\allowbreak{}export\_\allowbreak{}quarter\_\allowbreak{}net}}
run each side's pin steps on its network sum with \textnormal{\texttt{run\_\allowbreak{}chain}} (block \textnormal{\texttt{CHAIN\_\allowbreak{}BLOCK}}), then
replay the reduction with atoms from \textnormal{\texttt{\ensuremath{\langle}runL,\allowbreak{} [],\allowbreak{} []\ensuremath{\rangle}}} with \textnormal{\texttt{aruns\_\allowbreak{}chain}} (block
\textnormal{\texttt{CHAIN\_\allowbreak{}BLOCK}}), linking the two by \textnormal{\texttt{valA\_\allowbreak{}nil}}.

\paragraph{Qudit networks.} \textnormal{\texttt{export\_\allowbreak{}qudit\_\allowbreak{}net(\allowbreak{}left,\allowbreak{} right,\allowbreak{} name="main")}} covers Z, X, and F
nodes on one odd prime dimension $p$ (file importing \file{ArchytasZX.PathSumQuditNet},
\textnormal{\texttt{QNET\_\allowbreak{}OPTIONS}}): \textnormal{\texttt{\_\allowbreak{}qnet\_\allowbreak{}side}} labels boundary port $i$ by $i$ and a summed wire by its
class's variable when free, else from \textnormal{\texttt{QPIN\_\allowbreak{}OFFSET}}, anchors each Z spider, writes pin steps
setting the other summed labels to their anchors deepest first, then the reduction trace;
\textnormal{\texttt{\_\allowbreak{}qnet\_\allowbreak{}lines}} writes \textnormal{\texttt{NL}}, \textnormal{\texttt{psL : List (\allowbreak{}Node \ensuremath{\times} QFrag)}}, and
\textnormal{\texttt{netL : QNetPS E NL p psL}}; \textnormal{\texttt{qnet\_\allowbreak{}body}} closes with \textnormal{\texttt{equal\_\allowbreak{}qps}}, the Legendre forms, and
\textnormal{\texttt{QSem.\allowbreak{}final\_\allowbreak{}sound}}. It raises \textnormal{\texttt{ScalarUnsupported}} when the constants' identity has no
proof.

\subsection{\texorpdfstring{\module{archytaszx.semantics.lean_amp}}{lean\_amp}}\label{sec:lean-pathsum-python-amp}

\textnormal{\texttt{Atom = tuple[int,\allowbreak{} F2Poly]}} (coefficient index, polynomial). \textnormal{\texttt{f2\_\allowbreak{}text(\allowbreak{}poly)}} prints an
$\mathbb F_2$ polynomial as a Lean list of sorted monomials sorted by their lists.

\begin{signature}
\textnormal{\texttt{@dataclass class Coefs: m: int; values: list[Cyclo]}}
\end{signature}
\noindent The atoms' bases by index over $\omega_m$. \textnormal{\texttt{index(\allowbreak{}c)}} returns the index of an equal
base, appending new ones; \textnormal{\texttt{model(\allowbreak{}i)}} is base $i$ over \textnormal{\texttt{Om(\allowbreak{}m)}} (raising when $m$ is not a
multiple of the base's order); \textnormal{\texttt{lean\_\allowbreak{}def(\allowbreak{}sc)}} is the definition of \textnormal{\texttt{A}}.

\begin{signature}
\textnormal{\texttt{@dataclass class AmpRun: coefs,\allowbreak{} cur,\allowbreak{} held,\allowbreak{} atoms,\allowbreak{} steps,\allowbreak{} reqs,\allowbreak{} facs,\allowbreak{} omega,\allowbreak{} qreqs,\allowbreak{} support}}
\end{signature}
\noindent The Lean state of a path sum with atoms: inner path \textnormal{\texttt{cur}}, held variables, atoms,
the \textnormal{\texttt{AStep}} texts written so far, the coefficient requirements \textnormal{\texttt{(\allowbreak{}k,\allowbreak{} j,\allowbreak{} ps)}}, the factor
models, whether an [omega] step occurred, the root requirements \textnormal{\texttt{(\allowbreak{}e,\allowbreak{} j,\allowbreak{} b)}}, and the leftover
support. \textnormal{\texttt{initial(\allowbreak{}coefs,\allowbreak{} path,\allowbreak{} amps)}} holds the path variables an atom reads (ascending) and
puts the others on the inner path (ascending). \textnormal{\texttt{snapshot}} and \textnormal{\texttt{restore}} copy and restore the
state (coefficients shared). The methods write the following steps.
\begin{itemize}
\item \textnormal{\texttt{inner(\allowbreak{}texts)}}: \textnormal{\texttt{.\allowbreak{}inner (\allowbreak{}st)}} per text, with factor $2$ for \textnormal{\texttt{.\allowbreak{}elim}},
  \textnormal{\texttt{.\allowbreak{}hhCons}}, \textnormal{\texttt{.\allowbreak{}hhPivot}}, \textnormal{\texttt{.\allowbreak{}hhNl}} and the [omega] factors; \textnormal{\texttt{trace(\allowbreak{}step)}} uses
  \textnormal{\texttt{trace\_\allowbreak{}step}} of \textnormal{\texttt{lean\_\allowbreak{}pathsum}}.
\item \textnormal{\texttt{release(\allowbreak{})}}: for each held variable no atom reads, \textnormal{\texttt{.\allowbreak{}hperm [y,\allowbreak{} rest]}} when it is not
  the head, then \textnormal{\texttt{.\allowbreak{}release}}.
\item \textnormal{\texttt{fold(\allowbreak{}path,\allowbreak{} unit=False)}}: \textnormal{\texttt{.\allowbreak{}adrop}} when some atom is on $0$; then in list order an atom
  on $1$ gives \textnormal{\texttt{.\allowbreak{}apick i}}, \textnormal{\texttt{.\allowbreak{}aconst j}} (factor $A_j$); with \textnormal{\texttt{unit}} an atom of base $1$
  gives \textnormal{\texttt{.\allowbreak{}aone j}} (requirement $(1,j,[[]])$); an atom of base $0$ free of \textnormal{\texttt{path}} gives
  \textnormal{\texttt{.\allowbreak{}azero j}} (requirement $(1,j,[])$).
\item \textnormal{\texttt{average(\allowbreak{}y,\allowbreak{} guard)}}: requires every atom reading $y$ to be on $(1+y)g$ or $yg$; the new base
  is $(\alpha+\beta)/2$; writes \textnormal{\texttt{.\allowbreak{}hperm}}, \textnormal{\texttt{.\allowbreak{}average y g j alpha beta}}, requirement
  $(2,j,[\alpha,\beta])$, factor $2$, and replaces the atoms reading $y$ by $(j,g)$.
\item \textnormal{\texttt{average\_\allowbreak{}half(\allowbreak{}y,\allowbreak{} q)}}: atoms on $1+y$ or $y$; bases $(\alpha+\beta)/2$, $(\alpha-\beta)/2$, and
  $-1$; writes \textnormal{\texttt{.\allowbreak{}inner (\allowbreak{}.\allowbreak{}norm)}}, \textnormal{\texttt{.\allowbreak{}averageH y j0 j1 jn alpha beta}} and the three
  requirements; atoms $(j_0, 1+q)$, $(j_1, q)$.
\item \textnormal{\texttt{quarter(\allowbreak{}y,\allowbreak{} b0,\allowbreak{} b1)}}: requires $4\mid m$; writes \textnormal{\texttt{.\allowbreak{}inner (\allowbreak{}.\allowbreak{}perm [y,\allowbreak{} rest])}},
  \textnormal{\texttt{.\allowbreak{}quarter b0 b1 j}} with $j$ the index of \textnormal{\texttt{QUARTER\_\allowbreak{}BASE}} ($(1+i)/2$), root requirement
  $(m/4,j,1)$, factor $2$, atom $(j,b_0)$.
\item \textnormal{\texttt{turn(\allowbreak{}y,\allowbreak{} bs)}}: requires $2^{|bs|}\mid m$; pairs $(s, \text{index of } \textnormal{\texttt{turn\_\allowbreak{}base(\allowbreak{}s,\allowbreak{} k)}})$
  for $s$ in \textnormal{\texttt{turn\_\allowbreak{}values(\allowbreak{}k)}}; writes \textnormal{\texttt{.\allowbreak{}turn [bs] [(\allowbreak{}s,\allowbreak{} j),\allowbreak{} .\allowbreak{}.\allowbreak{}.\allowbreak{}]}}, root requirements
  $(s\,m/2^k, j, 1)$, factor $2$, and the atoms $(j,\textnormal{\texttt{turn\_\allowbreak{}indicator(\allowbreak{}bs,\allowbreak{} s)}})$.
\item \textnormal{\texttt{hold(\allowbreak{}vs)}}, \textnormal{\texttt{turn\_\allowbreak{}held(\allowbreak{}y,\allowbreak{} bs)}}: hold the inner variables the digits read
  (\textnormal{\texttt{.\allowbreak{}inner (\allowbreak{}.\allowbreak{}perm .\allowbreak{}.\allowbreak{}.\allowbreak{})}}, \textnormal{\texttt{.\allowbreak{}hold}}), then \textnormal{\texttt{turn}}; when the constraint $N\ne2^{k-1}$ reads a
  live variable, \textnormal{\texttt{.\allowbreak{}acons j}} with $j$ the index of base $0$ (requirement $(1,j,[])$).
\item \textnormal{\texttt{average\_\allowbreak{}ctx(\allowbreak{}y,\allowbreak{} q,\allowbreak{} ctx)}}: per atom reading $y$ the values $n<2^{|ctx|}$ where it is $1$ at
  $y=0$ and at $y=1$ (\textnormal{\texttt{ctx\_\allowbreak{}value}}); the values $Q_s$ of $q$; per $n$ a new base
  $(\prod_{p_0}+\prod_{p_1})/2$; writes \textnormal{\texttt{.\allowbreak{}avctx y [ctx] [(\allowbreak{}j,\allowbreak{} Z0,\allowbreak{} Z1),\allowbreak{} .\allowbreak{}.\allowbreak{}.\allowbreak{}] Qs js jn}} with the
  requirements of \textnormal{\texttt{avctxReqs}}.
\item \textnormal{\texttt{pin(\allowbreak{}y,\allowbreak{} value)}}: \textnormal{\texttt{.\allowbreak{}inner (\allowbreak{}.\allowbreak{}norm)}}; when an atom reads $y$, hold the inner variables of
  the value, write \textnormal{\texttt{.\allowbreak{}asubB y value}}, substitute, and release; then
  \textnormal{\texttt{.\allowbreak{}inner (\allowbreak{}.\allowbreak{}perm [y,\allowbreak{} rest])}}, \textnormal{\texttt{.\allowbreak{}inner (\allowbreak{}.\allowbreak{}pin value)}}.
\item \textnormal{\texttt{substitute(\allowbreak{}rows)}}: \textnormal{\texttt{.\allowbreak{}asub v}} per row pivot an atom reads; \textnormal{\texttt{.\allowbreak{}adropC}} when the
  support removes monomials; \textnormal{\texttt{.\allowbreak{}anorm}}; \textnormal{\texttt{fold(\allowbreak{}set(\allowbreak{}),\allowbreak{} True)}}; then repeatedly merge two atoms
  on one polynomial (\textnormal{\texttt{.\allowbreak{}apick}}, \textnormal{\texttt{.\allowbreak{}amerge j j1 j2}}, requirement $(1,j,[[j_1,j_2]])$);
  finally \textnormal{\texttt{complement(\allowbreak{})}}.
\item \textnormal{\texttt{complement(\allowbreak{})}}: two atoms of one index on polynomials adding to $1$ give \textnormal{\texttt{.\allowbreak{}acomp j}}
  (factor $A_j$), repeated.
\item \textnormal{\texttt{split(\allowbreak{})}}: on a support, an atom whose monomials multiply pairwise to $0$ gives
  \textnormal{\texttt{.\allowbreak{}asplit}}, then \textnormal{\texttt{substitute(\allowbreak{}[])}}.
\item \textnormal{\texttt{normalize(\allowbreak{}phase,\allowbreak{} support,\allowbreak{} split=False)}}: with the support's lift $a+ab$, \textnormal{\texttt{.\allowbreak{}psup a b}}
  and \textnormal{\texttt{.\allowbreak{}inner (\allowbreak{}.\allowbreak{}norm)}}; each nonconstant phase term $(m,c)$ gives \textnormal{\texttt{.\allowbreak{}aphase m e j}}
  ($e=\lfloor cm\rfloor\bmod m$, $j$ the index of $e^{2\pi ic}$, root requirement $(e,j,0)$); each
  atom on a polynomial with a constant term other than $1$ gives \textnormal{\texttt{.\allowbreak{}aflip j inv one}}
  (requirements $(1,\mathit{one},[[]])$, $(1,\mathit{one},[[j,\mathit{inv}]])$, factor $A_j$); then
  \textnormal{\texttt{substitute(\allowbreak{}[])}} and, with \textnormal{\texttt{split}}, \textnormal{\texttt{split(\allowbreak{})}}.
\end{itemize}
\textnormal{\texttt{reduce\_\allowbreak{}run(\allowbreak{}run,\allowbreak{} trace)}} folds and releases, then dispatches each trace step: \textnormal{\texttt{average}},
\textnormal{\texttt{average\_\allowbreak{}half}}, \textnormal{\texttt{pin}}, \textnormal{\texttt{quarter}}, \textnormal{\texttt{turn}}, \textnormal{\texttt{turn\_\allowbreak{}held}}, \textnormal{\texttt{average\_\allowbreak{}ctx}} to the
methods above, every other rule to \textnormal{\texttt{trace}}; after each step it folds (except after inner trace
steps) and releases. It raises when held variables remain at the end.

\subsection{\texorpdfstring{\module{archytaszx.semantics.lean_qdigits}}{lean\_qdigits}}

\textnormal{\texttt{export\_\allowbreak{}qudit\_\allowbreak{}digits\_\allowbreak{}net(\allowbreak{}left,\allowbreak{} right,\allowbreak{} name="main")}} covers Z, X, F, binder, and splitter nodes on
the powers of one odd prime $p$ with some leg of dimension $p^k$, $k>1$. A label of dimension $p^k$
is $k$ digit variables, a label of dimension $p$ its own digit variable; a summed wire of one digit
is labeled by that digit (the class's variable when free, else from \textnormal{\texttt{QPIN\_\allowbreak{}OFFSET}}), a summed
wire of $k>1$ digits by \textnormal{\texttt{LABEL\_\allowbreak{}OFFSET + w}} (\textnormal{\texttt{LABEL\_\allowbreak{}OFFSET = 1 << 19}}). Z spiders on $p^k$,
binders, and splitters become \textnormal{\texttt{zFragQD}} and \textnormal{\texttt{bFragQ}}; Z, X, and F on dimension $p$ use
the fragments of \file{PathSumQuditNet} through \textnormal{\texttt{lift\_\allowbreak{}q}}. Per side it writes
\textnormal{\texttt{ownL : PSem.\allowbreak{}Own}}, \textnormal{\texttt{labsL}}, \textnormal{\texttt{digL}}, \textnormal{\texttt{ownLOK}}, \textnormal{\texttt{NL}},
\textnormal{\texttt{psL : List (\allowbreak{}Node \ensuremath{\times} QFrag)}}, \textnormal{\texttt{netL : QNetPSD E NL p ownL labsL digL psL}}, and \textnormal{\texttt{cL}};
then the steps, \textnormal{\texttt{lg<k>}}, \textnormal{\texttt{qg<a>}}, \textnormal{\texttt{facL}}/\textnormal{\texttt{facR}}, \textnormal{\texttt{consts}}, and the theorem
proved by \textnormal{\texttt{equal\_\allowbreak{}qpsd2 \_\allowbreak{} \_\allowbreak{} netL netR [boundary variables] (\allowbreak{}by decide) (\allowbreak{}by decide) rfl}},
\textnormal{\texttt{QSem.\allowbreak{}val\_\allowbreak{}smul}}, \textnormal{\texttt{QSem.\allowbreak{}run\_\allowbreak{}sound}}, and \textnormal{\texttt{QSem.\allowbreak{}final\_\allowbreak{}sound}}. The file imports
\file{ArchytasZX.PathSumQuditDigits}.

\subsection{\texorpdfstring{\module{archytaszx.semantics.lean_chain}}{lean\_chain}}

\textnormal{\texttt{export\_\allowbreak{}chain}}, \textnormal{\texttt{export\_\allowbreak{}tree}}, and \textnormal{\texttt{export\_\allowbreak{}tree\_\allowbreak{}theorem}} target \file{ChainInduct},
\file{ChainMixed}, and \file{TreeInduct} (\cref{ch:lean-export-net}); they prove each base pair and
each copy pair of a family by the network bodies of this section: \textnormal{\texttt{\_\allowbreak{}bodies}} uses
\textnormal{\texttt{qnet\_\allowbreak{}body}} on an odd prime and otherwise \textnormal{\texttt{net\_\allowbreak{}body}} over one phase modulus, and, with atoms,
\textnormal{\texttt{\_\allowbreak{}atom\_\allowbreak{}bodies}} uses \textnormal{\texttt{net\_\allowbreak{}amp\_\allowbreak{}body}} for pairs with amplitude atoms, \textnormal{\texttt{net\_\allowbreak{}quarter\_\allowbreak{}body}} for
pairs whose reductions take [Quarter] or [Turn] steps, and \textnormal{\texttt{net\_\allowbreak{}body}} for the rest, over one
coefficient vector \textnormal{\texttt{A}}. The bodies' \textnormal{\texttt{Equal}} theorems are the only interface between
\module{archytaszx.semantics.lean_chain} and the modules of this chapter.
\endgroup
